% Les problèmes d'héritage — Allemand Tle A — Cours signé OnBuch+
% Programme officiel MINESEC. Compiler avec : tectonic AL02-les-problemes-d-heritage.tex
\newcommand{\DOCMATIERE}{Allemand}
\newcommand{\DOCNIVEAU}{Tle A}
\newcommand{\DOCMODULE}{Chapitre 1 : Vie familiale et sociale}
\newcommand{\DOCLECON}{AL02}
\newcommand{\DOCTITRE}{Les problèmes d'héritage}
% OnBuch+ — préambule « cours signés » ENRICHI (compilé avec tectonic / XeLaTeX)
% Système visuel « Pop » d'OnBuch+ : fond crème (jamais blanc), bordures encre
% épaisses, ombres « dures » décalées, titres Archivo Black en capitales.
% Version MATHÉMATIQUES : graphes (pgfplots/tikz), boîtes pédagogiques
% (définition, propriété, méthode, attention, le savais-tu, exercice résolu,
% exercice, TP, bilan), page de garde et signature OnBuch+ ancrée en bas.
%
% Macros attendues dans main.tex AVANT \input{preamble} :
%   \DOCMATIERE  \DOCNIVEAU  \DOCMODULE  \DOCLECON (numéro)  \DOCTITRE
\documentclass[11pt]{article}

\usepackage{fontspec}
\usepackage[ngerman,french]{babel} % français = langue principale ; l'allemand via \all{..} / textebox
\usepackage[a4paper,margin=2cm,top=3.4cm,bottom=2.8cm,headheight=1.4cm,footskip=1.3cm]{geometry}
\usepackage{amsmath,amssymb,amsfonts}
\usepackage{mathtools} % \xmapsto, \coloneqq... souvent utilisés par les modèles
\usepackage{cancel} % simplifications barrées (bilans de forces)
% \abs{z} (module, valeur absolue) et \norm{u} (norme), utilisés par les modèles
\providecommand{\abs}{}\let\abs\relax\DeclarePairedDelimiter{\abs}{\lvert}{\rvert}
\providecommand{\norm}{}\let\norm\relax\DeclarePairedDelimiter{\norm}{\lVert}{\rVert}
\usepackage[table]{xcolor} % [table] : fournit aussi \rowcolors (alternance de lignes)
\usepackage{pagecolor}
\usepackage{tikz}
\usetikzlibrary{arrows.meta,positioning,calc,shapes.geometric,shapes.misc,shapes.arrows,decorations.pathmorphing,decorations.pathreplacing,decorations.markings,patterns,shadows,mindmap,trees,fit,backgrounds,matrix,intersections,angles,quotes}
\usepackage{pgfplots}
\usepackage[version=4]{mhchem} % équations nucléaires \ce{^{235}_{92}U}
\usepackage[european,siunitx]{circuitikz} % schémas électriques
\pgfplotsset{compat=1.18}
\usepgfplotslibrary{fillbetween,groupplots} % aires entre courbes (calcul intégral)
\usepackage{enumitem}
\usepackage{fancyhdr}
% [most] charge toutes les bibliothèques tcolorbox (listings, minted, raster,
% documentation...) et gonfle inutilement la mémoire TeX sur un document long
% (~300 boîtes) : on ne charge que ce qui est réellement utilisé.
\usepackage[skins,breakable]{tcolorbox}
\usepackage{graphicx}
\usepackage{colortbl}
\usepackage{booktabs}
\usepackage{tabularx}
\usepackage{diagbox} % case d'en-tête diagonale des tableaux à double entrée
% Colonnes extensibles centrée / gauche / droite (C, L, R), souvent utilisées par les modèles
\newcolumntype{C}{>{\centering\arraybackslash}X}
\newcolumntype{L}{>{\raggedright\arraybackslash}X}
\newcolumntype{R}{>{\raggedleft\arraybackslash}X}
\usepackage{array}
\usepackage{multicol}
\usepackage{multirow}
\usepackage{siunitx}
% parse-numbers=false : accepte aussi \SI{2,5 \times 10^{-3}}{...}
\sisetup{locale=FR,output-decimal-marker={,},group-minimum-digits=4,parse-numbers=false}
\DeclareSIUnit\ohm{\ensuremath{\symup{\Omega}}} % U+2126 absent de la police
\DeclareSIUnit\division{div} % réglages d'oscilloscope (ms/div, V/div)
\DeclareSIUnit\tr{tr} % tours (vitesse de rotation, ex. \SI{1500}{\tr\per\minute})
\DeclareSIUnit\molar{M} % concentration molaire (ex. \SI{3}{\molar}, \SI{1}{\micro\molar})
\DeclareSIUnit\an{an} % année (le modèle confond parfois avec \year, un compteur TeX) — ex. \SI{5}{\kilo\meter\per\an}
\DeclareSIUnit\FCFA{FCFA} % franc CFA (ex. \SI{500}{\FCFA\per\kilo\gram})
\DeclareSIUnit\degreeBrix{\textdegree Brix} % teneur en sucre d'un sirop/jus (réfractométrie)
\DeclareSIUnit\month{mois} % le modèle utilise parfois \month, un compteur TeX, comme unité de durée
\usepackage{qrcode}
% \degree : commande fréquemment utilisée par les modèles pour un angle ou une
% température en dehors de \SI{}{} (non fournie par les paquets chargés).
\providecommand{\degree}{^{\circ}}
% \qed (fin de démonstration), utilisé par les modèles sans amsthm
\providecommand{\qed}{\hfill$\square$}
% \barre : double barre des valeurs interdites dans un tableau de variations (tkz-tab)
\providecommand{\barre}{\ensuremath{\|}}
% environnement proof (amsthm non chargé) utilisé par les modèles
\ifdefined\proof\else\newenvironment{proof}[1][Démonstration]{\par\noindent\textit{#1.}\ }{\qed\par}\fi
\usepackage{lastpage}
\usepackage{hyperref}
\hypersetup{hidelinks}

% ============================================================
% Palette « Pop » OnBuch+ — mêmes hex que la classe P (plus_ui.dart)
% ============================================================
\definecolor{poporange}{HTML}{F59321}
\definecolor{popdark}{HTML}{E07A0C}
\definecolor{popcream}{HTML}{FFF6E8}
\definecolor{popink}{HTML}{1C1714}
\definecolor{poppurple}{HTML}{7A5AE0}
\definecolor{popgreen}{HTML}{2E9E6A}
\definecolor{poppink}{HTML}{E0457B}
\definecolor{popblue}{HTML}{2F7DE1}
\definecolor{popgold}{HTML}{C98A12}
\definecolor{popmuted}{HTML}{8A7F76}
\definecolor{popline}{HTML}{EADFD2}
\definecolor{popgray}{HTML}{8A7F76}
\colorlet{popgrey}{popgray}
% Alias tolérants (le modèle nomme parfois les couleurs différemment)
\colorlet{popwhite}{white}
\colorlet{popblack}{popink}
\colorlet{popred}{poppink}
\colorlet{popyellow}{popgold}
% Teintes claires (fonds de tableaux/encadrés) — le modèle nomme parfois
% les couleurs avec un suffixe L (ex. popblueL) sans les avoir définies.
\colorlet{poporangeL}{poporange!20}
\colorlet{popdarkL}{popdark!20}
\colorlet{poppurpleL}{poppurple!20}
\colorlet{popgreenL}{popgreen!20}
\colorlet{poppinkL}{poppink!20}
\colorlet{popblueL}{popblue!20}
\colorlet{popgoldL}{popgold!20}
\colorlet{popredL}{popred!20}
\colorlet{popgrayL}{popgray!20}
% Teintes claires pour fonds de tableaux / zones
\colorlet{popblueL}{popblue!10!white}
\colorlet{poporangeL}{poporange!14!white}
\colorlet{popgreenL}{popgreen!12!white}
\colorlet{poppurpleL}{poppurple!12!white}
\colorlet{poppinkL}{poppink!12!white}
\colorlet{popgoldL}{popgold!14!white}

\pagecolor{popcream}

% ============================================================
% Typographie
% ============================================================
\defaultfontfeatures{Ligatures=TeX}
\setmainfont{PlusJakartaSans-Medium.ttf}[Path=../../../fonts/,BoldFont=PlusJakartaSans-Bold.ttf]
% Maths en sans-serif (Fira Math) avec chiffres et lettres droites de Plus
% Jakarta, pour que les formules se fondent dans le texte.
\usepackage[math-style=ISO,bold-style=ISO]{unicode-math}
\setmathfont{FiraMath-Regular.otf}[Path=../../../fonts/]
\setmathfont{PlusJakartaSans-Medium.ttf}[Path=../../../fonts/,range={up/num,up/{Latin,latin},\mathup}]
\usepackage{icomma} % virgule décimale sans espace en mode math
% Caractères mathématiques Unicode absents des polices de texte (Plus Jakarta,
% Archivo Black) : rendus par leur équivalent mathématique, en texte comme en
% formule — sinon ils s'affichent en carré vide (ex. « Arithmétique dans ℤ »).
\usepackage{newunicodechar}
\newunicodechar{ℝ}{\ensuremath{\mathbb{R}}}
\newunicodechar{ℤ}{\ensuremath{\mathbb{Z}}}
\newunicodechar{ℂ}{\ensuremath{\mathbb{C}}}
\newunicodechar{ℕ}{\ensuremath{\mathbb{N}}}
\newunicodechar{ℚ}{\ensuremath{\mathbb{Q}}}
\newunicodechar{𝔻}{\ensuremath{\mathbb{D}}}
\newunicodechar{α}{\ensuremath{\alpha}}
\newunicodechar{β}{\ensuremath{\beta}}
\newunicodechar{γ}{\ensuremath{\gamma}}
\newunicodechar{δ}{\ensuremath{\delta}}
\newunicodechar{ε}{\ensuremath{\varepsilon}}
\newunicodechar{θ}{\ensuremath{\theta}}
\newunicodechar{λ}{\ensuremath{\lambda}}
\newunicodechar{μ}{\ensuremath{\mu}}
\newunicodechar{σ}{\ensuremath{\sigma}}
\newunicodechar{τ}{\ensuremath{\tau}}
\newunicodechar{φ}{\ensuremath{\varphi}}
\newunicodechar{ψ}{\ensuremath{\psi}}
\newunicodechar{ω}{\ensuremath{\omega}}
\newunicodechar{Σ}{\ensuremath{\Sigma}}
% majuscules produites par \MakeUppercase dans les titres (α-aminés -> Α)
\newunicodechar{Α}{\ensuremath{\alpha}}
\newunicodechar{Β}{\ensuremath{\beta}}
\newunicodechar{Γ}{\ensuremath{\Gamma}}
\newunicodechar{Δ}{\ensuremath{\Delta}}
\newunicodechar{Ε}{\ensuremath{\varepsilon}}
\newunicodechar{Θ}{\ensuremath{\Theta}}
\newunicodechar{Λ}{\ensuremath{\Lambda}}
\newunicodechar{Μ}{\ensuremath{\mu}}
\newunicodechar{Τ}{\ensuremath{\tau}}
\newunicodechar{Φ}{\ensuremath{\Phi}}
\newunicodechar{Ψ}{\ensuremath{\Psi}}
\newunicodechar{Ω}{\ensuremath{\Omega}}
\newunicodechar{↦}{\ensuremath{\mapsto}}
\newunicodechar{∈}{\ensuremath{\in}}
\newunicodechar{∉}{\ensuremath{\notin}}
\newunicodechar{∧}{\ensuremath{\wedge}}
\newunicodechar{⇔}{\ensuremath{\Leftrightarrow}}
\newunicodechar{⟺}{\ensuremath{\Longleftrightarrow}}
\newunicodechar{⇒}{\ensuremath{\Rightarrow}}
\newunicodechar{⟹}{\ensuremath{\Longrightarrow}}
\newunicodechar{∘}{\ensuremath{\circ}}
\newunicodechar{∪}{\ensuremath{\cup}}
\newunicodechar{∩}{\ensuremath{\cap}}
\newunicodechar{≡}{\ensuremath{\equiv}}
\newunicodechar{⊂}{\ensuremath{\subset}}
\newunicodechar{⊥}{\ensuremath{\perp}}
\newunicodechar{∃}{\ensuremath{\exists}}
\newunicodechar{∀}{\ensuremath{\forall}}
\newunicodechar{∛}{\ensuremath{\sqrt[3]{\,}}}
\newunicodechar{∜}{\ensuremath{\sqrt[4]{\,}}}
\newunicodechar{⃗}{\ensuremath{{}^{\to}}}
\newunicodechar{ⁿ}{\ifmmode^{n}\else\textsuperscript{\ensuremath{n}}\fi}
\newunicodechar{ˣ}{\ifmmode^{x}\else\textsuperscript{\ensuremath{x}}\fi}
\newunicodechar{ᵇ}{\ifmmode^{b}\else\textsuperscript{\ensuremath{b}}\fi}
\newunicodechar{ʳ}{\ifmmode^{r}\else\textsuperscript{\ensuremath{r}}\fi}
\newunicodechar{ᵘ}{\ifmmode^{u}\else\textsuperscript{\ensuremath{u}}\fi}
\newunicodechar{ʸ}{\ifmmode^{y}\else\textsuperscript{\ensuremath{y}}\fi}
\newunicodechar{ᵃ}{\ifmmode^{a}\else\textsuperscript{\ensuremath{a}}\fi}
\newunicodechar{ᵉ}{\ifmmode^{e}\else\textsuperscript{\ensuremath{e}}\fi}
\newunicodechar{ᵏ}{\ifmmode^{k}\else\textsuperscript{\ensuremath{k}}\fi}
\newunicodechar{ᵖ}{\ifmmode^{p}\else\textsuperscript{\ensuremath{p}}\fi}
\newunicodechar{ᴺ}{\ifmmode^{N}\else\textsuperscript{\ensuremath{N}}\fi}
\newunicodechar{ᶜ}{\ifmmode^{c}\else\textsuperscript{\ensuremath{c}}\fi}
\newunicodechar{ʰ}{\ifmmode^{h}\else\textsuperscript{\ensuremath{h}}\fi}
\newunicodechar{ᵠ}{\ifmmode^{q}\else\textsuperscript{\ensuremath{q}}\fi}
\newunicodechar{ₙ}{\ifmmode_{n}\else\textsubscript{\ensuremath{n}}\fi}
\newunicodechar{ₐ}{\ifmmode_{a}\else\textsubscript{\ensuremath{a}}\fi}
\newunicodechar{ᵢ}{\ifmmode_{i}\else\textsubscript{\ensuremath{i}}\fi}
\newunicodechar{ᵣ}{\ifmmode_{r}\else\textsubscript{\ensuremath{r}}\fi}
\newunicodechar{ₖ}{\ifmmode_{k}\else\textsubscript{\ensuremath{k}}\fi}
\newunicodechar{ᵦ}{\ifmmode_{\beta}\else\textsubscript{\ensuremath{\beta}}\fi}
\newunicodechar{ₓ}{\ifmmode_{x}\else\textsubscript{\ensuremath{x}}\fi}
\newfontfamily\popdisplay{ArchivoBlack-Regular.ttf}[Path=../../../fonts/]
\newfontfamily\poplogo{SpaceGrotesk-Bold.ttf}[Path=../../../fonts/]
\newfontfamily\popsemi{PlusJakartaSans-SemiBold.ttf}[Path=../../../fonts/]
\newfontfamily\popxbold{PlusJakartaSans-ExtraBold.ttf}[Path=../../../fonts/]
\newfontfamily\popxboldls{PlusJakartaSans-ExtraBold.ttf}[Path=../../../fonts/,LetterSpace=3.0]

\setlist[enumerate]{leftmargin=1.7em,itemsep=5pt,topsep=4pt}
\setlist[itemize]{leftmargin=1.7em,itemsep=4pt,topsep=4pt,label={\color{poporange}\textbullet}}
\setlength{\parindent}{0pt}
\setlength{\parskip}{5pt}
\renewcommand{\arraystretch}{1.25}

% pgfplots : style Pop par défaut
\pgfplotsset{
  every axis/.append style={axis line style={popink,line width=1pt},tick style={popink},
    grid style={popline,line width=0.5pt},label style={font=\small},tick label style={font=\footnotesize}},
  popcurve/.style={poporange,line width=1.8pt,no markers,smooth},
}
% Même style disponible dans un \draw TikZ simple (hors pgfplots)
\tikzset{popcurve/.style={poporange,line width=1.8pt}}
\tikzset{popgrid/.style={popline,very thin}}
\colorlet{popcurve}{poporange} % le nom du style est parfois utilisé comme couleur

% ============================================================
% Pill / squiggle
% ============================================================
\newcommand{\poppill}[3]{%
  \tikz[baseline=-0.55ex]\node[draw=popink,line width=1.2pt,rounded corners=2.6pt,%
    fill=#2,inner xsep=6.5pt,inner ysep=3pt]{%
    \popxboldls\fontsize{7}{7}\selectfont\color{#3}\MakeUppercase{#1}};%
}
\newcommand{\popsquiggle}[1]{%
  \tikz[baseline=-0.4ex]{
    \draw[#1,line width=2.6pt,line cap=round]
      plot [smooth] coordinates {(0,0.09) (0.34,0.01) (0.68,0.21) (1.02,0.01) (1.36,0.10)};
  }%
}
% Mot-clé surligné (marqueur orange sous le texte)
\newcommand{\cle}[1]{{\popxbold\color{popdark}#1}}

% ============================================================
% En-tête / pied
% ============================================================
\pagestyle{fancy}
\fancyhf{}
\renewcommand{\headrulewidth}{0pt}
\renewcommand{\footrulewidth}{0pt}
\lhead{\raisebox{-0.15ex}{\poplogo\fontsize{13}{13}\selectfont\color{popink}OnBuch\color{poporange}+}}
\rhead{\poppill{\DOCMATIERE\ \textbullet\ \DOCNIVEAU\ \textbullet\ Leçon \DOCLECON}{white}{popink}}
\lfoot{\popxbold\fontsize{7.6}{7.6}\selectfont\color{popmuted}\MakeUppercase{Cours signé OnBuch+}\ \textbullet\ {\popsemi\DOCTITRE}}
\rfoot{\tikz[baseline=-0.6ex]\node[rounded corners=8.5pt,fill=popink,minimum height=17pt,minimum width=17pt,inner xsep=5pt,inner sep=0]{\popxbold\fontsize{8}{8}\selectfont\color{popcream}\thepage\,/\,\pageref*{LastPage}};}

% ============================================================
% Titres
% ============================================================
\newcommand{\chaptitre}[2]{%
  \begin{tcolorbox}[enhanced,colback=poporange,colframe=popink,boxrule=2.6pt,arc=11pt,
      drop shadow=popink,left=15pt,right=15pt,top=12pt,bottom=11pt,before skip=3pt,after skip=18pt]
    \poppill{#2}{popink}{popcream}\par\vspace{9pt}
    {\raggedright\hyphenpenalty=10000\exhyphenpenalty=10000\popdisplay\fontsize{22}{24}\selectfont\color{popcream}\MakeUppercase{#1}\par}
  \end{tcolorbox}
}
% Section numérotée : pastille orange avec le numéro + titre Archivo
\newcounter{popsec}
\newcommand{\coursec}[1]{%
  \stepcounter{popsec}\setcounter{popsub}{0}%
  \par\vspace{16pt}\noindent
  \begin{minipage}{\linewidth}
  \tikz[baseline=-0.7ex]\node[draw=popink,line width=1.6pt,fill=poporange,rounded corners=4pt,minimum size=22pt,inner sep=0]{\popdisplay\fontsize{12}{12}\selectfont\color{popcream}\thepopsec};%
  \hspace{8pt}{\popdisplay\fontsize{15}{17}\selectfont\color{popink}\MakeUppercase{#1}}\par
  \vspace{4pt}\noindent{\color{poporange}\rule{36pt}{3.6pt}}
  \end{minipage}\par\nopagebreak\vspace{8pt}
}
\newcounter{popsub}
\newcommand{\courssub}[1]{%
  \stepcounter{popsub}%
  \par\vspace{9pt}\noindent{\popxbold\fontsize{11.5}{13}\selectfont\color{popink}{\color{poporange}\thepopsec.\thepopsub}\hspace{6pt}#1}\par\nopagebreak\vspace{4pt}
}

% ============================================================
% Boîtes pédagogiques — même recette : fond blanc, bordure épaisse colorée,
% ombre dure encre, badge plein. Titre optionnel [#1].
% ============================================================
\tcbset{popbase/.style={breakable,enhanced,colback=white,boxrule=2.3pt,arc=9pt,
  shadow={3pt}{-3pt}{0pt}{popink},left=14pt,right=14pt,top=11pt,bottom=11pt,before skip=11pt,after skip=15pt,
  fontupper=\popsemi\fontsize{10.6}{14.5}\selectfont\color{popink},
  fontlower=\popsemi\fontsize{10.6}{14.5}\selectfont\color{popink}}}
% Coche dessinée (le glyphe ✓ n'existe pas dans les polices du document)
\newcommand{\popcheck}[1]{\tikz[baseline=-0.5ex]\draw[#1,line width=1.6pt,line cap=round,line join=round](-0.13,0.0)--(-0.04,-0.09)--(0.13,0.1);}
\providecommand{\coche}{\popcheck{popgreen}}
\providecommand{\resultat}[1]{\par\smallskip\noindent\fcolorbox{popgreen}{popgreenL}{\parbox{\dimexpr\linewidth-2\fboxsep-2\fboxrule\relax}{\textbf{Résultat.} #1}}\par\smallskip}
\newcommand{\popboxhead}[3]{\poppill{#1}{#2}{white}\ifx\relax#3\relax\else\hspace{6pt}{\popxbold\fontsize{10.6}{12}\selectfont\color{popink}#3}\fi\par\vspace{7pt}}

\newtcolorbox{aretenir}[1][]{popbase,colframe=popgreen,before upper={\popboxhead{À retenir}{popgreen}{#1}}}
% Texte en allemand (document de lecture, dialogue, extrait) : cadre
% d'étude. Un titre optionnel (source, auteur) s'affiche dans l'en-tête.
\newtcolorbox{textebox}[1][]{popbase,colframe=popblue,before upper={\popboxhead{Text}{popblue}{#1}\begin{otherlanguage}{ngerman}},after upper={\end{otherlanguage}}}
% Marques ✓ / ✗ (forme correcte / fautive) souvent utilisées par les modèles
\providecommand{\cross}{\textcolor{poppink}{\ensuremath{\times}}}
\providecommand{\xmark}{\cross}\providecommand{\crossmark}{\cross}\providecommand{\cmark}{\ensuremath{\checkmark}}
% Ligne de réponse / justification à compléter (inventée par les modèles)
\providecommand{\justification}[1]{\par\noindent\textit{Justification~:} \dotfill\par}
% \textipa (package tipa non chargé) : la police ne couvre pas l'API -> simple passage
\providecommand{\textipa}[1]{#1}
% Soulignement (ulem non chargé) : \uline, \ul
\providecommand{\uline}[1]{\underline{#1}}\providecommand{\ul}[1]{\underline{#1}}
% variantes ASCII d'ordinaux écrites en macro
\providecommand{\iere}{\textsuperscript{re}}\providecommand{\ieme}{\textsuperscript{e}}\providecommand{\ere}{\textsuperscript{re}}\providecommand{\eme}{\textsuperscript{e}}\providecommand{\er}{\textsuperscript{er}}\providecommand{\ie}{\textsuperscript{e}}
% Ordinaux français écrits en macro par les modèles : 1\ière, 3\ième
\newcommand{\ière}{\textsuperscript{re}}\newcommand{\ième}{\textsuperscript{e}}
% \exemple (inventée par les modèles) : étiquette « Exemple : »
\providecommand{\exemple}{\textbf{Exemple~:}\ }
% \square utilisable aussi hors mode math (cases à cocher)
\AtBeginDocument{\let\squareMath\square\renewcommand{\square}{\ensuremath{\squareMath}}}
% Mot ou phrase en allemand à l'intérieur d'un paragraphe français
\newcommand{\all}[1]{\ifmmode\text{\textit{#1}}\else\begingroup\selectlanguage{ngerman}\itshape #1\endgroup\fi}
\let\esp\all % alias toléré
\newtcolorbox{exemplebox}[1][]{popbase,colframe=poporange,before upper={\popboxhead{Exemple}{poporange}{#1}}}
\newtcolorbox{definition}[1][]{popbase,colframe=popblue,before upper={\popboxhead{Définition}{popblue}{#1}}}
\newtcolorbox{propriete}[1][]{popbase,colframe=poppurple,before upper={\popboxhead{Propriété}{poppurple}{#1}}}
\newtcolorbox{methode}[1][]{popbase,colframe=popgold,before upper={\popboxhead{Méthode}{popgold}{#1}}}
\newtcolorbox{attention}[1][]{popbase,colframe=poppink,before upper={\popboxhead{Attention}{poppink}{#1}}}
\newtcolorbox{experience}[1][]{popbase,colframe=popblue,colback=popblueL,before upper={\popboxhead{Expérience}{popblue}{#1}}}
% « Le savais-tu ? » : carte violette pleine (contexte, culture, Cameroun)
\newtcolorbox{savaistu}[1][]{popbase,colback=poppurple,colframe=popink,
  fontupper=\popsemi\fontsize{10.4}{14.2}\selectfont\color{white},
  before upper={\poppill{Le savais-tu\,?}{white}{poppurple}\ifx\relax#1\relax\else\hspace{6pt}{\popxbold\color{white}#1}\fi\par\vspace{7pt}}}
% Exercice résolu : énoncé en haut, \tcblower puis la solution sur fond crème
\newcounter{popexo}\newcounter{popexr}
\newtcolorbox{exoresolu}[1][]{popbase,colframe=poporange,colbacklower=popcream,
  segmentation style={popink,line width=1.2pt,dashed},
  before upper={\stepcounter{popexr}\popboxhead{Exercice résolu \thepopexr}{poporange}{#1}},
  before lower={\poppill{Solution}{popink}{popcream}\par\vspace{6pt}}}
% Exercice (non résolu en place) — la correction est dans la partie Corrigés
\newtcolorbox{exercice}[1][]{popbase,colframe=popink,
  before upper={\stepcounter{popexo}\popboxhead{Exercice \thepopexo}{popink}{#1}}}
% Corrigé
\newtcolorbox{corrige}[1][]{popbase,colframe=popgreen,colback=popgreenL,
  before upper={\popboxhead{Corrigé}{popgreen}{#1}}}
% Objectifs / prérequis / bilan
\newtcolorbox{objectifs}{popbase,colframe=popink,colback=poporangeL,
  before upper={\popboxhead{Objectifs de la leçon}{popink}{}}}
\newtcolorbox{prerequis}{popbase,colframe=popink,colback=popblueL,
  before upper={\popboxhead{Prérequis}{popblue}{}}}
\newtcolorbox{bilan}{popbase,colframe=popink,colback=white,boxrule=2.8pt,
  before upper={\popboxhead{Fiche bilan}{poporange}{L'essentiel à connaître pour l'examen}}}
% Figure encadrée avec légende
\newtcolorbox{popfigure}[1][]{enhanced,breakable=false,colback=white,colframe=popline,boxrule=1.4pt,arc=8pt,
  left=10pt,right=10pt,top=10pt,bottom=8pt,before skip=10pt,after skip=12pt,halign=center,
  lower separated=false,halign lower=center,
  fontlower=\popsemi\fontsize{9}{11.5}\selectfont\color{popmuted},
  #1}
\newcommand{\legende}[1]{\par\vspace{6pt}{\popsemi\fontsize{9}{11.5}\selectfont\color{popmuted}#1}}

% ============================================================
% Signature OnBuch+
% ============================================================
\newcommand{\poplogoinline}[1]{{\poplogo\fontsize{#1}{#1}\selectfont\color{popink}OnBuch\color{poporange}+}}
\newcommand{\signatureonbuch}{%
  \begin{tcolorbox}[enhanced,colback=popink,colframe=popink,boxrule=2.2pt,arc=10pt,
      drop shadow=poporange,left=14pt,right=14pt,top=10pt,bottom=10pt,before skip=0pt,after skip=0pt]
    \tikz[baseline=-0.7ex]\node[circle,fill=popgreen,draw=popcream,line width=1.3pt,minimum size=22pt,inner sep=0]{\popcheck{white}};%
    \hspace{9pt}\begin{minipage}[c]{0.62\linewidth}
      {\popxbold\fontsize{12}{13}\selectfont\color{popcream}Signé\ \textbullet\ {\poplogo OnBuch\color{poporange}+}}\par\vspace{2pt}
      {\raggedright\popsemi\fontsize{8}{10.5}\selectfont\color{popline}\MakeUppercase{Équipe pédagogique OnBuch+}\\\MakeUppercase{Conforme au programme officiel MINESEC}\par}
    \end{minipage}\hfill
    {\poplogo\fontsize{22}{22}\selectfont\color{popcream}OnBuch\color{poporange}+}
  \end{tcolorbox}%
}

% ============================================================
% Page de garde
% #1 titre · #2 matière · #3 niveau · #4 module · #5 n° leçon · #6 durée ·
% #7 description / objectifs (texte ou itemize)
% ============================================================
\newcommand{\pagegarde}[7]{%
  \thispagestyle{empty}
  \newgeometry{a4paper,margin=1.7cm,top=1.6cm,bottom=1.4cm}%
  \begin{tikzpicture}[remember picture,overlay]
    \fill[poporange!18] ([xshift=-3.2cm,yshift=-3.6cm]current page.north east) circle (4.6cm);
    \fill[poppurple!14] ([xshift=2.4cm,yshift=7.6cm]current page.south west) circle (3.8cm);
  \end{tikzpicture}%
  {\poplogo\fontsize{30}{30}\selectfont\color{popink}OnBuch\color{poporange}+}\hfill
  \raisebox{0.6ex}{\poppill{Cours signé \textbullet\ Premium}{popink}{popcream}}\par
  \vspace{1pt}\popsquiggle{poppurple}\par
  \vspace{8pt}
  \begin{tcolorbox}[enhanced,colback=poporange,colframe=popink,boxrule=2.8pt,arc=13pt,
      drop shadow=popink,left=18pt,right=18pt,top=12pt,bottom=13pt,after skip=10pt]
    \poppill{#2 \textbullet\ #3}{popink}{popcream}\hspace{5pt}\poppill{#4}{white}{popink}\par\vspace{8pt}
    {\popdisplay\fontsize{14}{14}\selectfont\color{popink}LEÇON #5}\par\vspace{4pt}
    {\raggedright\hyphenpenalty=10000\exhyphenpenalty=10000\popdisplay\fontsize{25}{27}\selectfont\color{popcream}\MakeUppercase{#1}\par}
    \vspace{8pt}
    {\popxbold\fontsize{9.5}{9.5}\selectfont\color{popink}\MakeUppercase{Durée conseillée : #6}}
  \end{tcolorbox}
  \begin{tcolorbox}[enhanced,colback=white,colframe=popink,boxrule=2.2pt,arc=10pt,
      drop shadow=popink,left=16pt,right=16pt,top=10pt,bottom=10pt,after skip=8pt]
    \poppill{Dans cette leçon}{poppurple}{white}\par\vspace{5pt}
    \popsemi\fontsize{9.3}{12.6}\selectfont\color{popink}#7
  \end{tcolorbox}
  \noindent\begin{minipage}[t]{0.487\linewidth}
    \begin{tcolorbox}[enhanced,colback=popgreen,colframe=popink,boxrule=2.2pt,arc=10pt,
        drop shadow=popink,left=11pt,right=11pt,top=10pt,bottom=10pt,height=3.1cm,valign=center]
      \tcbox[colback=white,colframe=popink,boxrule=1.3pt,arc=5pt,boxsep=0pt,nobeforeafter,
        left=3pt,right=3pt,top=3pt,bottom=3pt,valign=center]{\qrcode[height=1.9cm]{https://play.google.com/store/apps/details?id=com.luvvix.onbuch&hl=fr}}%
      \hspace{8pt}\begin{minipage}[c]{0.5\linewidth}
        {\popdisplay\fontsize{11}{12.5}\selectfont\color{white}\MakeUppercase{Télécharge OnBuch}}\par\vspace{4pt}
        {\raggedright\popsemi\fontsize{8.4}{11}\selectfont\color{white}Gratuit \textbullet\ Google Play\\Tuteur IA, annales, concours, résultats\par}
      \end{minipage}
    \end{tcolorbox}
  \end{minipage}\hfill
  \begin{minipage}[t]{0.487\linewidth}
    \begin{tcolorbox}[enhanced,colback=poppurple,colframe=popink,boxrule=2.2pt,arc=10pt,
        drop shadow=popink,left=11pt,right=11pt,top=10pt,bottom=10pt,height=3.1cm,valign=center]
      \tikz[baseline=-0.7ex]\node[circle,fill=white,draw=popink,line width=1.3pt,minimum size=1.6cm,inner sep=0]{\popdisplay\fontsize{24}{24}\selectfont\color{poppurple}+};%
      \hspace{8pt}\begin{minipage}[c]{0.6\linewidth}
        {\popdisplay\fontsize{11}{12.5}\selectfont\color{white}\MakeUppercase{Passe à OnBuch+}}\par\vspace{4pt}
        {\raggedright\popsemi\fontsize{8.4}{11}\selectfont\color{white}Tous les cours signés, Tuteur IA illimité, fiches PDF — onglet OnBuch+ de l'app\par}
      \end{minipage}
    \end{tcolorbox}
  \end{minipage}
  \vspace{8pt}
  \popcontenu
  \vfill
  \signatureonbuch
  \restoregeometry
  \newpage
}

% Rangée « Ce cours contient » (page de garde)
\newcommand{\poptile}[3]{% #1 couleur #2 gros texte #3 légende
  \begin{tcolorbox}[enhanced,colback=#1,colframe=popink,boxrule=2pt,arc=9pt,shadow={2.5pt}{-2.5pt}{0pt}{popink},
      left=6pt,right=6pt,top=9pt,bottom=9pt,halign=center,valign=center,height=1.75cm,width=\linewidth,nobeforeafter]
    {\popdisplay\fontsize{12.5}{13}\selectfont\color{white}\MakeUppercase{#2}}\par\vspace{3pt}
    {\centering\popsemi\fontsize{7.8}{9.5}\selectfont\color{white}#3\par}
  \end{tcolorbox}}
\newcommand{\popcontenu}{%
  \par\noindent{\popxboldls\fontsize{7.5}{7.5}\selectfont\color{popmuted}CE COURS SIGNÉ CONTIENT}\par\vspace{7pt}
  \noindent\begin{minipage}[t]{0.235\linewidth}\poptile{popblue}{Cours}{complet, illustré, pas à pas}\end{minipage}\hfill
  \begin{minipage}[t]{0.235\linewidth}\poptile{poporange}{Méthodes}{et exercices résolus}\end{minipage}\hfill
  \begin{minipage}[t]{0.235\linewidth}\poptile{poppink}{Exercices}{type examen + corrigés}\end{minipage}\hfill
  \begin{minipage}[t]{0.235\linewidth}\poptile{popgreen}{Fiche bilan}{l'essentiel à retenir}\end{minipage}\par}

% Sommaire
\newtcolorbox{sommaire}{enhanced,colback=white,colframe=popink,boxrule=2.2pt,arc=10pt,
  drop shadow=popink,left=18pt,right=18pt,top=13pt,bottom=13pt,before skip=6pt,after skip=20pt,
  before upper={\poppill{Sommaire}{poppurple}{white}\par\vspace{9pt}\popsemi\fontsize{10.6}{14.5}\selectfont\color{popink}}}

% Signature de fin de cours — poussée tout en bas de la dernière page
\newcommand{\findecours}{%
  \par\vfill\nopagebreak
  \begin{center}\popsquiggle{poporange}\end{center}\vspace{4pt}
  \signatureonbuch
}
\AtBeginDocument{\def\llbracket{\mathopen{[\mkern-3.5mu[}}\def\rrbracket{\mathclose{]\mkern-3.5mu]}}} % intervalles d'entiers (glyphe ⟦ absent de la police)
% Glyphes absents de Fira Math : redéfinis à partir de glyphes disponibles
\AtBeginDocument{%
  \def\setminus{\mathbin{\backslash}}\let\smallsetminus\setminus
  \def\vdots{\mathord{\vbox{\baselineskip3.2pt\lineskiplimit0pt\kern4pt\hbox{.}\hbox{.}\hbox{.}}}}%
  \def\ddots{\mathinner{\mkern1mu\raise6.5pt\hbox{.}\mkern2mu\raise3.6pt\hbox{.}\mkern2mu\raise0.7pt\hbox{.}\mkern1mu}}%
  \def\triangle{\mathord{\scalebox{0.9}{$\symup{\Delta}$}}}%
  \def\nabla{\mathord{\raisebox{\depth}{\scalebox{1}[-1]{$\symup{\Delta}$}}}}%
  \def\checkmark{\popcheck{popdark}}%
  \def\star{\mathbin{\ast}}\def\bigstar{\mathord{\ast}}%
  \def\diamond{\mathbin{\scalebox{0.6}{$\Diamond$}}}%
  \def\bigcup{\mathop{\mathchoice{\raisebox{-0.3em}{\scalebox{1.7}{$\cup$}}}{\scalebox{1.25}{$\cup$}}{\cup}{\cup}}\displaylimits}%
  \def\bigcap{\mathop{\mathchoice{\raisebox{-0.3em}{\scalebox{1.7}{$\cap$}}}{\scalebox{1.25}{$\cap$}}{\cap}{\cap}}\displaylimits}%
  \def\lesseqqgtr{\mathrel{\lessgtr}}\def\bigoplus{\mathop{\oplus}\displaylimits}%
}
\providecommand{\ssi}{si et seulement si} % abréviation fréquente
\AtBeginDocument{\providecommand{\wideparen}{\overparen}} % arc de cercle

%% FIGFIX-BEGIN v4 — correctifs globaux des figures (docs/onbuch-generation/tools/figfix)
\usepackage{adjustbox}\usepackage{etoolbox}
% toute figure TikZ/pgfplots plus large que la ligne est réduite à la largeur disponible
\BeforeBeginEnvironment{tikzpicture}{\begin{adjustbox}{max width=\linewidth}}
\AfterEndEnvironment{tikzpicture}{\end{adjustbox}}
% légendes pgfplots lisibles par-dessus les courbes
\pgfplotsset{every axis/.append style={legend style={fill=white,fill opacity=0.92,text opacity=1,draw=black!15,inner sep=3pt}}}
% étiquettes d'arêtes (syntaxe edge["..."]) avec fond blanc
\tikzset{every edge quotes/.append style={fill=white,inner sep=1.2pt}}
%% FIGFIX-END
\begin{document}

\pagegarde{Les problèmes d'héritage}{Allemand}{Tle A}{Chapitre 1 : Vie familiale et sociale}{AL02}{2 h}{Cette leçon de type « texte » propose la lecture et le commentaire d'un texte allemand sur les conflits d'héritage dans une famille. Elle permet d'acquérir le vocabulaire thématique (Erbe, Testament, Streit, Gesetz), de maîtriser le génitif et l'alternance parfait/prétérit pour raconter, et de produire un résumé et un court texte argumenté.
\vspace{4pt}
\begin{multicols}{2}\small
\begin{itemize}
\item À la fin de cette leçon, je sais comprendre globalement puis en détail un texte allemand sur les problèmes d'héritage.
\item À la fin de cette leçon, je sais employer correctement le vocabulaire du thème (das Erbe, der Erbe, das Testament, der Streit, das Gesetz) avec articles et pluriels.
\item À la fin de cette leçon, je sais former et utiliser le génitif pour exprimer la possession et les compléments du nom.
\item À la fin de cette leçon, je sais conjuguer les verbes au parfait avec haben ou sein selon les règles.
\item À la fin de cette leçon, je sais conjuguer les verbes faibles au prétérit et choisir entre parfait et prétérit pour raconter.
\item À la fin de cette leçon, je sais répondre par des phrases complètes à des questions de compréhension.
\end{itemize}\end{multicols}}

\begin{sommaire}
\begin{multicols}{2}
\begin{enumerate}
\item Texte support : Der Streit um das Erbe
\item Compréhension du texte
\item Lexique thématique : Erbe, Testament, Familie, Gesetz
\item Grammaire 1 : le génitif
\item Grammaire 2 : parfait et prétérit pour raconter
\item Commentaire et production guidée
\item Activité d'intégration
\item Méthodes et exercices résolus
\item Exercices
\item Corrigés des exercices
\item Fiche bilan
\end{enumerate}
\end{multicols}
\end{sommaire}

\begin{objectifs}
\begin{itemize}
\item À la fin de cette leçon, je sais comprendre globalement puis en détail un texte allemand sur les problèmes d'héritage.
\item À la fin de cette leçon, je sais employer correctement le vocabulaire du thème (das Erbe, der Erbe, das Testament, der Streit, das Gesetz) avec articles et pluriels.
\item À la fin de cette leçon, je sais former et utiliser le génitif pour exprimer la possession et les compléments du nom.
\item À la fin de cette leçon, je sais conjuguer les verbes au parfait avec haben ou sein selon les règles.
\item À la fin de cette leçon, je sais conjuguer les verbes faibles au prétérit et choisir entre parfait et prétérit pour raconter.
\item À la fin de cette leçon, je sais répondre par des phrases complètes à des questions de compréhension.
\item À la fin de cette leçon, je sais rédiger un résumé et un court paragraphe d'opinion sur un conflit d'héritage.
\item À la fin de cette leçon, je sais commenter un texte selon la méthode (introduction, résumé, analyse, opinion).
\end{itemize}
\end{objectifs}

\begin{prerequis}
\begin{itemize}
\item Les articles définis et indéfinis au nominatif, accusatif et datif (classe de Première)
\item Le parfait des verbes réguliers avec haben (Première)
\item Le vocabulaire de la famille (die Familie, die Eltern, die Geschwister)
\item La structure de la phrase allemande : place du verbe conjugué en position 2
\item Les pronoms personnels et possessifs (mein, dein, sein, ihr...)
\end{itemize}
\end{prerequis}

\newpage

\coursec{Texte support : Der Streit um das Erbe}

À Douala comme à Berlin, la mort d'un parent déclenche souvent de vives discussions : qui hérite de la maison, du terrain, de l'argent ? En Allemagne, un testament écrit règle généralement la succession, mais les conflits familiaux restent fréquents. Au Cameroun, la coexistence du droit moderne et des coutumes rend la question encore plus délicate. Comment parler en allemand de l'héritage, du testament et des disputes familiales ? Cette leçon nous y prépare à travers un texte authentique et le vocabulaire juridique de base.

\courssub{Lecture globale et situation de communication}

\begin{textebox}[Der Streit um das Erbe]
Als der Großvater \cle{starb}, hinterließ er ein großes Haus in München, ein gut gefülltes Konto und wertvollen Schmuck. Sein \cle{Testament} war jedoch nicht eindeutig formuliert: Das Haus der Familie sollte an die älteste Tochter, Anna, gehen, aber der Sohn des Großvaters, Thomas, bestand auf seinem gesetzlichen \cle{Pflichtteil}. Die \cle{Erben} stritten monatelang vor Gericht. Die Mutter der Kinder versuchte zu \cle{vermitteln}, aber der \cle{Streit} wurde immer größer und teurer. Schließlich hat ein \cle{Notar} die \cle{Erbschaft} geregelt: Das Haus wurde verkauft und das Geld wurde unter den Erben aufgeteilt. Am Ende waren alle froh, dass der Konflikt vorbei war, aber die Familie hat viel Zeit, viel Geld und viele \cle{Nerven} verloren. Der \cle{Erblasser} hatte wohl gehofft, mit seinem Testament Frieden zu stiften, doch das Gegenteil war der Fall.
\end{textebox}

\begin{center}
\begin{tabularx}{\linewidth}{|X|X|}\hline
\rowcolor{popblueL} \textbf{Deutsch} & \textbf{Français} \\ \hline
\all{sterben} (starb, ist gestorben) & mourir \\ \hline
\all{hinterlassen} (hinterließ, hat hinterlassen) & laisser (après sa mort) \\ \hline
\all{das Testament, die Testamente} & le testament \\ \hline
\all{der Erbe, die Erben / die Erbin, die Erbinnen} & l'héritier / l'héritière \\ \hline
\all{die Erbschaft, die Erbschaften} & l'héritage, la succession \\ \hline
\all{der Erblasser, die Erblasser} & le défunt (celui qui laisse l'héritage) \\ \hline
\all{der Pflichtteil, die Pflichtteile} & la réserve héréditaire (part minimale légale) \\ \hline
\all{der Notar, die Notare} & le notaire \\ \hline
\all{vermitteln} & servir de médiateur, tenter de concilier \\ \hline
\all{der Streit, die Streite} & le conflit, la dispute \\ \hline
\all{regeln} & régler, organiser \\ \hline
\all{aufteilen} (teilte auf, hat aufgeteilt) & partager, répartir \\ \hline
\all{die Nerven (Pl.)} & les nerfs (au fig. : le calme, la patience) \\ \hline
\all{stiften} & fonder, créer (ici : apporter) \\ \hline
\all{das Gegenteil, die Gegenteile} & le contraire \\ \hline
\all{eindeutig} & clair, sans ambiguïté \\ \hline
\all{vor Gericht} & devant le tribunal \\ \hline
\all{Frieden stiften} & rétablir la paix, reconcilier \\ \hline
\end{tabularx}
\end{center}

\courssub{Repérage des personnages, du lieu et du problème central}

\begin{center}
\begin{tabularx}{\linewidth}{|X|X|X|}\hline
\rowcolor{poporangeL} \textbf{Élément} & \textbf{Information du texte} & \textbf{Mots-clés allemands} \\ \hline
Personnages principaux & Le grand-père (défunt), Anna (fille aînée), Thomas (fils), la mère, le notaire & \all{der Großvater, die Tochter, der Sohn, die Mutter, der Notar} \\ \hline
Lieu & Munich (Allemagne), tribunal, étude notariale & \all{München, vor Gericht, beim Notar} \\ \hline
Objets de l'héritage & Une grande maison, un compte en banque, des bijoux & \all{das Haus, das Konto, der Schmuck} \\ \hline
Problème central & Testament ambigu, dispute sur la part légale (Pflichtteil), vente forcée & \all{uneindeutig, der Pflichtteil, der Verkauf} \\ \hline
Issue & Vente de la maison, partage de l'argent, relations familiales fatiguées & \all{verkauft, aufgeteilt, verloren} \\ \hline
\end{tabularx}
\end{center}

\courssub{Lecture détaillée par paragraphe et repérage grammatical}

\textbf{Paragraphe 1 :} \all{„Als der Großvater starb, hinterließ er ein großes Haus in München …“}
\begin{itemize}
    \item \textbf{Information clé :} Mort du grand-père, énumération des biens.
    \item \textbf{Grammaire :} \cle{Präteritum} : \all{starb} (verbe fort), \all{hinterließ} (verbe fort). Observez la place du verbe : en position 2 dans la principale (\all{hinterließ er}) et en position finale dans la subordonnée introduite par \all{als} (\all{Als der Großvater starb}).
\end{itemize}

\textbf{Paragraphe 1 (suite) :} \all{„Sein Testament war jedoch nicht eindeutig formuliert …“}
\begin{itemize}
    \item \textbf{Information clé :} Le testament est flou (\all{nicht eindeutig}). Conflit : Anna (la maison) contre Thomas (la part légale).
    \item \textbf{Grammaire :} \cle{Génitif} : \all{das Haus der Familie} (la maison \textbf{de la famille}), \all{der Sohn des Großvaters} (le fils \textbf{du grand-père}). \cle{Präteritum} : \all{war} (\all{sein}), \all{sollte} (\all{sollen}), \all{bestand} (\all{bestehen auf} + Datif : insister sur).
\end{itemize}

\textbf{Paragraphe 2 :} \all{„Die Erben stritten monatelang vor Gericht.“}
\begin{itemize}
    \item \textbf{Information clé :} Procès long (\all{monatelang} = pendant des mois).
    \item \textbf{Grammaire :} \cle{Präteritum} fort : \all{stritten} (\all{streiten}, \all{stritt}, \all{hat gestritten} — verbe fort, malgré sa ressemblance avec un verbe faible). \cle{Präteritum} faible : \all{versuchte} (\all{versuchen}).
\end{itemize}

\textbf{Paragraphe 3 :} \all{„Schließlich hat ein Notar die Erbschaft geregelt …“}
\begin{itemize}
    \item \textbf{Information clé :} Intervention du notaire, vente de la maison, partage de l'argent.
    \item \textbf{Grammaire :} \cle{Perfekt} actif avec auxiliaire \all{haben} : \all{hat ... geregelt}. \cle{Präteritum} du passif : \all{wurde verkauft}, \all{wurde aufgeteilt} (auxiliaire \all{werden} au Präteritum + participe II). \cle{Datif pluriel} : \all{unter den Erben} (parmi les héritiers).
\end{itemize}

\textbf{Paragraphe 4 :} \all{„Am Ende waren alle froh, dass der Konflikt vorbei war …“}
\begin{itemize}
    \item \textbf{Information clé :} Soulagement, mais coût élevé (temps, argent, nerfs). Ironie : le testament voulait apporter la paix (\all{Frieden stiften}), il a provoqué la guerre.
    \item \textbf{Grammaire :} \cle{Präteritum} : \all{waren}, \all{war}. \cle{Perfekt} : \all{hat ... verloren}. \cle{Plusquamperfekt} (antériorité) : \all{hatte ... gehofft}. Subordonnée en \all{dass} : verbe conjugué à la fin (\all{dass der Konflikt vorbei war}).
\end{itemize}

\begin{methode}[Comment lire un texte allemand efficacement]
\begin{enumerate}
    \item \textbf{Titre et hypothèses :} Lisez le titre (\all{Der Streit um das Erbe}). De quoi s'agit-il ? Activez votre vocabulaire sur la famille, la mort, le droit.
    \item \textbf{Lecture globale (Skimming) :} Lisez une première fois \textbf{sans dictionnaire}. Cherchez : Qui ? Quand ? Où ? Quel dénouement ? Repérez les verbes conjugués (indices temporels).
    \item \textbf{Lecture détaillée (Scanning) :} Relisez paragraphe par paragraphe avec le glossaire. Soulignez les \textbf{mots-clés} (noms, verbes, connecteurs logiques : \all{jedoch, schließlich, aber, doch}).
    \item \textbf{Repérage grammatical ciblé :} Surlignez d'une couleur les formes du \textbf{Präteritum} (récit), d'une autre les formes du \textbf{Perfekt} (résultat, oral), entourez les \textbf{Génitifs} (possession).
    \item \textbf{Reformulation orale :} Fermez le texte. Résumez l'histoire en 3 ou 4 phrases allemandes simples (au Präteritum ou au Perfekt).
\end{enumerate}
\end{methode}

\begin{attention}[Piège de traduction : mot à mot contre sens global]
Ne traduisez jamais mot à mot !
\begin{itemize}
    \item \all{„Er hinterließ ein Haus.“} ne signifie pas « Il a laissé derrière une maison », mais \textbf{« Il a laissé une maison (à sa mort). »}
    \item \all{„das Haus der Familie“} : \all{der Familie} est un \textbf{génitif singulier féminin} (la maison de la famille), pas un pluriel.
    \item \all{„Er bestand auf seinem Pflichtteil.“} ne signifie pas « Il a consisté sur sa part », mais \textbf{« Il a insisté sur sa réserve héréditaire. »} (\all{bestehen auf} + Datif).
    \item \all{„die Nerven verlieren“} : expression idiomatique, \textbf{« perdre patience, craquer nerveusement »}, et non « perdre les nerfs » au sens propre.
    \item Faux ami : \all{das Gesetz} = la loi (juridique), et non « le gaz » ; \all{der Erbe} = l'héritier (personne), \all{das Erbe} = l'héritage (biens).
\end{itemize}
\end{attention}

\begin{savaistu}[Le \all{Pflichtteil} : une protection légale en Allemagne]
En droit successoral allemand (\all{das Erbrecht}), les descendants (enfants, petits-enfants), le conjoint survivant et, à défaut, les parents du défunt ont droit à un \textbf{minimum légal} : le \cle{Pflichtteil}. Il correspond à \textbf{la moitié de la part légale} (\all{der gesetzliche Erbteil}). Même si le testament les déshérite (\all{enterben}), ils peuvent réclamer cette somme en \textbf{argent} (et non en nature). Au Cameroun, la loi organise aussi la réserve héréditaire, mais la coutume (droit coutumier) prévaut souvent dans la pratique, notamment pour la transmission des terres ancestrales : d'où de fréquents conflits successoraux, comme dans notre texte.
\end{savaistu}

\begin{definition}[Famille de mots : \all{erben} — \all{vererben} — \all{das Erbe}]
\begin{itemize}
    \item \all{das Erbe, die Erben} : l'héritage (les biens) ; \all{der Erbe, die Erben} : l'héritier ; \all{die Erbin, die Erbinnen} : l'héritière. Attention : même pluriel \all{Erben}, mais articles différents au singulier (\all{das} / \all{der}).
    \item \all{erben} (verbe fort : \all{erbte, hat geerbt}) : hériter (recevoir). \all{Ich habe das Haus geerbt.} « J'ai hérité de la maison. »
    \item \all{vererben} (verbe faible : \all{vererbte, hat vererbt}) : léguer, transmettre par testament (donner). \all{Der Großvater vererbte das Haus an Anna.} « Le grand-père a légué la maison à Anna. »
    \item \all{die Erbschaft, die Erbschaften} : la succession, l'héritage (aspect juridique, patrimoine).
    \item \all{der Erblasser, die Erblasser} : le défunt, l'auteur de la succession (celui qui « laisse »).
    \item \all{enterben} (verbe faible) : déshériter. \all{Er enterbte seinen Sohn.} « Il a déshérité son fils. »
\end{itemize}
\end{definition}

\begin{popfigure}
\begin{tikzpicture}[
    mindmap/.style={concept, execute at begin node=\hskip0pt},
    concept color=popblue!80,
    level 1/.style={level distance=3.5cm, sibling angle=90},
    level 2/.style={level distance=2.5cm, sibling angle=45},
    every node/.style={font=\small, text=white}
]
\node [mindmap] {das Erbe}
    child [concept color=poporange] {node {Personen}
        child {node {der Erblasser}}
        child {node {der Erbe / die Erbin}}
        child {node {der Notar}}
        child {node {der Pflichtteilsberechtigte}}
    }
    child [concept color=popgreen] {node {Objekte}
        child {node {das Haus / die Immobilie}}
        child {node {das Konto / das Geld}}
        child {node {der Schmuck / die Wertsachen}}
        child {node {das Grundstück}}
    }
    child [concept color=poppurple] {node {Probleme}
        child {node {der Streit}}
        child {node {uneindeutiges Testament}}
        child {node {die Enterbung}}
        child {node {hohe Kosten}}
    }
    child [concept color=poppink] {node {Lösungen}
        child {node {klare Regelung}}
        child {node {Testament / Erbvertrag}}
        child {node {Verkauf und Teilung}}
        child {node {Pflichtteil geltend machen}}
    };
\end{tikzpicture}
\legende{Carte mentale lexicale : le champ sémantique de \all{das Erbe}.}
\end{popfigure}

\begin{popfigure}
\begin{tikzpicture}[
    every node/.style={font=\small, align=center},
    male/.style={rectangle, draw=popblue, thick, fill=popblueL, rounded corners, minimum width=2.2cm, minimum height=1cm},
    female/.style={rectangle, draw=poppink, thick, fill=poppinkL, rounded corners, minimum width=2.2cm, minimum height=1cm},
    arrow/.style={-Stealth, thick, popdark}
]
% Generation 1
\node[male, align=center] (gf) at (0, 3){Großvater \\ (Erblasser, †)};
% Generation 2
\node[female, align=center] (mother) at (-4, 1){Mutter \\ (vermittelt)};
\node[male, align=center] (thomas) at (0, 1){Thomas \\ (Sohn, Erbe)};
\node[female, align=center] (anna) at (4, 1){Anna \\ (Tochter, Erbin)};
% Generation 3
\node[female] (enkelin) at (4, -1) {Enkelin};
% Arrows
\draw[arrow] (gf) -- (mother);
\draw[arrow] (gf) -- (thomas);
\draw[arrow] (gf) -- (anna);
\draw[arrow] (anna) -- (enkelin);
% Labels
\node[font=\bfseries\small, text=popdark] at (-6, 3) {1. Generation};
\node[font=\bfseries\small, text=popdark] at (-6, 1) {2. Generation};
\node[font=\bfseries\small, text=popdark] at (-6, -1) {3. Generation};
\end{tikzpicture}
\legende{Stammbaum simplifié de la famille du texte.}
\end{popfigure}

\begin{exoresolu}[Repérage des temps du récit dans le texte]
\textbf{Énoncé :} Relevez dans le texte \all{Der Streit um das Erbe} \textbf{6 formes du Präteritum} et \textbf{2 formes du Perfekt} (auxiliaire + participe II). Classez-les dans un tableau en indiquant l'infinitif et le type de verbe.
\tcblower
\textbf{Solution détaillée :}

\begin{center}
\begin{tabularx}{\linewidth}{|X|X|X|}\hline
\rowcolor{popblueL} \textbf{Forme trouvée} & \textbf{Infinitif} & \textbf{Temps / Type} \\ \hline
\all{starb} & \all{sterben} & Präteritum (fort, 3e p. sg.) \\ \hline
\all{hinterließ} & \all{hinterlassen} & Präteritum (fort, 3e p. sg.) \\ \hline
\all{war} & \all{sein} & Präteritum (irrégulier, 3e p. sg.) \\ \hline
\all{sollte} & \all{sollen} & Präteritum (verbe modal, 3e p. sg.) \\ \hline
\all{bestand} & \all{bestehen} & Präteritum (fort, 3e p. sg.) \\ \hline
\all{stritten} & \all{streiten} & Präteritum (fort, 3e p. pl.) \\ \hline
\all{versuchte} & \all{versuchen} & Präteritum (faible, 3e p. sg.) \\ \hline
\all{wurde} & \all{werden} & Präteritum (auxiliaire du passif, 3e p. sg.) \\ \hline
\all{waren} & \all{sein} & Präteritum (irrégulier, 3e p. pl.) \\ \hline
\hline
\all{hat ... geregelt} & \all{regeln} & Perfekt (\all{haben} + Part. II) \\ \hline
\all{hat ... verloren} & \all{verlieren} & Perfekt (\all{haben} + Part. II) \\ \hline
\hline
\all{hatte ... gehofft} & \all{hoffen} & Plusquamperfekt (\all{hatte} + Part. II) \\ \hline
\all{wurde ... verkauft} & \all{verkaufen} & Präteritum passif (\all{wurde} + Part. II) \\ \hline
\all{wurde ... aufgeteilt} & \all{aufteilen} & Präteritum passif (\all{wurde} + Part. II) \\ \hline
\end{tabularx}
\end{center}

\textbf{Observation :} Le \textbf{Präteritum} domine pour la narration des événements passés (\all{starb, hinterließ, stritten, war}) : c'est le temps du récit écrit (presse, littérature). Le \textbf{Perfekt} (\all{hat geregelt, hat verloren}) insiste sur le résultat encore sensible au présent ; c'est aussi le temps privilégié à l'oral. Le \textbf{Plusquamperfekt} (\all{hatte gehofft}) exprime l'antériorité : l'espoir du grand-père est \emph{antérieur} au conflit. Attention : \all{wurde verkauft} n'est pas un Perfekt, mais le \textbf{Präteritum du passif} (\all{werden} conjugué au Präteritum + participe II).
\end{exoresolu}

\coursec{Compréhension du texte}

Après la lecture du texte support \all{Der Streit um das Erbe} (section précédente), nous passons maintenant à l'étape essentielle de toute leçon de texte : la \cle{compréhension}. Lire un texte allemand ne suffit pas ; il faut savoir répondre à des questions précises, justifier ses réponses par le texte et donner son opinion. C'est exactement ce qui est demandé au baccalauréat. Nous allons apprendre une méthode rigoureuse, puis l'appliquer pas à pas.

Pour travailler de manière autonome, voici d'abord un court texte parallèle, sur le même thème, qui servira de support d'entraînement tout au long de cette section.

\begin{textebox}[Ein Erbe in Douala]
Herr Mbappe, ein bekannter Händler aus Douala, ist letzten Monat gestorben. Er hat ein großes Erbe hinterlassen: ein Haus in Bonanjo, zwei Geschäfte am Markt und viel Geld auf der Bank. Leider hat er kein Testament geschrieben. Deshalb gab es sofort einen Streit in der Familie. Sein ältester Sohn Jean wollte das Haus haben, seine Tochter Marie wollte die Geschäfte, und seine Frau wollte das Geld für die Enkelkinder behalten. Die Erben stritten drei Monate lang. Sie sprachen nicht mehr miteinander, und die Familie war sehr traurig. Dann kam der Onkel des Verstorbenen, ein alter und weiser Mann. Er lud alle Erben zu einem Gespräch ein. „Das Gesetz ist klar“, sagte er, „aber wir können auch eine gerechte Lösung finden.“ Nach langen Diskussionen einigten sich die Erben: Jean bekam das Haus, Marie die Geschäfte, und das Geld wurde unter der Frau und den Enkelkindern geteilt. Am Ende umarmten sich alle, und die Familie feierte zusammen das Gedächtnis des Vaters. So wurde die Succession geregelt, ohne dass ein Gericht nötig war. Heute sagen alle: „Der Onkel hat uns gerettet.“
\end{textebox}

\textbf{Glossaire :} der Händler, - (le commerçant) ; das Erbe, -n (l'héritage) ; hinterlassen (laisser, léguer) ; das Testament, -e (le testament) ; der Streit, -e (la dispute) ; die Erben (les héritiers) ; der Verstorbene, -n (le défunt) ; sich einigen (se mettre d'accord) ; teilen (partager) ; das Gericht, -e (le tribunal) ; die Lösung, -en (la solution).

\courssub{Compréhension globale : de quoi parle le texte ?}

La première lecture sert à répondre à deux questions simples : \emph{de quoi parle le texte ?} et \emph{qui sont les personnages ?}. En allemand :

\begin{itemize}
\item \all{Worum geht es in dem Text?} — De quoi parle le texte ?
\item \all{Wer sind die Personen?} — Qui sont les personnages ?
\item \all{Wo und wann spielt die Geschichte?} — Où et quand se déroule l'histoire ?
\end{itemize}

\begin{exemplebox}[Réponses globales]
\all{Der Text handelt von einem Erbe und einem Familienstreit in Douala.} — Le texte parle d'un héritage et d'une dispute familiale à Douala.

\all{Die Personen sind Herr Mbappe, sein Sohn Jean, seine Tochter Marie, seine Frau und der Onkel.} — Les personnages sont M. Mbappe, son fils Jean, sa fille Marie, sa femme et l'oncle.
\end{exemplebox}

Notez la construction \all{es geht um + Akkusativ} (il s'agit de) : \all{In dem Text geht es um ein Erbe.} — Il s'agit d'un héritage dans le texte.

\courssub{La méthode de réponse à une question de compréhension}

Avant de répondre aux questions détaillées, apprenons la technique professionnelle du germaniste.

\begin{methode}[Répondre à une question de compréhension en quatre étapes]
\begin{enumerate}
\item \textbf{Lire la question} et souligner le mot interrogatif (\all{wer}, \all{was}, \all{warum}, \all{wann}, \all{wo}).
\item \textbf{Repérer la ligne du texte} (\all{die Textstelle finden}) qui contient la réponse : cherchez les mots-clés de la question dans le texte.
\item \textbf{Formuler une phrase complète} (\all{ein vollständiger Satz}) en reprenant les éléments de la question : sujet + verbe en deuxième position. Ne répondez jamais par un seul mot.
\item \textbf{Justifier par une courte citation} entre guillemets allemands „ … “ si la question le demande (\all{Belegen Sie Ihre Antwort mit einem Zitat}).
\end{enumerate}
\end{methode}

\begin{popfigure}
\begin{center}
\begin{tikzpicture}[
node distance=0.55cm,
box/.style={draw=popblue, fill=popblueL, rounded corners, align=center, text width=4.6cm, minimum height=0.9cm, font=\small},
arr/.style={-{Stealth[length=3mm]}, thick, poporange}]
\node[box, align=center] (q){\textbf{1. Question lesen}\\ \all{W-Frage unterstreichen}};
\node[box, below=of q, align=center] (t){\textbf{2. Textstelle finden}\\ Stichwörter suchen};
\node[box, below=of t, align=center] (s){\textbf{3. Vollständiger Satz}\\ Subjekt + Verb an Position 2};
\node[box, below=of s, align=center] (z){\textbf{4. Zitat}\\ kurz, zwischen „ … “};
\draw[arr] (q) -- (t);
\draw[arr] (t) -- (s);
\draw[arr] (s) -- (z);
\end{tikzpicture}
\end{center}
\legende{Organigramme de la méthode de réponse : de la question à la citation}
\end{popfigure}

\courssub{Compréhension détaillée : les questions classiques}

Appliquons la méthode aux questions types de la fiche de progression.

\begin{exemplebox}[Questions et réponses modèles]
\textbf{1.} \all{Wer ist gestorben?} — Qui est mort ?\\
\all{Herr Mbappe, ein Händler aus Douala, ist gestorben.} — M. Mbappe, un commerçant de Douala, est mort.

\textbf{2.} \all{Was hat er hinterlassen?} — Qu'a-t-il laissé ?\\
\all{Er hat ein Haus, zwei Geschäfte und viel Geld hinterlassen.} — Il a laissé une maison, deux boutiques et beaucoup d'argent.

\textbf{3.} \all{Warum gab es einen Streit?} — Pourquoi y a-t-il eu une dispute ?\\
\all{Es gab einen Streit, weil der Verstorbene kein Testament geschrieben hatte.} — Il y a eu une dispute parce que le défunt n'avait pas écrit de testament.

\textbf{4.} \all{Wer hat das Problem gelöst?} — Qui a résolu le problème ?\\
\all{Der Onkel des Verstorbenen hat das Problem gelöst.} — L'oncle du défunt a résolu le problème.
\end{exemplebox}

Observez la question 3 : la réponse à une question en \all{warum} se construit avec \all{weil}, et le verbe conjugué va \emph{à la fin} de la subordonnée.

\begin{propriete}[La réponse à une question avec \all{warum}]
Une question introduite par \all{warum} (pourquoi) exige une réponse avec \all{weil} (parce que). \all{Weil} est une conjonction de subordination : le verbe conjugué se place \textbf{en dernière position} de la proposition.

\all{Warum stritten die Erben? — Die Erben stritten, weil das Testament nicht klar war.} — Pourquoi les héritiers se disputaient-ils ? — Les héritiers se disputaient parce que le testament n'était pas clair.

Forme : \all{weil + sujet + … + verbe conjugué (fin)}. Au parfait, le participe précède l'auxiliaire : \all{…, weil er kein Testament geschrieben \textbf{hat}.}
\end{propriete}

\begin{center}
\begin{tabularx}{\linewidth}{|X|X|X|}
\hline
\rowcolor{popblueL} Français & Deutsch & Remarque \\
\hline
Pourquoi… ? & \all{Warum… ?} & mot interrogatif en tête, verbe en 2\textsuperscript{e} position \\
\hline
parce que + verbe conjugué en 2\textsuperscript{e} position & \all{weil} + verbe à la fin & différence majeure avec le français \\
\hline
car (phrase principale) & \all{denn} + verbe en 2\textsuperscript{e} position & \all{denn} ne déplace pas le verbe \\
\hline
citer le texte & \all{den Text zitieren} & citation courte entre „ … “ \\
\hline
justifier sa réponse & \all{die Antwort begründen} & verbe faible : \all{begründete} \\
\hline
\end{tabularx}
\end{center}

\begin{attention}[Ne recopiez pas le texte !]
Erreur fréquente des francophones : recopier trois ou quatre lignes du texte comme réponse. Le correcteur veut \emph{votre} phrase. Citez seulement quelques mots entre guillemets : \all{Der Text sagt: „kein Testament geschrieben“ — deshalb gab es einen Streit.} Autre piège : oublier la majuscule aux noms (\all{das Testament}, \all{der Streit}) et placer le verbe en deuxième position après \all{weil} — non ! Après \all{weil}, le verbe va à la fin.
\end{attention}

\courssub{Vrai ou faux : justifier par une citation}

L'exercice \all{Richtig oder falsch?} demande de juger des affirmations et de prouver son jugement par une citation du texte.

\begin{exoresolu}[Richtig oder falsch ?]
Énoncé. Jugez les affirmations suivantes (richtig / falsch) et justifiez par une courte citation du texte \all{Ein Erbe in Douala}.\\
a) \all{Herr Mbappe hat ein Testament geschrieben.}\\
b) \all{Die Familie stritt drei Monate lang.}\\
c) \all{Ein Gericht hat die Succession geregelt.}
\tcblower
Solution détaillée.\\
a) \textbf{Falsch.} Le texte dit : \all{„Leider hat er kein Testament geschrieben.“} — Hélas, il n'a pas écrit de testament.\\
b) \textbf{Richtig.} Le texte dit : \all{„Die Erben stritten drei Monate lang.“} — Les héritiers se sont disputés pendant trois mois.\\
c) \textbf{Falsch.} Le texte dit : \all{„So wurde die Succession geregelt, ohne dass ein Gericht nötig war.“} — La succession a été réglée sans qu'un tribunal soit nécessaire.
\end{exoresolu}

\courssub{Trouver les équivalents allemands dans le texte}

Compétence clé de la version et du thème : retrouver dans le texte l'équivalent allemand d'une expression française. Cherchez le sens, pas la traduction mot à mot.

\begin{center}
\begin{tabularx}{\linewidth}{|X|X|}
\hline
\rowcolor{popblueL} Français & Équivalent dans le texte \\
\hline
laisser un héritage & \all{ein Erbe hinterlassen} \\
\hline
régler une succession & \all{die Succession regeln} \\
\hline
partager l'argent & \all{das Geld teilen} \\
\hline
se disputer (pendant trois mois) & \all{(drei Monate lang) streiten} \\
\hline
se mettre d'accord & \all{sich einigen} \\
\hline
trouver une solution juste & \all{eine gerechte Lösung finden} \\
\hline
\end{tabularx}
\end{center}

Remarquez le verbe à particule séparable \all{hinterlassen} : c'est un verbe \emph{inséparable} (participe \all{hinterlassen}, sans \all{ge-}) : \all{Er hat ein Erbe hinterlassen.} En revanche \all{sich einigen} est réfléchi : \all{Die Erben einigten sich.}

\courssub{Questions d'opinion : s'exprimer personnellement}

Les questions d'opinion (\all{Meinungsfragen}) n'ont pas de réponse unique : elles évaluent votre capacité à produire un allemand correct et argumenté.

\begin{exemplebox}[Donner son avis]
\textbf{1.} \all{Findest du die Lösung gerecht?} — Trouves-tu la solution juste ?\\
\all{Ja, ich finde die Lösung gerecht, weil jeder Erbe etwas bekommen hat.} — Oui, je trouve la solution juste parce que chaque héritier a reçu quelque chose.

\textbf{2.} \all{Was würdest du machen?} — Que ferais-tu ?\\
\all{Ich würde ein Testament schreiben, um einen Streit zu vermeiden.} — J'écrirais un testament pour éviter une dispute.
\end{exemplebox}

Notez le conditionnel (\all{Konjunktiv II}) : \all{ich würde + infinitif} pour exprimer une hypothèse, et les tournures d'opinion : \all{Ich finde …}, \all{Meiner Meinung nach …} (à mon avis), \all{Ich bin der Meinung, dass …}.

\begin{exercice}[À vous de jouer]
Répondez par des phrases complètes en allemand au texte \all{Ein Erbe in Douala}.\\
1. \all{Wo spielt die Geschichte?}\\
2. \all{Warum war die Familie traurig?}\\
3. Richtig oder falsch (avec citation) : \all{Marie bekam das Haus.}\\
4. Trouvez dans le texte l'équivalent de : « inviter tous les héritiers à une discussion ».\\
5. \all{Findest du den Onkel weise? Begründe deine Antwort.}
\end{exercice}

\begin{corrige}[Exercice : À vous de jouer]
1. \all{Die Geschichte spielt in Douala, in Kamerun.} — L'histoire se déroule à Douala, au Cameroun.\\
2. \all{Die Familie war traurig, weil die Erben nicht mehr miteinander sprachen.} — La famille était triste parce que les héritiers ne se parlaient plus. (Verbe à la fin après \all{weil}.)\\
3. \textbf{Falsch.} Citation : \all{„Jean bekam das Haus, Marie die Geschäfte.“} — C'est Jean qui a reçu la maison, Marie les boutiques.\\
4. \all{„Er lud alle Erben zu einem Gespräch ein.“} — Il a invité tous les héritiers à une discussion. (Verbe séparable \all{einladen} : \all{lud … ein}.)\\
5. Réponse possible : \all{Ja, ich finde den Onkel weise, weil er eine gerechte Lösung gefunden hat und die Familie gerettet hat.} — Oui, je trouve l'oncle sage parce qu'il a trouvé une solution juste et sauvé la famille.
\end{corrige}

\begin{aretenir}[L'essentiel de la compréhension de texte]
\begin{itemize}
\item Toujours répondre par une \textbf{phrase complète} qui reprend les éléments de la question.
\item \all{Warum} appelle \all{weil} + verbe conjugué \textbf{à la fin}.
\item Justifier par une \textbf{citation courte} entre guillemets allemands „ … “, jamais par un long passage recopié.
\item Pour les équivalents, chercher le \textbf{sens} dans le texte, pas le mot à mot.
\item Pour l'opinion : \all{Ich finde …}, \all{Meiner Meinung nach …}, \all{ich würde …} + argument avec \all{weil}.
\end{itemize}
\end{aretenir}

Nous savons désormais lire, comprendre et exploiter le texte. La section suivante approfondira le lexique thématique de l'héritage : \all{das Erbe}, \all{der Erbe}, \all{die Erbin}, \all{das Testament}, \all{der Streit}, \all{das Gesetz}.

\coursec{Lexique thématique : Erbe, Testament, Familie, Gesetz}

Après avoir lu et compris le texte \all{Der Streit um das Erbe}, nous avons rencontré de nombreux mots nouveaux : \all{das Erbe}, \all{das Testament}, \all{der Notar}, \all{sich streiten}... Cette section a pour but de \cle{maîtriser ce vocabulaire} de façon systématique : connaître chaque mot avec son \cle{article} et son \cle{pluriel}, comprendre les \cle{familles de mots}, savoir employer les \cle{expressions} dans des phrases correctes. Un bon lexique, c'est la moitié du travail pour parler et écrire sur les problèmes d'héritage.

\courssub{Observation : le vocabulaire dans son contexte}

Avant d'apprendre des listes de mots, observons-les d'abord dans un texte authentique. Lisez le texte suivant et repérez tous les mots qui appartiennent au thème de l'héritage.

\begin{textebox}[Ein Testament bringt Ruhe]
Herr Mbarga, ein bekannter Kaufmann aus Yaoundé, hatte drei Kinder: zwei Söhne und eine Tochter. Er besaß ein großes Haus, ein Geschäft im Zentrum und ein Stück Land vor der Stadt. Weil er schon viele Konflikte in anderen Familien gesehen hatte, wollte er nach seinem Tod keinen Streit zwischen seinen Kindern. Deshalb ging er eines Tages zu einem Notar und machte ein Testament. Darin schrieb er genau, wer was bekommen sollte: Der älteste Sohn sollte das Geschäft erben, die Tochter das Haus und der jüngste Sohn das Land. Außerdem hinterließ er jedem Kind eine kleine Summe Geld.

Als Herr Mbarga starb, las der Notar das Testament vor der ganzen Familie vor. Die Verwandten kamen aus Douala und Bafoussam. Zuerst waren alle traurig, aber niemand war überrascht: Die Regelung war klar und gerecht. Ein Cousin, der auch etwas bekommen wollte, protestierte. „Ich bin auch ein Verwandter!“, rief er. Aber der Anwalt der Familie erklärte ihm das Gesetz: Nur die Kinder haben einen Pflichtteil. Der Cousin musste schweigen.

Die Geschwister stritten nicht. Sie teilten alles so, wie der Vater es gewollt hatte. Die Tochter sagte später: „Unser Vater hat uns nicht nur ein Erbe hinterlassen, er hat uns auch den Frieden hinterlassen.“ Heute erinnern sich die Enkel oft an ihren Großvater. Sie wissen: Ein gutes Testament schützt die Familie.
\end{textebox}

\textbf{Glossaire des mots difficiles :}

\begin{center}
\begin{tabularx}{\linewidth}{|X|X|}
\hline
\rowcolor{popblueL} \textbf{Deutsch} & \textbf{Français} \\
\hline
der Kaufmann, die Kaufleute & le commerçant \\
\hline
das Geschäft, die Geschäfte & le magasin, le commerce \\
\hline
das Stück Land & le terrain, la parcelle \\
\hline
der Notar, die Notare & le notaire \\
\hline
die Regelung, die Regelungen & le règlement, l'arrangement \\
\hline
gerecht & juste, équitable \\
\hline
der Pflichtteil, die Pflichtteile & la part réservataire (part légale obligatoire) \\
\hline
der Frieden & la paix \\
\hline
schützen & protéger \\
\hline
\end{tabularx}
\end{center}

\textbf{Questions de compréhension rapide :}
\begin{enumerate}
\item \all{Warum macht Herr Mbarga ein Testament?} (Pourquoi M. Mbarga fait-il un testament ?)
\item \all{Wer erbt das Haus?} (Qui hérite de la maison ?)
\item \all{Warum bekommt der Cousin nichts?} (Pourquoi le cousin ne reçoit-il rien ?)
\end{enumerate}

\courssub{Le champ lexical de l'héritage}

\begin{definition}[das Erbe]
\all{das Erbe, die Erben} signifie « l'héritage » (ce qu'on reçoit d'un défunt : biens, argent, maison). Le mot \all{die Erbschaft, die Erbschaften} est un synonyme plus juridique. Le verbe \all{erben} (faible) signifie « hériter (de) », et \all{vererben} (faible) signifie « léguer, transmettre par héritage ». Le verbe \all{hinterlassen} (fort : \all{hinterließ}, \all{hat hinterlassen}) signifie « laisser (après sa mort) ».
\end{definition}

\begin{exemplebox}[Les verbes en action]
\begin{itemize}
\item \all{Die Tochter erbt das Haus.} « La fille hérite de la maison. »
\item \all{Der Großvater hat seiner Enkelin das Haus vererbt.} « Le grand-père a légué la maison à sa petite-fille. »
\item \all{Er hat seiner Familie ein großes Vermögen hinterlassen.} « Il a laissé une grande fortune à sa famille. »
\item \all{Wem hat der Onkel sein Geld vererbt?} « À qui l'oncle a-t-il légué son argent ? »
\end{itemize}
\end{exemplebox}

\begin{attention}[erben ou vererben ?]
Les francophones confondent souvent les deux verbes. Retenez la direction de l'action : \all{erben} = \textbf{recevoir} (l'héritier hérite), \all{vererben} = \textbf{donner} (le défunt lègue). La personne qui reçoit est au \textbf{datif} avec \all{vererben} : \all{Er vererbt \emph{seiner Tochter} (Dat.) \emph{das Haus} (Akk.).} « Il lègue la maison \emph{à sa fille}. »
\end{attention}

\courssub{Les personnes}

Voici les acteurs de toute histoire d'héritage. Apprenez chaque nom avec son article et son pluriel.

\begin{center}
\begin{tabularx}{\linewidth}{|X|X|X|}
\hline
\rowcolor{popblueL} \textbf{Deutsch (Singulier / Pluriel)} & \textbf{Français} & \textbf{Exemple} \\
\hline
der Erbe / die Erben & l'héritier & \all{Der Erbe bekommt alles.} \\
\hline
die Erbin / die Erbinnen & l'héritière & \all{Sie ist die einzige Erbin.} \\
\hline
der Verstorbene / die Verstorbenen & le défunt & \all{Der Verstorbene hatte drei Kinder.} \\
\hline
der Notar / die Notare & le notaire & \all{Der Notar liest das Testament vor.} \\
\hline
der Anwalt / die Anwälte & l'avocat & \all{Der Anwalt erklärt das Gesetz.} \\
\hline
die Familie / die Familien & la famille & \all{Die Familie kommt zusammen.} \\
\hline
die Geschwister (Pl.) & les frères et sœurs & \all{Die Geschwister teilen das Erbe.} \\
\hline
der Enkel / die Enkel & le petit-fils & \all{Der Enkel erbt das Land.} \\
\hline
die Enkelin / die Enkelinnen & la petite-fille & \all{Die Enkelin bekommt das Haus.} \\
\hline
die Verwandten (Pl.) & les parents (famille élargie) & \all{Alle Verwandten sind gekommen.} \\
\hline
\end{tabularx}
\end{center}

\begin{propriete}[La formation du féminin : -in]
En allemand, la plupart des noms de personnes masculins forment leur féminin avec le suffixe \cle{-in}, et le pluriel féminin en \cle{-nen} : \all{der Erbe} → \all{die Erbin} → \all{die Erbinnen} ; \all{der Enkel} → \all{die Enkelin} → \all{die Enkelinnen} ; \all{der Notar} → \all{die Notarin} → \all{die Notarinnen}. C'est beaucoup plus régulier qu'en français (héritier/héritière, petit-fils/petite-fille).
\end{propriete}

\courssub{Les documents et le droit}

\begin{definition}[das Testament]
\all{das Testament, die Testamente} : document écrit dans lequel une personne décide de la répartition de ses biens après sa mort. Expression clé : \all{ein Testament machen} « faire un testament ».
\end{definition}

\begin{center}
\begin{tabularx}{\linewidth}{|X|X|X|}
\hline
\rowcolor{popblueL} \textbf{Deutsch (Singulier / Pluriel)} & \textbf{Français} & \textbf{Remarque} \\
\hline
das Testament / die Testamente & le testament & \all{ein Testament machen} \\
\hline
das Gesetz / die Gesetze & la loi (texte juridique) & \all{nach dem Gesetz} « selon la loi » \\
\hline
das Recht / die Rechte & le droit & \all{das Erbrecht} « le droit des successions » \\
\hline
der Pflichtteil / die Pflichtteile & la part réservataire & part légale garantie aux proches \\
\hline
die Regelung / die Regelungen & le règlement, la solution & \all{eine klare Regelung} \\
\hline
\end{tabularx}
\end{center}

\begin{attention}[das Gesetz ≠ das Recht]
Piège classique pour les francophones, car le français « droit » et « loi » se recouvrent partiellement. \all{das Gesetz} = la \textbf{loi}, c'est-à-dire un texte précis voté et publié (\all{ein neues Gesetz} « une nouvelle loi »). \all{das Recht} = le \textbf{droit}, c'est-à-dire l'ensemble du système juridique ou un droit subjectif : \all{das Erbrecht} « le droit d'hériter / le droit des successions », \all{Ich habe ein Recht darauf.} « J'y ai droit. »
\end{attention}

\courssub{Le conflit et la solution}

\begin{center}
\begin{tabularx}{\linewidth}{|X|X|X|}
\hline
\rowcolor{popblueL} \textbf{Deutsch (Singulier / Pluriel)} & \textbf{Français} & \textbf{Exemple} \\
\hline
der Streit / die Streite & la dispute, le conflit & \all{ein Streit um das Erbe} \\
\hline
sich streiten (stritt, hat gestritten) & se disputer & \all{Die Brüder streiten um das Geld.} \\
\hline
der Konflikt / die Konflikte & le conflit & \all{ein langer Konflikt} \\
\hline
teilen (faible) & partager & \all{Sie teilen das Erbe.} \\
\hline
verkaufen (faible, inséparable) & vendre & \all{Sie verkaufen das Haus.} \\
\hline
vermitteln (faible, inséparable) & servir de médiateur & \all{Der Notar vermittelt.} \\
\hline
gerecht / ungerecht & juste / injuste & \all{eine gerechte Lösung} \\
\hline
\end{tabularx}
\end{center}

\begin{exemplebox}[Exemples traduits]
\begin{itemize}
\item \all{Die Brüder streiten um das Geld.} « Les frères se disputent l'argent. »
\item \all{Nach dem Tod des Vaters gab es einen langen Streit.} « Après la mort du père, il y a eu une longue dispute. »
\item \all{Die Kinder verkaufen das Haus und teilen das Geld.} « Les enfants vendent la maison et se partagent l'argent. »
\item \all{Die Lösung ist nicht gerecht.} « La solution n'est pas juste. »
\end{itemize}
\end{exemplebox}

\begin{attention}[bekommen ≠ devenir]
Faux ami célèbre ! \all{bekommen} signifie « \textbf{recevoir, obtenir} », jamais « devenir ». \all{Die Tochter bekommt das Haus.} = « La fille reçoit la maison. » Pour dire « devenir », l'allemand utilise \all{werden} : \all{Sie wird Ärztin.} « Elle devient médecin. »
\end{attention}

\begin{propriete}[sich streiten : un verbe fort réfléchi]
\all{sich streiten} est un verbe fort : \all{streiten – stritt – hat gestritten}. Il se construit avec \all{um} + accusatif pour l'objet de la dispute et \all{mit} + datif pour la personne : \all{Ich habe mit meinem Bruder um das Erbe gestritten.} « Je me suis disputé avec mon frère à propos de l'héritage. »
\end{propriete}

\courssub{Familles de mots : un radical, plusieurs mots}

L'allemand construit des familles entières autour d'un radical. Connaître le radical, c'est connaître dix mots à la fois.

\begin{center}
\begin{tabularx}{\linewidth}{|X|X|X|}
\hline
\rowcolor{popgreenL} \textbf{Radical} & \textbf{Mot} & \textbf{Français} \\
\hline
\multirow{5}{*}{erb-} & erben & hériter \\
 & der Erbe / die Erbin & l'héritier / l'héritière \\
 & das Erbe & l'héritage \\
 & die Erbschaft & la succession \\
 & erblich & héréditaire \\
\hline
\multirow{2}{*}{streit-} & streiten / der Streit & disputer / la dispute \\
 & streitsüchtig & querelleur (qui aime se disputer) \\
\hline
\end{tabularx}
\end{center}

\begin{savaistu}[Le saviez-tu ?]
Le mot \all{Erbe} existe aussi dans \all{die Erbkrankheit} « la maladie héréditaire » — même racine \all{erb-}, autre domaine ! De même, on trouve \all{die Erbsubstanz} en biologie (la matière héréditaire, l'ADN). Quand vous rencontrez un mot allemand inconnu commençant par \all{Erb-}, vous savez déjà qu'il parle de transmission. C'est la force des familles de mots allemandes.
\end{savaistu}

\courssub{Carte mentale du vocabulaire}

\begin{popfigure}
\begin{center}
\begin{tikzpicture}[mindmap, grow cyclic,
  every node/.style={concept, align=center, text=white},
  root concept/.append style={concept color=popblue, font=\Large\bfseries},
  level 1/.append style={level distance=3.4cm, sibling angle=90},
  level 2/.append style={level distance=2.6cm, sibling angle=45, font=\small},
  concept connection/.append style={color=popmuted}]
\node[root concept]{ERBEN}
  child[concept color=poporange]{ node {Personen}
    child { node[concept color=poporange!70, align=center]{der Erbe\\die Erbin} }
    child { node[concept color=poporange!70, align=center]{der Notar\\der Anwalt} }
  }
  child[concept color=popgreen]{ node {Dokumente}
    child { node[concept color=popgreen!70, align=center]{das\\Testament} }
    child { node[concept color=popgreen!70, align=center]{das Gesetz\\das Recht} }
  }
  child[concept color=poppurple]{ node {Konflikt}
    child { node[concept color=poppurple!70, align=center]{der Streit\\sich streiten} }
    child { node[concept color=poppurple!70] {der Konflikt} }
  }
  child[concept color=poppink]{ node {Lösung}
    child { node[concept color=poppink!70, align=center]{teilen\\verkaufen} }
    child { node[concept color=poppink!70, align=center]{vermitteln\\gerecht} }
  };
\end{tikzpicture}
\end{center}
\legende{Carte mentale du champ lexical de l'héritage : quatre branches à mémoriser.}
\end{popfigure}

\courssub{Expressions utiles à réemployer}

Ces expressions figées sont indispensables pour l'expression écrite et orale. Apprenez-les comme des blocs.

\begin{center}
\begin{tabularx}{\linewidth}{|X|X|}
\hline
\rowcolor{popgoldL} \textbf{Deutsch} & \textbf{Français} \\
\hline
ein Testament machen & faire un testament \\
\hline
um das Erbe streiten & se disputer à propos de l'héritage \\
\hline
das Erbe teilen & partager l'héritage \\
\hline
jmdm. etwas vererben & léguer quelque chose à quelqu'un \\
\hline
jmdm. etwas hinterlassen & laisser quelque chose à quelqu'un (à sa mort) \\
\hline
eine gerechte Lösung finden & trouver une solution juste \\
\hline
nach dem Gesetz & selon la loi \\
\hline
\end{tabularx}
\end{center}

\begin{attention}[jmdm. = datif !]
Dans \all{jmdm. etwas vererben}, la personne (\all{jmdm.} = jemandem) est au \textbf{datif} et la chose (\all{etwas}) à l'\textbf{accusatif} : \all{Er hat \emph{seinem Enkel} (Dat.) \emph{das Land} (Akk.) vererbt.} Ne dites jamais \all{Er hat seinen Enkel das Haus vererbt} — cela signifierait « il a légué son petit-fils à la maison » !
\end{attention}

\courssub{Exercice résolu et entraînement}

\begin{exoresolu}[Complétez avec le mot qui convient]
Énoncé — Complétez les phrases avec : \all{Testament}, \all{Erbin}, \all{streiten}, \all{vererbt}, \all{Gesetz}, \all{Pflichtteil}.
\begin{enumerate}
\item Der Großvater hat ein \dots\ beim Notar gemacht.
\item Meine Tante ist die einzige \dots\ : sie bekommt alles.
\item Die Geschwister \dots\ seit Monaten um das Haus.
\item Der Vater hat seiner Tochter das Geschäft \dots.
\item Nach dem \dots\ haben die Kinder einen \dots.
\end{enumerate}
\tcblower
Solution détaillée :
\begin{enumerate}
\item \all{Testament} — on « fait » un testament chez le notaire : \all{ein Testament machen}.
\item \all{Erbin} — « tante » est féminin, donc le féminin de \all{der Erbe} : \all{die Erbin}.
\item \all{streiten} — sujet pluriel (\all{die Geschwister}), présent : \all{sie streiten um} + accusatif.
\item \all{vererbt} — le père \emph{donne} (et ne reçoit pas) : participe passé de \all{vererben}, verbe faible : \all{hat vererbt}.
\item \all{Gesetz} / \all{Pflichtteil} — \all{nach dem Gesetz} « selon la loi » ; la part légale obligatoire des enfants s'appelle \all{der Pflichtteil}.
\end{enumerate}
\end{exoresolu}

\begin{exercice}[À vous de jouer]
\begin{enumerate}
\item Traduisez en allemand : a) « Le notaire lit le testament devant la famille. » b) « Les sœurs se disputent l'argent. » c) « Mon grand-père m'a légué sa maison. »
\item Formez le féminin et le pluriel : \all{der Erbe}, \all{der Enkel}, \all{der Notar}.
\item Construisez trois phrases avec : \all{das Erbe teilen}, \all{ein Testament machen}, \all{gerecht}.
\end{enumerate}
\end{exercice}

\begin{corrige}[Exercice « À vous de jouer »]
\begin{enumerate}
\item a) \all{Der Notar liest das Testament vor der Familie vor.} (Attention : \all{vorlesen} est un verbe à particule séparable, \all{vor} va en fin de phrase.) b) \all{Die Schwestern streiten um das Geld.} c) \all{Mein Großvater hat mir sein Haus vererbt.} (\all{mir} = datif.)
\item \all{der Erbe} → \all{die Erbin}, Pl. \all{die Erbinnen} ; \all{der Enkel} → \all{die Enkelin}, Pl. \all{die Enkelinnen} ; \all{der Notar} → \all{die Notarin}, Pl. \all{die Notarinnen}.
\item Exemples : \all{Die Kinder teilen das Erbe gerecht.} / \all{Mein Onkel macht ein Testament beim Notar.} / \all{Die Regelung ist gerecht.}
\end{enumerate}
\end{corrige}

\begin{aretenir}[L'essentiel du lexique]
\begin{itemize}
\item \all{erben} = recevoir ; \all{vererben} = donner (personne au datif).
\item \all{das Gesetz} = la loi ; \all{das Recht} = le droit (\all{das Erbrecht}).
\item \all{bekommen} = recevoir (jamais « devenir »).
\item Féminin en \all{-in}, pluriel en \all{-nen} : \all{die Erbin, die Erbinnen}.
\item Expressions clés : \all{ein Testament machen}, \all{um das Erbe streiten}, \all{das Erbe teilen}, \all{jmdm. etwas vererben}.
\end{itemize}
\end{aretenir}

Ce vocabulaire est maintenant votre boîte à outils. Dans la section suivante, nous verrons comment la grammaire — et d'abord le \cle{génitif} — permet d'exprimer précisément la possession : \all{das Haus \emph{des Großvaters}}, \all{die Kinder \emph{der Verstorbenen}}...

\coursec{Grammaire 1 : le génitif}

Dans la section précédente, nous avons découvert le lexique de l'héritage : \all{das Erbe}, \all{der Erbe}, \all{die Erbin}, \all{das Testament}, \all{der Streit}, \all{das Gesetz}. Or, dans le texte \all{Der Streit um das Erbe}, ces mots n'apparaissent presque jamais seuls : on lit \all{das Haus der Familie}, \all{der Sohn des Großvaters}, \all{die Mutter der Kinder}. Que remarque-t-on ? Les articles changent de forme (\all{der} devient \all{des}, \all{die} devient \all{der}) et les noms masculins et neutres prennent une terminaison \all{-s} ou \all{-es}. C'est exactement le phénomène que nous allons étudier : le \cle{génitif}, le cas de la possession et du complément du nom. C'est le cas roi de tout texte juridique ou familial sur l'héritage.

\courssub{Observation : le génitif dans le texte}

\begin{textebox}[Mini-texte d'observation : Nach dem Tod des Großvaters]
Nach dem Tod des Großvaters begann der Streit der Familie. Das Testament des Verstorbenen war verschwunden. Der Sohn des Großvaters, Onkel Karl, behauptete, das Haus der Familie gehöre ihm. Die Mutter der Kinder, Tante Rosa, zeigte den Anwälten eine Kopie des Testaments. Darin stand der Wille des Verstorbenen: Das Haus des Großvaters sollte an die Kinder der Erbin gehen. Der Bruder des Notars bestätigte die Echtheit des Dokuments. Am Ende akzeptierten alle Mitglieder der Familie die Entscheidung des Richters. Der Frieden der Familie war wichtiger als das Geld des Erbes.
\end{textebox}

\textbf{Glossaire :}\par
\begin{itemize}
\item \all{der Großvater, die Großväter} — le grand-père
\item \all{verschwinden} — disparaître (participe \all{verschwunden} = disparu)
\item \all{behaupten} — prétendre, affirmer
\item \all{der Anwalt, die Anwälte} — l'avocat
\item \all{die Kopie, -n} — la copie
\item \all{der Wille, -n} — la volonté
\item \all{bestätigen} — confirmer
\item \all{die Echtheit} — l'authenticité
\item \all{der Richter, -} — le juge
\item \all{der Frieden} — la paix
\end{itemize}

\textbf{Questions d'observation :}\par
\begin{enumerate}
\item Souligne dans le texte tous les groupes de mots au génitif (article \all{des} ou \all{der} + nom).
\item Que devient le nom \all{der Großvater} après \all{das Haus} ? Et \all{der Verstorbene} après \all{das Testament} ?
\item À quelle question française répondent ces groupes : « de qui ? », « de quoi ? »
\end{enumerate}

\textbf{Observation :}\par
Observation : dans \all{das Haus des Großvaters}, l'article \all{der} est devenu \all{des} et le nom a pris un \all{-s} : \all{Großvaters}. Dans \all{die Mutter der Kinder}, l'article \all{die} (pluriel) est devenu \all{der}, sans terminaison sur le nom. Le génitif marque donc à la fois l'\cle{article} et parfois le \cle{nom}.

\courssub{Formation du génitif}

\begin{propriete}[La règle d'or du génitif]
Au génitif singulier, les noms \textbf{masculins} et \textbf{neutres} prennent la désinence \cle{-(e)s} :
\begin{itemize}
\item \all{der Vater} $\rightarrow$ \all{des Vaters} (du père)
\item \all{das Kind} $\rightarrow$ \all{des Kindes} (de l'enfant)
\item \all{der Notar} $\rightarrow$ \all{des Notars} (du notaire)
\end{itemize}
Les noms \textbf{féminins} et les noms au \textbf{pluriel} ne prennent \textbf{aucune désinence} : seul l'article change (\all{die} $\rightarrow$ \all{der}) :
\begin{itemize}
\item \all{die Mutter} $\rightarrow$ \all{der Mutter} (de la mère)
\item \all{die Kinder} $\rightarrow$ \all{der Kinder} (des enfants)
\end{itemize}
Choix entre \all{-s} et \all{-es} : on ajoute \all{-es} aux noms d'une syllabe et aux noms se terminant par une consonne dure (\all{-s, -ß, -z, -sch, -st}) : \all{des Kindes}, \all{des Hauses}, \all{des Gesetzes}. On ajoute \all{-s} aux noms de plusieurs syllabes et aux noms en voyelle : \all{des Vaters}, \all{des Testaments}, \all{des Erbes}.
\end{propriete}

\begin{center}
\begin{tabular}{|l|l|l|l|l|}
\hline
\rowcolor{popblueL} & \textbf{Masculin} & \textbf{Féminin} & \textbf{Neutre} & \textbf{Pluriel} \\
\hline
Nominatif & der Vater & die Mutter & das Kind & die Kinder \\
\hline
Génitif défini & des Vaters & der Mutter & des Kindes & der Kinder \\
\hline
Génitif indéfini & eines Vaters & einer Mutter & eines Kindes & (keiner Kinder) \\
\hline
Traduction & d'un père / du père & d'une mère / de la mère & d'un enfant / de l'enfant & des enfants \\
\hline
\end{tabular}
\end{center}

\begin{exemplebox}[Exemples traduits]
\begin{itemize}
\item \all{das Haus der Familie} — « la maison de la famille »
\item \all{der Sohn des Großvaters} — « le fils du grand-père »
\item \all{die Mutter der Kinder} — « la mère des enfants »
\item \all{der Wille des Verstorbenen} — « la volonté du défunt »
\item \all{die Kinder der Erbin} — « les enfants de l'héritière »
\item \all{das Testament eines reichen Mannes} — « le testament d'un homme riche »
\item \all{die Regelung der Erbschaft} — « le règlement de la succession »
\item \all{die Entscheidung des Richters} — « la décision du juge »
\end{itemize}
\end{exemplebox}

\courssub{Le génitif parmi les quatre cas}

Pour bien situer le génitif dans le système des cas allemands, voici le tableau des quatre cas avec l'article défini :

\begin{popfigure}
\begin{center}
\begin{tikzpicture}[
  case/.style={draw=popblue, fill=popblueL, rounded corners=2pt, minimum width=2.6cm, minimum height=1.5cm, align=center, font=\small},
  gen/.style={draw=poporange, fill=poporangeL, rounded corners=2pt, minimum width=2.6cm, minimum height=1.5cm, align=center, font=\small},
  lbl/.style={font=\small\bfseries, text=popdark}
]
\node[lbl] at (0,2.2) {Masculin};
\node[lbl] at (0,0.6) {Féminin};
\node[lbl] at (0,-1.0) {Neutre};
\node[lbl] at (0,-2.6) {Pluriel};
\node[case, align=center] at (3,2.2){N\\ der Vater};
\node[case, align=center] at (6,2.2){A\\ den Vater};
\node[case, align=center] at (9,2.2){D\\ dem Vater};
\node[gen, align=center] at (12,2.2){G\\ des Vaters};
\node[case, align=center] at (3,0.6){N\\ die Mutter};
\node[case, align=center] at (6,0.6){A\\ die Mutter};
\node[case, align=center] at (9,0.6){D\\ der Mutter};
\node[gen, align=center] at (12,0.6){G\\ der Mutter};
\node[case, align=center] at (3,-1.0){N\\ das Kind};
\node[case, align=center] at (6,-1.0){A\\ das Kind};
\node[case, align=center] at (9,-1.0){D\\ dem Kind};
\node[gen, align=center] at (12,-1.0){G\\ des Kindes};
\node[case, align=center] at (3,-2.6){N\\ die Erben};
\node[case, align=center] at (6,-2.6){A\\ die Erben};
\node[case, align=center] at (9,-2.6){D\\ den Erben};
\node[gen, align=center] at (12,-2.6){G\\ der Erben};
\node[font=\small\itshape, text=popmuted] at (12,-4.0) {Le génitif (colonne orange) : cas de la possession};
\end{tikzpicture}
\end{center}
\legende{Les quatre cas de l'article défini ; le génitif est la seule colonne où le nom lui-même se modifie (-s/-es au masculin et au neutre).}
\end{popfigure}

\courssub{Emploi 1 : la possession}

\begin{propriete}[Le génitif de possession]
Le génitif exprime le rapport de possession ou d'appartenance, exactement comme le français « de » : \all{das Testament des Vaters} = « le testament du père ». En allemand, le possesseur au génitif se place \textbf{après} le nom possédé, sans préposition :
\all{das Erbe der Großmutter} — « l'héritage de la grand-mère » ; \all{das Haus meines Bruders} — « la maison de mon frère ».
\end{propriete}

Ordre soutenu possible : on peut inverser et placer le génitif \textbf{avant} le nom (style écrit, soutenu) : \all{meines Vaters Haus} = « la maison de mon père » ; \all{des Großvaters Testament}. À l'oral courant, on préfère souvent \all{von} + datif : \all{das Haus von meinem Vater}, mais à l'écrit et à l'examen, le génitif est exigé.

\courssub{Emploi 2 : le complément du nom}

\begin{propriete}[Le génitif, complément du nom]
Le génitif précise aussi le contenu, la matière ou l'objet d'un nom abstrait, là où le français utilise « de » : \all{die Regelung der Erbschaft} = « le règlement de la succession » ; \all{die Frage des Erbes} = « la question de l'héritage » ; \all{der Anfang des Streits} = « le début du conflit » ; \all{die Lösung des Problems} = « la solution du problème ».
\end{propriete}

\courssub{Les noms propres et les cas particuliers}

\begin{propriete}[Noms propres au génitif]
Les noms propres de personnes prennent \all{-s} \textbf{sans article} et se placent souvent \textbf{avant} le nom : \all{Peters Erbe} (l'héritage de Pierre), \all{Marias Testament} (le testament de Maria), \all{Kameruns Geschichte} (l'histoire du Cameroun). Si le nom propre se termine déjà par \all{-s}, \all{-ß}, \all{-z} ou \all{-x}, on met une apostrophe : \all{Hans' Haus}, \all{Fritz' Erbe}.
\end{propriete}

\begin{propriete}[Cas particuliers]
\begin{itemize}
\item Noms en \all{-s}, \all{-ß}, \all{-z}, \all{-sch} : désinence \all{-es} : \all{das Haus} $\rightarrow$ \all{des Hauses} ; \all{der Fluss} $\rightarrow$ \all{des Flusses} ; \all{das Gesetz} $\rightarrow$ \all{des Gesetzes}.
\item \textbf{Noms faibles} (masculins en \all{-e} ou désignant des êtres animés) : ils prennent \all{-(e)n} et non \all{-s} : \all{der Erbe} $\rightarrow$ \all{des Erben} ; \all{der Mensch} $\rightarrow$ \all{des Menschen} ; \all{der Student} $\rightarrow$ \all{des Studenten} ; \all{der Junge} $\rightarrow$ \all{des Jungen}.
\item \textbf{Adjectif épithète au génitif} : après article défini, l'adjectif prend \all{-en} ; après article indéfini au masculin/neutre, il prend \all{-en} aussi : \all{der Erbe meines verstorbenen Vaters} (l'héritier de mon père défunt) ; \all{das Testament eines fremden Mannes} (le testament d'un homme étranger).
\end{itemize}
\end{propriete}

\begin{attention}[Noms faibles : attention aux exceptions !]
Les noms faibles prennent \all{-en} : \all{der Erbe} $\rightarrow$ \all{des Erben}, \all{der Mensch} $\rightarrow$ \all{des Menschen}. Mais il existe des exceptions notables : \all{der Notar} $\rightarrow$ \all{des Notars} (avec \all{-s} !), \all{der Name} $\rightarrow$ \all{des Namens}, \all{der Buchstabe} $\rightarrow$ \all{des Buchstabens}. Quand un doute existe, vérifie dans le dictionnaire : le génitif y est toujours indiqué.
\end{attention}

\courssub{Comparaison français / allemand}

\begin{center}
\begin{tabularx}{\linewidth}{|X|X|X|}
\hline
\rowcolor{popblueL} Français & Deutsch & Remarque \\
\hline
la maison de la famille & das Haus der Familie & « de la » = \all{der} (féminin), sans désinence sur le nom \\
\hline
le testament du père & das Testament des Vaters & « du » = \all{des} + \all{-s} sur le nom \\
\hline
les enfants de l'héritière & die Kinder der Erbin & féminin : \all{der}, pas de \all{-s} \\
\hline
la volonté du défunt & der Wille des Verstorbenen & adjectif nominalisé : \all{-en} (nom faible) \\
\hline
la maison de mon père & meines Vaters Haus & ordre inversé possible, style soutenu \\
\hline
l'héritage de Pierre & Peters Erbe & nom propre : \all{-s} sans article, placé avant \\
\hline
\end{tabularx}
\end{center}

\begin{aretenir}[L'essentiel]
Le français marque la possession avec la préposition « de » ; l'allemand la marque avec une \textbf{désinence} : article \all{des/der/des/der} et \all{-(e)s} sur les noms masculins et neutres singuliers. Question d'identification : \all{wessen ?} (de qui ? de quoi ?).
\end{aretenir}

\courssub{Méthode : reconnaître et former le génitif}

\begin{methode}[Comment reconnaître et construire un génitif]
\begin{enumerate}
\item \textbf{Poser la question} \all{wessen ?} (de qui ? de quoi ?) : si le groupe nominal répond à cette question, c'est un génitif. Exemple : \all{das Testament des Vaters} — le testament \emph{de qui} ? du père.
\item \textbf{Identifier le genre et le nombre} du nom possesseur (dictionnaire si besoin) : \all{der Vater} = masculin singulier.
\item \textbf{Choisir l'article} : masculin/neutre $\rightarrow$ \all{des} (défini) ou \all{eines} (indéfini) ; féminin/pluriel $\rightarrow$ \all{der} / \all{einer}.
\item \textbf{Ajouter la désinence} au nom si masculin ou neutre singulier : \all{-es} (une syllabe, consonne finale dure) ou \all{-s} (plusieurs syllabes) ; \all{-(e)n} pour les noms faibles.
\item \textbf{Vérifier l'adjectif} éventuel : \all{-en} après article au génitif : \all{des verstorbenen Vaters}.
\end{enumerate}
\end{methode}

\begin{exoresolu}[Mets au génitif]
Énoncé : complète avec le génitif. 1. Das Haus (\all{die Familie}) ist groß. 2. Das Testament (\all{der Großvater}) liegt beim Notar. 3. Die Kinder (\all{die Erbin}) streiten sich. 4. Der Wille (\all{der Verstorbene}) ist klar. 5. Das Erbe (\all{ein reicher Mann}) ist groß.
\tcblower
Solution détaillée : 1. \all{der Familie} — féminin : article \all{der}, pas de désinence : « la maison de la famille ». 2. \all{des Großvaters} — masculin plurisyllabique : \all{des} + \all{-s} : « le testament du grand-père ». 3. \all{der Erbin} — féminin : \all{der} : « les enfants de l'héritière ». 4. \all{des Verstorbenen} — adjectif nominalisé, décliné comme un nom faible : \all{des} + \all{-en} : « la volonté du défunt ». 5. \all{eines reichen Mannes} — indéfini masculin : \all{eines} + adjectif \all{-en} + nom avec \all{-es} (une syllabe) : « l'héritage d'un homme riche ».
\end{exoresolu}

\begin{attention}[L'erreur n° 1 des francophones]
Ne jamais oublier le \all{-s} du nom masculin ou neutre ! Le français ne marque que l'article (« du père »), donc le réflexe francophone produit : *\all{das Haus der Großvater} ou *\all{das Testament des Vater}. Formes correctes : \all{das Haus des Großvaters}, \all{das Testament des Vaters}. Autre piège : ne pas confondre le génitif \all{der} (féminin/pluriel) avec le nominatif masculin \all{der} : dans \all{die Mutter der Kinder}, \all{der Kinder} est un génitif pluriel, pas un nominatif !
\end{attention}

\begin{exercice}[À ton tour]
1. Traduis en allemand : a) la maison de la famille ; b) le fils du notaire ; c) les enfants de l'héritière ; d) le règlement de la succession ; e) l'héritage de Maria.
2. Mets au génitif : \all{die Entscheidung (der Richter)} ; \all{das Geld (das Kind)} ; \all{der Streit (die Familien)} ; \all{die Kopie (das Testament)}.
\end{exercice}

\begin{corrige}[Exercice « À ton tour »]
1. a) \all{das Haus der Familie} ; b) \all{der Sohn des Notars} (exception : \all{Notar} prend \all{-s}) ; c) \all{die Kinder der Erbin} ; d) \all{die Regelung der Erbschaft} ; e) \all{Marias Erbe}.
2. \all{die Entscheidung des Richters} ; \all{das Geld des Kindes} ; \all{der Streit der Familien} ; \all{die Kopie des Testaments}.
\end{corrige}

\begin{savaistu}[Le génitif disparaît-il ?]
À l'oral, beaucoup d'Allemands remplacent le génitif par \all{von} + datif : \all{das Haus von meinem Vater}. Un livre humoristique à succès, \all{Der Dativ ist dem Genitiv sein Tod} (« le datif est la mort du génitif » — phrase volontairement fautive !), se moque de cette tendance. Mais à l'écrit, dans les testaments, les articles de presse et les examens, le génitif reste la norme : maîtrise-le, c'est la marque d'un allemand soigné.
\end{savaistu}

Maintenant que tu sais exprimer « de qui » et « de quoi » on parle grâce au génitif, passons à la section suivante : le parfait et le prétérit, les deux temps qui permettent de \emph{raconter} l'histoire de cette succession — qui a dit quoi, qui a trouvé le testament, et comment le conflit s'est terminé.

\coursec{Grammaire 2 : parfait et prétérit pour raconter}

Dans la section précédente, nous avons étudié le génitif, ce cas qui permet d'exprimer la possession : \all{das Testament des Großvaters}, « le testament du grand-père ». Mais un récit d'héritage, ce n'est pas seulement des possessions : c'est avant tout une \emph{histoire}. Le grand-père est mort, le notaire a lu le testament, les héritiers se sont disputés. Pour raconter ces événements passés, l'allemand dispose de deux temps principaux : le \cle{Perfekt} et le \cle{Präteritum}. C'est ce que nous allons découvrir maintenant, en partant, comme toujours, de l'observation du texte.

\courssub{Observation : deux temps pour un même passé}

Reprenons deux phrases tirées de notre texte support \emph{Der Streit um das Erbe} :

\begin{exemplebox}[Deux temps dans le texte]
\all{Als der Großvater starb, hinterließ er ein Haus und ein großes Grundstück.}\\
« Quand le grand-père mourut, il laissa une maison et un grand terrain. »

\medskip
\all{Ein Notar hat die Erbschaft geregelt und das Testament vorgelesen.}\\
« Un notaire a réglé la succession et a lu le testament. »
\end{exemplebox}

Observons attentivement :

\begin{itemize}
\item Dans la première phrase, le verbe est en \textbf{un seul mot} : \all{starb}, \all{hinterließ}. C'est le \cle{Präteritum} (prétérit).
\item Dans la deuxième phrase, le verbe est en \textbf{deux mots} : un auxiliaire conjugué (\all{hat}) en deuxième position et un participe (\all{geregelt}, \all{vorgelesen}) renvoyé \textbf{à la fin de la phrase}. C'est le \cle{Perfekt} (parfait).
\end{itemize}

Les deux temps expriment le passé. La différence n'est pas d'abord une différence de sens, mais une différence de \textbf{registre} : le prétérit est le temps du récit écrit (presse, littérature, textes narratifs), le parfait est le temps de la conversation et de la correspondance. Nous allons maintenant étudier chacun en détail.

\courssub{Le parfait (Perfekt) : formation}

\begin{propriete}[Formation du parfait]
Le parfait se forme avec l'auxiliaire \cle{haben} (ou \cle{sein}) conjugué au présent \textbf{en deuxième position}, suivi du \cle{Partizip II} (participe passé) placé \textbf{à la fin de la phrase} :

\begin{center}
Sujet + \textbf{haben/sein} (position 2) + \ldots\ + \textbf{Partizip II} (fin)
\end{center}

Formation du participe II :
\begin{itemize}
\item \textbf{Verbes faibles} : \all{ge-} + radical + \all{-t} : \all{machen} $\rightarrow$ \all{gemacht} ; \all{regeln} $\rightarrow$ \all{geregelt}.
\item \textbf{Verbes forts} : \all{ge-} + radical (souvent modifié) + \all{-en} : \all{gehen} $\rightarrow$ \all{gegangen} ; \all{sterben} $\rightarrow$ \all{gestorben}.
\item \textbf{Pas de ge-} pour les verbes en \all{-ieren} : \all{studieren} $\rightarrow$ \all{studiert} ; et pour les verbes à préfixe inséparable (ver-, be-, er-, ent-, ge-, zer-, miss-) : \all{verkaufen} $\rightarrow$ \all{verkauft} ; \all{bekommen} $\rightarrow$ \all{bekommen}.
\item \textbf{Verbes à particule séparable} : le \all{ge-} s'insère entre la particule et le verbe : \all{vorlesen} $\rightarrow$ \all{vorgelesen} ; \all{mitbekommen} $\rightarrow$ \all{mitbekommen}.
\end{itemize}
\end{propriete}

\begin{popfigure}
\begin{tikzpicture}[
  box/.style={draw=popblue, fill=popblueL, rounded corners=2pt, align=center, font=\small, inner sep=5pt},
  posS/.style={font=\footnotesize\itshape, text=popmuted}
]
\node[box, align=center] (s) at (0,0){Der Notar\\ \footnotesize Sujet};
\node[box, fill=poporangeL, draw=poporange, align=center] (aux) at (3.2,0){\textbf{hat}\\ \footnotesize Auxiliaire (pos. 2)};
\node[box, align=center] (m) at (6.4,0){die Erbschaft\\ gestern\\ \footnotesize compléments};
\node[box, fill=popgreenL, draw=popgreen, align=center] (p) at (9.8,0){\textbf{geregelt}.\\ \footnotesize Partizip II (fin)};
\draw[-{Stealth}, thick, poporange] (aux.south) to[bend right=25] node[posS, below]{cadre verbal} (p.south);
\end{tikzpicture}
\legende{Le cadre verbal au parfait : l'auxiliaire en position 2, le participe II à la toute fin.}
\end{popfigure}

\begin{attention}[Le participe à la fin, toujours !]
Les francophones oublient souvent de renvoyer le participe en fin de phrase, surtout quand la phrase est longue. Jamais : *\all{Der Notar hat geregelt die Erbschaft}. Correct : \all{Der Notar hat die Erbschaft geregelt.} Tout ce qui se trouve entre l'auxiliaire et le participe forme le « cadre verbal » (\all{Satzklammer}).
\end{attention}

\courssub{Le choix de l'auxiliaire : haben ou sein ?}

C'est LA grande question du parfait allemand. La règle est simple à énoncer, mais il faut la maîtriser parfaitement.

\begin{propriete}[Choix de l'auxiliaire]
\begin{itemize}
\item \textbf{haben} : c'est l'auxiliaire par défaut. Il s'emploie avec tous les verbes \textbf{transitifs} (qui ont un complément d'objet direct), les verbes réfléchis, les modaux et la grande majorité des verbes : \all{Er hat das Haus verkauft.} « Il a vendu la maison. »
\item \textbf{sein} : il s'emploie avec les verbes \textbf{intransitifs de déplacement} (gehen, kommen, fahren, laufen, reisen, fliegen...) et les verbes de \textbf{changement d'état} (sterben, werden, aufwachen, einschlafen, passieren, gelingen, mißlingen...) : \all{Der Großvater ist gestorben.} « Le grand-père est mort. »
\end{itemize}
\end{propriete}

\begin{exemplebox}[Exemples traduits]
\all{Die Erben sind nach Yaoundé gefahren.} « Les héritiers sont allés à Yaoundé. » (déplacement $\rightarrow$ sein)

\all{Ein Notar hat das Problem gelöst.} « Un notaire a résolu le problème. » (transitif $\rightarrow$ haben)

\all{Die Familie ist nach dem Tod des Vaters auseinandergegangen.} « La famille s'est dispersée après la mort du père. » (déplacement $\rightarrow$ sein)

\all{Wir haben lange über das Testament gesprochen.} « Nous avons longuement parlé du testament. » (intransitif sans déplacement $\rightarrow$ haben)
\end{exemplebox}

\begin{attention}[Jamais sein avec un verbe transitif !]
Erreur classique du francophone, influencé par l'accord du participe en français : *\all{Er ist das Haus verkauft.} est FAUX. Le verbe \all{verkaufen} a un COD (\all{das Haus}), donc : \all{Er \textbf{hat} das Haus verkauft.} « Il a vendu la maison. »
\end{attention}

\begin{attention}[sterben et werden se conjuguent avec sein]
\all{sterben} (mourir) exprime un changement d'état : \all{Der Großvater \textbf{ist} gestorben.} et jamais *\all{hat gestorben}. Même logique pour \all{werden} (devenir) : \all{Sie ist Anwältin geworden.} « Elle est devenue avocate. »
\end{attention}

\courssub{Le prétérit des verbes faibles}

Le prétérit (\cle{Präteritum}) est un temps \textbf{simple} : le verbe se conjugue en un seul mot, comme notre imparfait ou notre passé simple.

\begin{propriete}[Formation du prétérit des verbes faibles]
Radical du verbe + \textbf{-te} + désinence personnelle. Exemples : \all{machen} $\rightarrow$ \all{machte} ; \all{regeln} $\rightarrow$ \all{regelte}.

\begin{center}
\begin{tabular}{|l|l|l|}
\hline
\rowcolor{popblueL} Personne & machen (faire) & regeln (régler) \\
\hline
ich & machte & regelte \\
\hline
du & machtest & regeltest \\
\hline
er/sie/es & machte & regelte \\
\hline
wir & machten & regelten \\
\hline
ihr & machtet & regeltet \\
\hline
sie/Sie & machten & regelten \\
\hline
\end{tabular}
\end{center}

Remarque : la 1re et la 3e personne du singulier sont \textbf{identiques} (\all{ich machte / er machte}), comme en français « je faisais / il faisait ».
\end{propriete}

\textbf{Cas particuliers à connaître :}

\begin{itemize}
\item \textbf{Verbes en -ern / -eln} : on garde le radical tel quel : \all{regeln} $\rightarrow$ \all{regelte} ; \all{ändern} (changer) $\rightarrow$ \all{änderte}.
\item \textbf{Radical terminé par -d- ou -t-} (ou -m/-n après consonne) : on insère un \textbf{-e-} de liaison pour la prononciation : \all{arbeiten} $\rightarrow$ \all{arbeitete} ; \all{warten} $\rightarrow$ \all{wartete} ; \all{öffnen} $\rightarrow$ \all{öffnete}. Comparez : \all{machte} se prononce « ma-rte », mais *\all{arbeitte} serait imprononçable, d'où \all{arbeitete} « ar-baï-tè-te ».
\end{itemize}

\begin{exemplebox}[Exemples traduits au prétérit]
\all{Die Erben stritten monatelang.} « Les héritiers se disputèrent pendant des mois. » (verbe fort)

\all{Der Notar öffnete das Testament.} « Le notaire ouvrit le testament. »

\all{Die Familie wartete lange auf eine Lösung.} « La famille attendit longtemps une solution. »

\all{Das Haus wurde verkauft.} « La maison fut vendue. » (passif au prétérit : \all{wurde} + participe)
\end{exemplebox}

\courssub{Emplois : prétérit ou parfait ?}

\begin{propriete}[Répartition des deux temps]
\begin{itemize}
\item \textbf{Präteritum} : récit \textbf{écrit}, presse, littérature, contes, comptes rendus, exposés formels. C'est le temps du narrateur.
\item \textbf{Perfekt} : \textbf{oral}, conversation, dialogues, lettres et courriels personnels. C'est le temps du locuteur qui raconte.
\item \textbf{Exception majeure} : même à l'oral, on utilise le \textbf{prétérit} pour \all{sein} (\all{war}), \all{haben} (\all{hatte}) et les \textbf{verbes modaux} (\all{wollte, musste, konnte, durfte, sollte}). Personne ne dit *\all{ich bin gewesen} dans une conversation normale : on dit \all{ich war}.
\end{itemize}
\end{propriete}

\begin{exemplebox}[Les verbes fréquents au prétérit, même à l'oral]
\all{Der Großvater war sehr streng.} « Le grand-père était très sévère. »

\all{Die Familie hatte früher ein großes Grundstück.} « La famille avait autrefois un grand terrain. »

\all{Die Brüder wollten das Haus nicht teilen.} « Les frères ne voulaient pas partager la maison. »

\all{Sie konnten sich nicht einigen.} « Ils ne pouvaient pas se mettre d'accord. »
\end{exemplebox}

\begin{methode}[Comment choisir entre prétérit et parfait]
\begin{enumerate}
\item \textbf{Identifier le contexte} : s'agit-il d'un texte écrit narratif (article, récit, roman) ou d'un dialogue / d'une lettre personnelle ?
\item \textbf{Texte écrit narratif} $\rightarrow$ prétérit pour les verbes d'action : \all{Er hinterließ ein Haus.}
\item \textbf{Dialogue, oral, lettre} $\rightarrow$ parfait : \all{Ein Notar hat alles geregelt.}
\item \textbf{Vérifier le verbe} : s'il s'agit de \all{sein}, \all{haben} ou d'un modal, utiliser le \textbf{prétérit même à l'oral} : \all{Ich war gestern beim Notar.} « J'étais hier chez le notaire. »
\item \textbf{Au parfait, choisir l'auxiliaire} : verbe de déplacement ou de changement d'état $\rightarrow$ sein ; sinon $\rightarrow$ haben.
\end{enumerate}
\end{methode}

\begin{popfigure}
\begin{tikzpicture}[font=\small]
\draw[-{Stealth}, very thick, popdark] (-0.3,0) -- (12.5,0) node[below left=2pt and 0pt, font=\footnotesize\itshape]{Vergangenheit (le passé)};
\foreach \x in {0.5,2,3.5,5,6.5,8,9.5,11,12} \draw[popline] (\x,-0.08) -- (\x,0.08);
\node[draw=popblue, fill=popblueL, rounded corners=3pt, align=center, inner sep=6pt] at (3.2,1.5) {\textbf{Präteritum}\\ \footnotesize récit écrit, presse, littérature\\ \footnotesize\all{Er hinterließ ein Haus.}};
\node[draw=popgreen, fill=popgreenL, rounded corners=3pt, align=center, inner sep=6pt] at (9,-1.7) {\textbf{Perfekt}\\ \footnotesize oral, conversation, lettres\\ \footnotesize\all{Ein Notar hat alles geregelt.}};
\node[draw=poporange, fill=poporangeL, rounded corners=3pt, align=center, inner sep=5pt, font=\footnotesize] at (3.2,-1.6) {même à l'oral :\\ \all{war, hatte, wollte, musste, konnte}};
\draw[-{Stealth}, thick, popblue] (3.2,1.1) -- (3.2,0.12);
\draw[-{Stealth}, thick, popgreen] (9,-1.3) -- (9,-0.12);
\end{tikzpicture}
\legende{Frise des temps du passé : le prétérit domine l'écrit narratif, le parfait domine l'oral.}
\end{popfigure}

\courssub{Comparaison avec le français}

Le français possède plusieurs temps du passé (passé composé, imparfait, passé simple) avec des nuances d'aspect. L'allemand, lui, n'oppose pas le prétérit et le parfait par le sens : les deux traduisent une action passée accomplie. La nuance entre « il mourut » (passé simple) et « il mourait » (imparfait) n'existe pas en allemand : on dit dans les deux cas \all{er starb}, et le contexte (ou un adverbe comme \all{oft}, \all{plötzlich}) précise le sens.

\begin{center}
\begin{tabularx}{\linewidth}{|X|X|X|}
\hline
\rowcolor{popblueL} Français & Deutsch & Remarque \\
\hline
Il a vendu la maison. (passé composé) & Er hat das Haus verkauft. & parfait $\approx$ passé composé (oral) \\
\hline
Il vendit la maison. (passé simple) & Er verkaufte das Haus. & prétérit, récit écrit \\
\hline
Il vendait la maison. (imparfait) & Er verkaufte das Haus. & même forme allemande ! \\
\hline
Le grand-père est mort. & Der Großvater ist gestorben. & sein + participe, comme en français \\
\hline
Il était très malade. & Er war sehr krank. & sein au prétérit, même à l'oral \\
\hline
\end{tabularx}
\end{center}

\begin{aretenir}[L'essentiel]
\begin{itemize}
\item Parfait = haben/sein (position 2) + Partizip II (fin de phrase).
\item sein avec les déplacements et changements d'état : \all{Er ist gestorben.}
\item Prétérit des verbes faibles = radical + -te : \all{machte, regelte, arbeitete}.
\item Écrit narratif $\rightarrow$ prétérit ; oral et lettres $\rightarrow$ parfait.
\item \all{sein}, \all{haben} et les modaux : prétérit partout (\all{war, hatte, konnte}).
\item L'allemand ne connaît pas la nuance imparfait / passé simple.
\end{itemize}
\end{aretenir}

\courssub{Texte d'application : une succession à Douala}

\begin{textebox}[Nach dem Tod des Vaters]
Als Herr Mbappe starb, hinterließ er seiner Familie ein Haus in Bonabéri und ein Stück Land in der Nähe von Douala. Der Vater hatte kein Testament gemacht, und so begann der Streit. Der älteste Sohn wollte das Haus behalten, aber seine Schwester wollte es verkaufen und das Geld teilen. Monatelang stritten die Geschwister, und die Mutter weinte oft. Sie war sehr traurig, denn die Familie war früher immer einig gewesen.

Eines Tages ging die Mutter zu einem Notar. Der Notar öffnete eine Mappe, las das Gesetz und erklärte den Erben ihre Rechte. Er arbeitete lange an diesem Fall und fand endlich eine Lösung: das Haus wurde verkauft, und jeder Erbe bekam seinen Teil. „Ich habe viele solcher Fälle gesehen“, sagte der Notar, „aber diese Familie hat endlich Frieden gefunden.“ Die Geschwister haben sich versöhnt, und heute besuchen sie ihre Mutter wieder jeden Sonntag.

\medskip
\emph{Glossaire} : \all{hinterlassen} (hinterließ, hat hinterlassen) : laisser (en héritage) ; \all{das Stück Land} : le terrain ; \all{behalten} (behielt, hat behalten) : garder ; \all{teilen} : partager ; \all{die Geschwister} (Pl.) : les frères et sœurs ; \all{einig} : uni, d'accord ; \all{die Mappe, -n} : la chemise, le dossier ; \all{das Recht, -e} : le droit ; \all{der Fall, „e} : le cas, l'affaire ; \all{sich versöhnen} : se réconcilier ; \all{hatte ... gemacht / war ... gewesen} : formes du plus-que-parfait (antériorité dans le passé).
\end{textebox}

\begin{savaistu}[Le plus-que-parfait (Plusquamperfekt)]
Dans le texte, vous avez relevé \all{hatte ... gemacht} et \all{war ... gewesen}. Ce sont des formes du \textbf{plus-que-parfait} (Plusquamperfekt), le temps de l'antériorité dans le passé (ce qui s'est passé \emph{avant} un autre événement passé). Il se forme exactement comme le parfait, mais avec l'auxiliaire \textbf{au prétérit} : \all{hatte gemacht} (au lieu de \all{hat gemacht}), \all{war gewesen} (au lieu de \all{ist gewesen}). Il est fréquent dans les récits écrits pour situer des antécédents.
\end{savaistu}

\textbf{Questions de compréhension :}
\begin{enumerate}
\item Was hinterließ Herr Mbappe seiner Familie ?
\item Warum begann der Streit zwischen den Geschwistern ?
\item Wie löste der Notar das Problem ?
\item Relevez dans le texte trois verbes au prétérit et deux verbes au parfait, puis justifiez l'emploi de chaque temps.
\end{enumerate}

\courssub{Exercice résolu}

\begin{exoresolu}[Prétérit ou parfait ?]
Mettez les verbes entre parenthèses à la forme correcte (prétérit ou parfait) selon le contexte.

a) (Article de journal) Der Streit um das Erbe der Familie Eto'o (beginnen) \ldots\ vor zwei Jahren.

b) (Dialogue entre amis) „(Sehen) \ldots\ du den Notar gestern ?“ — „Ja, er (haben) \ldots\ keine Zeit.“

c) (Récit) Der Großvater (sterben) \ldots\ im Jahr 2015 und (hinterlassen) \ldots\ ein großes Haus.

d) (Lettre à un ami) Lieber Paul, mein Onkel (verkaufen) \ldots\ endlich sein Haus in Douala!

\tcblower
\textbf{Solution détaillée :}

a) Contexte : article de journal, donc récit écrit $\rightarrow$ prétérit. \all{beginnen} est un verbe fort : \textbf{begann}. \all{Der Streit um das Erbe der Familie Eto'o \textbf{begann} vor zwei Jahren.}

b) Contexte : dialogue oral $\rightarrow$ parfait. \all{sehen} : verbe fort, auxiliaire haben : \textbf{hast ... gesehen}. Pour \all{haben}, on préfère le prétérit même à l'oral : \textbf{hatte}. \all{„\textbf{Hast} du den Notar gestern \textbf{gesehen}?“ — „Ja, er \textbf{hatte} keine Zeit.“}

c) Contexte : récit écrit $\rightarrow$ prétérit. \all{sterben} : verbe fort, \textbf{starb} ; \all{hinterlassen} : verbe fort, \textbf{hinterließ}. \all{Der Großvater \textbf{starb} im Jahr 2015 und \textbf{hinterließ} ein großes Haus.}

d) Contexte : lettre personnelle $\rightarrow$ parfait. \all{verkaufen} : verbe faible à préfixe inséparable, auxiliaire haben (transitif) : \textbf{hat ... verkauft}. \all{Mein Onkel \textbf{hat} endlich sein Haus in Douala \textbf{verkauft}!}
\end{exoresolu}

\begin{exercice}[À vous de jouer]
Conjuguez au prétérit : 1. der Notar (regeln) ; 2. die Erben (warten) ; 3. wir (machen) ; 4. ihr (öffnen) ; 5. ich (arbeiten). Puis mettez les phrases suivantes au parfait en choisissant le bon auxiliaire : 6. Die Familie (kommen) zum Notar. 7. Der Sohn (verkaufen) das Grundstück. 8. Die Mutter (werden) sehr traurig.
\end{exercice}

\begin{corrige}[Exercice : prétérit et parfait]
1. regelte ; 2. warteten (e de liaison après -t-) ; 3. machten ; 4. öffnetet ; 5. arbeitete (e de liaison après -t-). 6. Die Familie \textbf{ist} zum Notar \textbf{gekommen} (déplacement $\rightarrow$ sein). 7. Der Sohn \textbf{hat} das Grundstück \textbf{verkauft} (transitif $\rightarrow$ haben ; pas de ge- devant ver-). 8. Die Mutter \textbf{ist} sehr traurig \textbf{geworden} (changement d'état $\rightarrow$ sein).
\end{corrige}

\begin{attention}[Récapitulatif des pièges pour les francophones]
\begin{itemize}
\item Ne pas traduire le passé composé français mot à mot : « il a vendu » peut être \all{er hat verkauft} (oral) ou \all{er verkaufte} (récit écrit).
\item Ne pas calquer « il est allé » sur *\all{er ist gegangen}... en fait ici le français et l'allemand coïncident : \all{er ist gegangen} est correct ! Mais attention : « il a vendu » n'est JAMAIS *\all{er ist verkauft}.
\item Ne pas oublier le -e- de liaison : \all{arbeitete}, \all{wartete}, et non *\all{arbeitte}, *\all{wartte}.
\item Ne pas dire à l'oral *\all{ich bin gewesen} ou *\all{ich habe gehabt} : on dit \all{ich war}, \all{ich hatte}.
\end{itemize}
\end{attention}

Nous savons désormais raconter les événements du passé, à l'écrit comme à l'oral. Dans la dernière section, nous mettrons toutes ces compétences en pratique : commentaire du texte et production guidée sur le thème de l'héritage.

\coursec{Commentaire et production guidée}

Dans la section précédente, vous avez appris à raconter au passé avec le \cle{Perfekt} et le \cle{Präteritum}. Ces deux temps vont maintenant vous servir à rédiger : un résumé fidèle du texte \all{Der Streit um das Erbe}, un commentaire structuré et une courte production personnelle. C'est exactement ce que l'on vous demandera au baccalauréat : comprendre, résumer, puis donner votre avis en allemand correct.

\courssub{La méthode du commentaire de texte en 4 étapes}

Un commentaire de texte (\all{die Textanalyse} ou \all{der Kommentar}) n'est pas une improvisation : il suit un plan fixe, en quatre parties, que les correcteurs attendent.

\begin{methode}[Rédiger un commentaire en 4 étapes]
\begin{enumerate}
\item \textbf{\all{Einleitung} (introduction)} : présenter le texte en une ou deux phrases — de quoi parle-t-il ? quel est le thème ? Formule type : \all{Der Text handelt von...} / \all{In dem Text geht es um...} (Le texte parle de... / Il s'agit dans le texte de...).
\item \textbf{\all{Inhalt} (contenu)} : résumer le texte en 4 à 5 phrases, au passé, avec ses propres mots, sans opinion.
\item \textbf{\all{Analyse} (analyse)} : décrire les personnages (\all{die Personen}), le conflit (\all{der Konflikt}) et les moyens linguistiques (\all{die sprachlichen Mittel}) : temps employés, mots du champ lexical de l'héritage, ton du texte.
\item \textbf{\all{Eigene Meinung} (opinion personnelle)} : dire ce que l'on pense du conflit et de sa solution, avec \all{meiner Meinung nach}, \all{ich finde}, \all{ich denke, dass...}.
\end{enumerate}
\end{methode}

\begin{popfigure}
\begin{tikzpicture}[node distance=0.55cm,
  etape/.style={rectangle, rounded corners=3pt, draw=popblue, fill=popblueL, text width=4.6cm, align=center, minimum height=1cm, font=\small},
  detail/.style={rectangle, draw=poppurple, fill=poppurpleL, text width=5.6cm, align=center, minimum height=1cm, font=\small},
  fleche/.style={-{Stealth[length=2.5mm]}, thick, popdark}]
\node[etape, align=center] (e1){\textbf{1. Einleitung}\\Thème du texte};
\node[detail, right=0.9cm of e1] (d1) {\all{Der Text handelt von einem Erbstreit...}};
\node[etape, below=of e1, align=center] (e2){\textbf{2. Inhalt}\\Résumé au passé};
\node[detail, right=0.9cm of e2] (d2) {\all{Nach dem Tod des Großvaters stritten die Erben...}};
\node[etape, below=of e2, align=center] (e3){\textbf{3. Analyse}\\Personnages, conflit, langue};
\node[detail, right=0.9cm of e3] (d3) {\all{Der Konflikt entsteht, weil das Testament nicht klar war...}};
\node[etape, below=of e3, align=center] (e4){\textbf{4. Meinung}\\Opinion personnelle};
\node[detail, right=0.9cm of e4] (d4) {\all{Meiner Meinung nach war die Lösung gerecht...}};
\draw[fleche] (e1) -- (e2);
\draw[fleche] (e2) -- (e3);
\draw[fleche] (e3) -- (e4);
\end{tikzpicture}
\legende{Organigramme du commentaire de texte : quatre étapes obligatoires, de l'introduction à l'opinion.}
\end{popfigure}

\courssub{Les connecteurs utiles}

Pour enchaîner les idées, le correcteur attend des \cle{connecteurs} (\all{die Konnektoren}). Voici les indispensables de cette leçon :

\begin{center}
\begin{tabularx}{\linewidth}{|l|X|X|}
\hline
\rowcolor{popblueL} \textbf{Connecteur} & \textbf{Français} & \textbf{Exemple} \\
\hline
\all{zuerst} & d'abord & \all{Zuerst las der Notar das Testament vor.} (D'abord, le notaire a lu le testament.) \\
\hline
\all{dann} & puis, ensuite & \all{Dann begann der Streit.} (Puis la dispute a commencé.) \\
\hline
\all{danach} & après cela & \all{Danach suchten die Erben eine Lösung.} (Après cela, les héritiers ont cherché une solution.) \\
\hline
\all{schließlich} & finalement & \all{Schließlich einigten sie sich.} (Finalement, ils se sont mis d'accord.) \\
\hline
\all{einerseits... andererseits} & d'une part... d'autre part & \all{Einerseits wollte Anna das Haus, andererseits wollte Paul es verkaufen.} (D'une part Anna voulait la maison, d'autre part Paul voulait la vendre.) \\
\hline
\all{meiner Meinung nach} & à mon avis & \all{Meiner Meinung nach war die Lösung gerecht.} (À mon avis, la solution était juste.) \\
\hline
\all{deshalb} & c'est pourquoi & \all{Das Testament war nicht klar, deshalb gab es einen Streit.} (Le testament n'était pas clair, c'est pourquoi il y a eu une dispute.) \\
\hline
\end{tabularx}
\end{center}

\begin{attention}[Place du verbe après les connecteurs]
Après \all{zuerst}, \all{dann}, \all{danach}, \all{schließlich}, \all{deshalb} et \all{meiner Meinung nach} placés en tête de phrase, le verbe conjugué reste en \textbf{deuxième position} : le sujet passe après le verbe (inversion). On dit \all{Deshalb gab es einen Streit} et jamais \all{Deshalb es gab...}. Attention aussi : \all{meiner Meinung nach} est un génitif figé — on ne dit pas \all{nach meiner Meinung} dans ce sens (c'est possible mais plus lourd et moins idiomatique).
\end{attention}

\courssub{Le résumé guidé (die Zusammenfassung)}

Le résumé est un exercice codifié : 5 phrases maximum, au \all{Perfekt} ou au \all{Präteritum}, sans citation longue, sans opinion.

\begin{methode}[Rédiger un résumé en 4 questions]
Répondez par une phrase au passé à chaque question ; ajoutez éventuellement une cinquième phrase de conclusion.
\begin{enumerate}
\item \textbf{Qui ?} (\all{Wer?}) : de qui parle le texte ? — \all{Der Text handelt von einer deutschen Familie.}
\item \textbf{Quoi ?} (\all{Was?}) : que s'est-il passé ? — \all{Nach dem Tod des Großvaters stritten die Erben um das Haus.}
\item \textbf{Problème ?} (\all{Was war das Problem?}) : — \all{Das Testament war nicht klar, deshalb gab es einen Konflikt.}
\item \textbf{Solution ?} (\all{Wie endete der Streit?}) : — \all{Der Notar fand eine Lösung, und alle Erben bekamen Geld.}
\end{enumerate}
\end{methode}

\begin{attention}[Pas d'opinion dans le résumé !]
Dans la \all{Zusammenfassung}, vous ne donnez \textbf{jamais} votre avis : ni \all{ich finde}, ni \all{meiner Meinung nach}. L'opinion personnelle apparaît uniquement dans la dernière partie du commentaire (\all{eigene Meinung}). Mélanger résumé et opinion est l'erreur la plus sanctionnée au baccalauréat.
\end{attention}

\begin{savaistu}[Le résumé au baccalauréat]
Au Bac, la \all{Zusammenfassung} est notée sur deux critères : la \textbf{fidélité au texte} (avez-vous repris les informations essentielles, sans rien inventer ?) et la \textbf{correction linguistique} (temps du passé, place du verbe, génitifs). L'originalité n'est pas demandée : un résumé simple, fidèle et correct vaut mieux qu'un texte brillant mais hors sujet.
\end{savaistu}

\courssub{Commentaire modèle}

\begin{textebox}[Début d'un commentaire modèle sur \all{Der Streit um das Erbe}]
Der Text handelt von einem Erbstreit in einer deutschen Familie. Nach dem Tod des Großvaters \textbf{stritten} die Erben um das Haus, weil das Testament nicht klar war. Zuerst wollte die Enkelin Anna das Haus behalten, dann verlangte ihr Bruder Paul den Verkauf. Danach schlug der Notar eine Lösung vor: Das Haus wurde verkauft, und das Geld wurde unter den Erben geteilt. Schließlich akzeptierten alle diesen Vorschlag. Einerseits ist der Konflikt typisch für viele Familien, andererseits zeigt der Text, wie wichtig ein klares Testament ist. Meiner Meinung nach war die Lösung des Notars gerecht, weil alle Erben Geld bekommen haben.
\end{textebox}

\textbf{Glossaire :} \all{der Erbstreit, die Erbstreite} (le conflit d'héritage) ; \all{behalten} (garder) ; \all{der Verkauf, die Verkäufe} (la vente) ; \all{vorschlagen (schlug vor, hat vorgeschlagen)} (proposer) ; \all{teilen} (partager) ; \all{der Vorschlag, die Vorschläge} (la proposition) ; \all{gerecht} (juste, équitable).

Dans ce modèle, repérez les connecteurs : \all{zuerst} ouvre le récit, \all{dann} et \all{danach} enchaînent les événements, \all{schließlich} conclut le résumé, \all{einerseits... andererseits} introduit l'analyse, et \all{meiner Meinung nach} annonce l'opinion finale. Notez aussi les génitifs : \all{die Lösung \textbf{des Notars}}, \all{nach dem Tod \textbf{des Großvaters}}, \all{einen Teil \textbf{des Geldes}}.

\begin{exemplebox}[Phrases d'opinion traduites]
\all{Meiner Meinung nach war die Lösung des Notars gerecht, weil alle Erben Geld bekommen haben.} « À mon avis, la solution du notaire était juste parce que tous les héritiers ont reçu de l'argent. »

\all{Ich finde es traurig, dass die Familie wegen des Geldes gestritten hat.} « Je trouve triste que la famille se soit disputée à cause de l'argent. »

\all{Ich denke, dass ein klares Testament viele Probleme vermeiden kann.} « Je pense qu'un testament clair peut éviter beaucoup de problèmes. »
\end{exemplebox}

\begin{exoresolu}[Rédiger un résumé en 5 phrases]
Énoncé : Résumez le texte \all{Der Streit um das Erbe} en 5 phrases maximum, au passé, sans opinion. Utilisez au moins deux connecteurs et un génitif.
\tcblower
Solution détaillée : on suit la méthode des 4 questions (qui ? quoi ? problème ? solution ?) plus une phrase de conclusion.

\all{1. Der Text handelt von einem Erbstreit in einer Familie. 2. Nach dem Tod des Großvaters stritten die Erben um das Haus. 3. Das Testament war nicht klar, deshalb begann ein langer Konflikt. 4. Zuerst wollte Anna das Haus behalten, dann verlangte Paul den Verkauf. 5. Schließlich fand der Notar eine Lösung, und alle Erben bekamen einen Teil des Geldes.}

Vérification : 5 phrases ; passé (\all{stritten}, \all{war}, \all{begann}, \all{wollte}, \all{verlangte}, \all{fand}, \all{bekamen}) ; connecteurs (\all{deshalb}, \all{zuerst}, \all{dann}, \all{schließlich}) ; génitifs (\all{des Großvaters}, \all{des Geldes}) ; aucune opinion. Le résumé est conforme.
\end{exoresolu}

\courssub{Production écrite guidée}

\begin{experience}[Erzählen Sie einen Erbstreit]
Consigne : \all{Erzählen Sie einen Erbstreit (echt oder erfunden) und geben Sie Ihre Meinung.} Racontez un conflit d'héritage (réel ou imaginé — par exemple dans une famille de Yaoundé ou de Douala) et donnez votre avis, en 8 à 10 phrases.

Plan suggéré : 1) présentation de la famille ; 2) le décès et l'héritage (maison, terrain, argent) ; 3) le conflit entre les héritiers ; 4) la solution (ou l'absence de solution) ; 5) votre opinion avec \all{meiner Meinung nach}.

Début possible : \all{In meiner Nachbarschaft gab es letztes Jahr einen Erbstreit. Nach dem Tod des Vaters stritten seine drei Kinder um das Haus und das Land...} (Dans mon quartier, il y a eu l'année dernière un conflit d'héritage. Après la mort du père, ses trois enfants se sont disputés la maison et le terrain...)
\end{experience}

\textbf{Grille d'auto-évaluation.} Avant de rendre votre production, vérifiez chaque point :

\begin{center}
\begin{tabularx}{\linewidth}{|X|l|}
\hline
\rowcolor{popgreenL} \textbf{Critère} & \textbf{Oui / Non} \\
\hline
J'ai employé au moins 3 mots du vocabulaire du thème (\all{das Erbe, der Erbe, das Testament, der Streit, das Gesetz}...) & \\
\hline
J'ai écrit au moins 2 génitifs corrects (\all{das Haus des Vaters, die Lösung des Notars}) & \\
\hline
J'ai conjugué au moins 3 verbes au passé (Perfekt ou Präteritum) & \\
\hline
J'ai utilisé au moins 1 connecteur d'opinion (\all{meiner Meinung nach, ich finde, ich denke}) & \\
\hline
Mon texte compte 8 à 10 phrases, avec le verbe en deuxième position & \\
\hline
\end{tabularx}
\end{center}

\begin{exercice}[À vous]
1. Remettez dans l'ordre les étapes du commentaire : \all{Meinung – Einleitung – Analyse – Inhalt}.

2. Complétez avec \all{zuerst}, \all{danach}, \all{schließlich} ou \all{deshalb} : \all{Das Testament war unklar, ... gab es einen Streit. ... stritten die Brüder, ... riefen sie den Notar, ... einigten sie sich.}

3. Rédigez votre production de 8 à 10 phrases selon la consigne de l'activité ci-dessus, puis auto-évaluez-vous avec la grille.
\end{exercice}

\begin{corrige}[Exercice : Commentaire et production guidée]
1. Ordre correct : \all{Einleitung → Inhalt → Analyse → Meinung} (introduction, résumé, analyse, opinion).

2. \all{Das Testament war unklar, \textbf{deshalb} gab es einen Streit. \textbf{Zuerst} stritten die Brüder, \textbf{danach} riefen sie den Notar, \textbf{schließlich} einigten sie sich.} (Le testament était peu clair, c'est pourquoi il y a eu une dispute. D'abord les frères se sont disputés, après cela ils ont appelé le notaire, finalement ils se sont mis d'accord.)

3. Production libre : le correcteur vérifiera la présence des 3 mots du thème, des 2 génitifs, des 3 verbes au passé et du connecteur d'opinion, ainsi que la place du verbe en deuxième position dans chaque phrase.
\end{corrige}

\begin{aretenir}[L'essentiel de la section]
Un commentaire de texte suit toujours quatre étapes : \all{Einleitung}, \all{Inhalt}, \all{Analyse}, \all{eigene Meinung}. Le résumé se rédige au passé, en 5 phrases maximum, fidèle au texte et sans opinion ; l'opinion n'apparaît qu'à la fin, introduite par \all{meiner Meinung nach}. Les connecteurs (\all{zuerst, dann, danach, schließlich, einerseits/andererseits, deshalb}) structurent le texte, et le verbe conjugué reste toujours en deuxième position.
\end{aretenir}

\newpage

\coursec{Activité d'intégration}

\coursec{Activité d'intégration : Le conflit d'héritage}

\courssub{Objectifs de l'activité}

À l'issue de cette activité d'intégration, l'élève sera capable de :
\begin{itemize}
\item mobiliser le \cle{vocabulaire thématique} de l'héritage (\all{das Erbe, der Erbe, die Erbin, das Testament, der Streit, das Gesetz, der Notar, gerecht, teilen, erben, streiten}) dans une situation de communication authentique ;
\item employer correctement le \cle{génitif} pour exprimer la possession et le complément du nom (\all{das Haus des Großvaters, der Wille des Verstorbenen}) ;
\item raconter des événements passés en choisissant correctement entre \cle{parfait} et \cle{prétérit}, avec l'auxiliaire approprié (\all{haben} ou \all{sein}) ;
\item rédiger une lettre semi-formelle à un notaire, en respectant la structure de la lettre allemande (adresse, lieu et date, formule d'appel, corps, formule de politesse) ;
\item débattre à l'oral, exprimer son accord ou son désaccord et proposer une solution de manière argumentée.
\end{itemize}

\courssub{Situation de départ}

\textbf{Contexte.} Vous êtes les membres de la famille Mbarga, à Yaoundé. Le grand-père, Herr Emmanuel Mbarga, est décédé le mois dernier. Il laisse derrière lui : une maison à Mvog-Ada, un compte en banque de 8 millions de FCFA et un champ de cacao près d'Obala. Le problème : son testament, écrit à la main, est ambigu. La phrase décisive peut être interprétée de deux façons différentes. La famille est divisée et décide de s'adresser à un notaire, Maître Rechtsanwalt Ndoumbe, pour trancher le conflit.

\begin{textebox}[Das Testament des Großvaters (Auszug)]
„Ich, Emmanuel Mbarga, erkläre hiermit meinen letzten Willen. Das Haus in Mvog-Ada soll an meine Kinder gehen, das Geld auf dem Konto an die Enkel, die mich gepflegt haben. Das Kakaofeld bei Obala gehört dem, der es bearbeitet hat. Ich hoffe, dass meine Familie in Frieden bleibt. Yaoundé, den 12. März 2024. — Emmanuel Mbarga“
\end{textebox}

\textbf{Glossaire :} \all{der letzte Wille} : la dernière volonté ; \all{pflegen} : soigner, prendre soin de ; \all{bearbeiten} : cultiver, travailler (la terre) ; \all{der Friede, -n} : la paix.

\textbf{Pourquoi ce testament est-il ambigu ?} Le grand-père avait trois enfants, mais seulement deux petits-enfants l'ont soigné pendant sa maladie — les autres petits-enfants réclament aussi leur part. Et deux personnes affirment avoir cultivé le champ de cacao : son fils aîné et son neveu.

\courssub{Déroulement de l'activité}

\begin{enumerate}
\item \textbf{Compréhension du document (10 min, à l'oral).} Lisez le testament à voix haute. Répondez en allemand : \all{Was hat der Großvater hinterlassen? Wer sind die möglichen Erben? Warum gibt es einen Streit?} Soulignez dans le texte tous les génitifs et tous les verbes au passé.

\item \textbf{Jeu de rôle : la réunion de famille (15 min, oral).} En groupes de quatre à cinq, répartissez les rôles : le fils aîné (\all{der älteste Sohn}), la fille cadette (\all{die jüngste Tochter}), le petit-fils qui a soigné le grand-père (\all{der Enkelsohn}), la petite-fille qui réclame sa part (\all{die Enkeltochter}), le neveu (\all{der Neffe}). Débattez de la répartition des biens. Utilisez au moins six mots du vocabulaire thématique et trois structures d'opinion : \all{Ich finde, dass … ; Meiner Meinung nach … ; Das ist nicht gerecht, weil … ; Ich schlage vor, dass …}.

\item \textbf{Préparation linguistique (10 min, individuel).} Complétez les phrases suivantes avec le génitif correct : \all{das Haus (der Großvater) ; der Wille (der Verstorbene) ; die Pflege (die Enkel) ; die Lösung (das Problem) ; das Konto (die Familie)}. Puis mettez au parfait : \all{der Großvater / sterben ; wir / das Testament / finden ; die Familie / streiten ; der Neffe / das Feld / bearbeiten}.

\item \textbf{Production écrite : la lettre au notaire (25 min).} Chaque élève rédige une lettre de 10 à 12 phrases à Maître Ndoumbe. La lettre doit : (a) raconter le conflit au parfait et au prétérit ; (b) décrire les biens avec au moins quatre génitifs ; (c) proposer une solution juste (\all{gerecht}) ; (d) respecter la forme de la lettre semi-formelle.

\item \textbf{Lecture et commentaire collectif (10 min).} Les deux meilleures lettres sont lues en classe. La classe identifie les génitifs, les temps du passé et les mots du thème, puis vote pour la solution la plus juste.
\end{enumerate}

\courssub{Ressources linguistiques à mobiliser}

\begin{aretenir}[Lexique thématique : L'héritage]
\begin{center}
\begin{tabularx}{\linewidth}{|X|X|}\hline
\rowcolor{popblueL} \textbf{Deutsch} & \textbf{Français} \\ \hline
das Erbe, -n & l'héritage \\ \hline
der Erbe, -n / die Erbin, -nen & l'héritier / l'héritière \\ \hline
das Testament, -e & le testament \\ \hline
der Verstorbene, -n & le défunt \\ \hline
der Notar, -e & le notaire \\ \hline
der Streit, -e & le conflit, la dispute \\ \hline
das Gesetz, -e & la loi \\ \hline
erben / vererben & hériter / léguer \\ \hline
teilen / aufteilen & partager / répartir \\ \hline
gerecht / ungerecht & juste / injuste \\ \hline
streiten (stritt, hat gestritten) & se disputer \\ \hline
\end{tabularx}
\end{center}
\end{aretenir}

\begin{propriete}[Le génitif : possession et complément du nom]
Le génitif marque l'appartenance ou la relation entre deux noms (complément du nom). En français, il correspond souvent à \emph{de} + nom.
\begin{itemize}
\item \textbf{Forme} : Article défini au génitif + nom (éventuellement avec terminaison).
\item \textbf{Noms masculins et neutres} : ajout de \all{-(e)s} au singulier (\all{des Großvaters, des Hauses, des Problems, des Testaments}).
\item \textbf{Noms féminins et pluriels} : pas de terminaison au nom, seul l'article change (\all{der Familie, der Enkel, der Kinder}).
\item \textbf{Adjectifs} : déclinaison forte après article défini (\all{meines Großvaters, dieses Problems, unserer Meinung}).
\end{itemize}
\all{Das Haus des Großvaters} « la maison du grand-père » ; \all{der Wille des Verstorbenen} « la volonté du défunt ».
\end{propriete}

\begin{propriete}[Le passé du récit : choix entre parfait et prétérit]
\begin{itemize}
\item \textbf{Prétérit (Präteritum)} : temps du récit écrit. Obligatoire pour \all{sein} (\all{war}), \all{haben} (\all{hatte}) et les verbes modaux (\all{wollte, konnte, musste}). Utilisé pour tous les verbes dans un texte narratif formel.
\item \textbf{Parfait (Perfekt)} : temps du récit oral et de la langue courante. Utilisé pour les actions accomplies. Auxiliaire \all{haben} (majorité des verbes) ou \all{sein} (verbes de mouvement \all{gehen, kommen} et changement d'état \all{sterben, einschlafen, aufwachen, werden}).
\item \textbf{Exemples} : \all{Der Großvater starb (Prät.) / ist gestorben (Perf.)} ; \all{Er hinterließ (Prät.) / hat hinterlassen (Perf.) ein Haus} ; \all{Die Familie stritt (Prät.) / hat gestritten (Perf.)}.
\end{itemize}
\end{propriete}

\begin{propriete}[Structure de la lettre semi-formelle]
\begin{enumerate}
\item \textbf{Lieu et date} : en haut à droite. Ex. \all{Yaoundé, den 15. Mai 2025}.
\item \textbf{Formule d'appel} : \all{Sehr geehrter Herr Ndoumbe,} (virgule, puis continuation en minuscule).
\item \textbf{Corps de la lettre} : paragraphes séparés, style soigné, pas de familier.
\item \textbf{Formule de politesse finale} : \all{Mit freundlichen Grüßen} (sans virgule, sans signature au-dessus).
\item \textbf{Signature} : prénom + nom manuscrit.
\end{enumerate}
\end{propriete}

\begin{attention}[Pièges fréquents pour francophones]
\begin{itemize}
\item \textbf{Génitif vs \all{von} + datif} : En langue soutenue et écrite, évitez \all{das Haus von dem Großvater}. Utilisez \all{das Haus des Großvaters}.
\item \textbf{Ordre des mots en subordonnée} : Après \all{weil, dass, ob, wenn}, le verbe conjugué va en \textbf{position finale}. Faux : \all{…, weil der Testament unklar ist}. Correct : \all{…, weil das Testament unklar ist}.
\item \textbf{Majuscules} : Tous les noms prennent une majuscule : \all{das Testament, der Streit, das Gesetz, der Friede, das Erbe}.
\item \textbf{Verbe \all{gehören}} : Il régit le \textbf{datif} (\all{Das Feld gehört dem Neffen}), pas l'accusatif.
\item \textbf{Faux ami} : \all{der Notar} = notaire (juriste), pas \emph{notaire} au sens de \emph{secrétaire} ; \all{das Gesetz} = la loi (texte légal), pas \emph{le règlement} intérieur.
\end{itemize}
\end{attention}

\courssub{Production écrite guidée : Lettre au notaire}

\begin{exoresolu}[Modèle de lettre au notaire]
\textbf{Énoncé :} Rédiger la lettre de Paul Mbarga, petit-fils, au notaire (10 à 12 phrases).
\tcblower
\textbf{Production modèle :}

Yaoundé, den 15. Mai 2025

Sehr geehrter Herr Ndoumbe,

ich schreibe Ihnen wegen des Erbes meines Großvaters Emmanuel Mbarga. Er ist vor einem Monat gestorben und hat ein Haus in Mvog-Ada, ein Konto und ein Kakaofeld bei Obala hinterlassen. Leider war das Testament des Verstorbenen nicht klar, und deshalb hat die Familie gestritten. Mein Onkel wollte das Haus des Großvaters allein behalten, aber das fanden wir ungerecht. Meine Schwester und ich haben unseren Großvater drei Jahre lang gepflegt, deshalb gehört das Geld des Kontos unserer Meinung nach uns. Das Kakaofeld hat mein Onkel zwanzig Jahre bearbeitet, aber der Neffe des Großvaters behauptet, er hat auch dort gearbeitet. Ich schlage vor, dass wir das Haus verkaufen und den Preis unter den Kindern des Verstorbenen gerecht teilen. Die Lösung dieses Problems ist für unsere Familie sehr wichtig, denn der Wille meines Großvaters war der Friede. Ich hoffe auf Ihre Hilfe und danke Ihnen im Voraus.

Mit freundlichen Grüßen

Paul Mbarga

\textbf{Commentaire des choix de langue :}
\begin{itemize}
\item \textbf{Génitifs (6 occurrences)} : \all{des Erbes meines Großvaters} (chaîne génitivale), \all{des Verstorbenen} (adjectif substantivé, n-déclinaison : \all{der Verstorbene, des Verstorbenen}), \all{des Großvaters, des Kontos, unter den Kindern des Verstorbenen, der Wille meines Großvaters, die Lösung dieses Problems}. Tous corrects avec \all{-s} (masc./neutre) ou article seul (fém./plur.).
\item \textbf{Temps du passé} : Prétérit pour \all{sein, wollen, finden} (\all{war, wollte, fanden}) — norme du récit écrit ; Parfait pour les actions ponctuelles (\all{ist gestorben} — aux. \all{sein}, changement d'état ; \all{hat hinterlassen, hat gestritten, haben gepflegt, hat bearbeitet, hat gearbeitet} — aux. \all{haben}).
\item \textbf{Syntaxe subordonnée} : Verbe en position finale après \all{weil, dass} (\all{weil das Testament … nicht klar war} — ici \all{war} final ; \all{dass wir … teilen}). Adverbe en tête + V2 : \all{deshalb hat die Familie gestritten}.
\item \textbf{Lexique du thème} : \all{das Erbe, das Testament, der Verstorbene, streiten, ungerecht, gerecht teilen, der Wille, der Friede, der Notar, erben, pflegen, bearbeiten}.
\item \textbf{Forme épistolaire} : Lieu/date, appel avec virgule/minuscule, formule de politesse sans virgule, signature.
\end{itemize}
\end{exoresolu}

\courssub{Bilan et prolongements}

Cette activité a permis d'intégrer les trois savoir-faire de la leçon dans une situation de communication réaliste et proche du vécu des élèves. À l'oral, le jeu de rôle a entraîné l'expression de l'opinion et du désaccord avec le vocabulaire de l'héritage. À l'écrit, la lettre au notaire a exigé la maîtrise combinée du génitif (possession), du parfait et du prétérit (récit au passé) et de la structure de la lettre allemande. Pour vérifier l'acquisition, l'enseignant s'assure que chaque production contient : au moins quatre génitifs corrects, un choix cohérent entre parfait et prétérit avec les bons auxiliaires, et six mots du lexique thématique. En prolongement interculturel, les élèves pourront comparer le droit successoral allemand (\all{das Erbrecht}, avec la \all{Pflichtteil}, la part réservataire) et les pratiques successorales camerounaises (coutume, droit écrit, rôle de la famille élargie), ce qui ouvre une perspective sur le thème « Vie familiale et sociale » au programme.

\newpage

\coursec{Méthodes et exercices résolus}

\coursec{Les problèmes d'héritage}

\courssub{Texte support : Der Streit um das Erbe}

\begin{textebox}[Der Streit um das Erbe]
Nach dem Tod seines Vaters fand Thomas ein Testament. Darin stand: „Das Haus gehört meinem Sohn, das Geld meiner Tochter.“ Aber seine Schwester Anna war nicht einverstanden. Sie sagte, der Vater sei krank gewesen und habe nicht klar denken können. Die Geschwister stritten monatelang. Schließlich entschied ein Gericht: Das Testament war gültig. Thomas behielt das Haus, Anna bekam das Geld.
\end{textebox}

\courssub{Méthode 1 : Comprendre un texte sur les problèmes d'héritage}

\begin{methode}[Lire et comprendre un texte : démarche en 5 étapes]
\begin{enumerate}
\item \textbf{Lire le titre et la source.} Le titre \all{Der Streit um das Erbe} annonce déjà le thème : un conflit (\all{der Streit}) autour d'un héritage (\all{das Erbe}). Repérer les mots du champ lexical avant même de lire.
\item \textbf{Première lecture globale.} Lire sans dictionnaire. Identifier : \emph{qui ?} (les personnages), \emph{où ?}, \emph{quand ?}, \emph{quel est le problème ?}. Ne pas s'arrêter sur les mots inconnus.
\item \textbf{Deuxième lecture détaillée.} Souligner les mots-clés du thème : \all{das Testament, der Erbe, die Erbin, das Gesetz, die Familie}. Deviner le sens des mots inconnus par le contexte ou par la famille de mots (\all{erben} $\rightarrow$ \all{das Erbe}, \all{der Erbe}).
\item \textbf{Repérer les connecteurs.} \all{zuerst, dann, danach, schließlich, aber, deshalb} structurent le récit et aident à suivre la chronologie.
\item \textbf{Répondre aux questions.} Reformuler avec ses propres mots, en phrases complètes, en reprenant le temps de la question (Präsens ou Perfekt). Ne jamais recopier une phrase entière du texte sans la justifier.
\end{enumerate}
\end{methode}

\begin{exoresolu}[Compréhension de texte — type Bac]
\textbf{Texte :} \all{Nach dem Tod seines Vaters fand Thomas ein Testament. Darin stand: „Das Haus gehört meinem Sohn, das Geld meiner Tochter.“ Aber seine Schwester Anna war nicht einverstanden. Sie sagte, der Vater sei krank gewesen und habe nicht klar denken können. Die Geschwister stritten monatelang. Schließlich entschied ein Gericht: Das Testament war gültig. Thomas behielt das Haus, Anna bekam das Geld.}

\textbf{Questions :} 1. \all{Was fand Thomas nach dem Tod seines Vaters?} 2. \all{Warum war Anna nicht einverstanden?} 3. \all{Wer entschied am Ende über den Streit?}
\tcblower
\textbf{Raisonnement :} La question 1 porte sur un fait explicite (première phrase). La question 2 demande la cause du désaccord : il faut repérer \all{Sie sagte, der Vater sei krank gewesen} — discours indirect au Konjunktiv I, à reformuler. La question 3 vise \all{ein Gericht entschied}.

\textbf{Réponses rédigées :}

1. \all{Nach dem Tod seines Vaters fand Thomas ein Testament.} (Après la mort de son père, Thomas a trouvé un testament.)

2. \all{Anna war nicht einverstanden, weil ihr Vater krank gewesen war und nicht klar hatte denken können.} (Anna n'était pas d'accord parce que son père avait été malade et n'avait pas pu penser clairement.) — Justification : subordonnée de cause avec \all{weil}, verbe conjugué à la fin ; le plus-que-parfait \all{gewesen war} marque l'antériorité.

3. \all{Am Ende entschied ein Gericht über den Streit.} (À la fin, c'est un tribunal qui a tranché le conflit.) — Justification : \all{entscheiden über + Akkusativ} ; inversion sujet-verbe après \all{am Ende} en tête de phrase.

\textbf{Conclusion :} chaque réponse est une phrase complète, au temps du texte (Präteritum/Perfekt), avec le vocabulaire du thème correctement employé.
\end{exoresolu}

\courssub{Méthode 2 : Employer le vocabulaire de l'héritage}

\begin{methode}[Maîtriser le lexique thématique : Erbe, Testament, Gesetz]
\begin{enumerate}
\item \textbf{Apprendre chaque nom avec son article et son pluriel :} \all{das Erbe (-n)}, \all{der Erbe (-n)}, \all{die Erbin (-nen)}, \all{das Testament (-e)}, \all{der Streit (-e)}, \all{das Gesetz (-e)}, \all{der Verstorbene / die Verstorbene (-n)} (le/la défunt(e)), \all{das Gericht (-e)} (le tribunal).
\item \textbf{Mémoriser les verbes associés :} \all{erben} (hériter), \all{vererben} (léguer), \all{hinterlassen} (laisser, en mourant), \all{streiten (stritt, hat gestritten)} (se disputer), \all{ein Testament machen} (faire un testament).
\item \textbf{Construire des familles de mots :} \all{das Erbe} $\rightarrow$ \all{der Erbe / die Erbin} $\rightarrow$ \all{erben} $\rightarrow$ \all{die Erbschaft} ; \all{das Gesetz} $\rightarrow$ \all{gesetzlich} $\rightarrow$ \all{der Gesetzgeber}.
\item \textbf{Utiliser les expressions figées :} \all{ein Erbe antreten} (recueillir un héritage), \all{um das Erbe streiten} (se disputer l'héritage), \all{vor Gericht gehen} (aller au tribunal), \all{das Gesetz sagt / bestimmt} (la loi dispose).
\item \textbf{Vérifier le genre dans la phrase :} le genre commande l'article et l'adjectif (\all{das gültige Testament}, \all{der älteste Erbe}).
\end{enumerate}
\end{methode}

\begin{exoresolu}[Vocabulaire — phrases à trous et production]
\textbf{Énoncé :} Complétez avec le mot qui convient (article + forme correcte) : \all{Erbe, Testament, Gesetz, Erbin, Streit, Gericht}.

1. \all{Vor seinem Tod hat der Großvater ein \_\_\_\_\_ gemacht.}
2. \all{Meine Tante ist die einzige \_\_\_\_\_ des Hauses.}
3. \all{Der \_\_\_\_\_ zwischen den Brüdern dauerte zwei Jahre.}
4. \all{Nach dem deutschen \_\_\_\_\_ erben die Kinder immer einen Teil.}
5. \all{Das \_\_\_\_\_ hat entschieden: Das Testament ist gültig.}
\tcblower
\textbf{Solutions justifiées :}

1. \all{ein Testament} — neutre (\all{das Testament}), accusatif après \all{machen} : \all{ein Testament}. « Avant sa mort, le grand-père a fait un testament. »

2. \all{die einzige Erbin} — il s'agit d'une femme (\all{meine Tante}), donc le féminin \all{die Erbin} et non \all{der Erbe}. « Ma tante est la seule héritière de la maison. »

3. \all{Der Streit} — masculin, nominatif sujet ; \all{der Streit zwischen + Dativ}. « La dispute entre les frères a duré deux ans. »

4. \all{Nach dem deutschen Gesetz} — neutre, datif après \all{nach} : \all{dem Gesetz}. « Selon la loi allemande, les enfants héritent toujours d'une part. »

5. \all{Das Gericht} — neutre, nominatif sujet de \all{hat entschieden}. « Le tribunal a décidé : le testament est valable. »

\textbf{Conclusion :} le choix du mot dépend du sens, mais sa forme dépend du genre et du cas : toujours vérifier les deux.
\end{exoresolu}

\courssub{Grammaire 1 : Le génitif — possession et compléments du nom}

\begin{propriete}[Formes et emplois du génitif]
\begin{enumerate}
\item \textbf{Fonction :} le génitif exprime la possession ou le lien entre deux noms : \all{das Haus \textbf{meines Vaters}} (la maison \emph{de mon père}), \all{der Streit \textbf{der Geschwister}} (le conflit \emph{des frères et sœurs}).
\item \textbf{Terminaisons (singulier) :} noms masculins et neutres ajoutent \textbf{-(e)s} ; noms féminins inchangés.
\begin{center}
\begin{tabular}{|l|l|l|l|}
\hline
\rowcolor{popblueL} & Masculin & Féminin & Neutre \\
\hline
Article défini & des Vater\textbf{s} & der Mutter & des Kind\textbf{es} \\
\hline
Article indéfini & eines Vater\textbf{s} & einer Mutter & eines Kind\textbf{es} \\
\hline
\end{tabular}
\end{center}
\item \textbf{Pluriel :} article \all{der} pour les trois genres (\all{der Eltern, der Geschwister, der Testamente}). Pas de -s au nom au pluriel (sauf pluriel en -n/-s déjà existant).
\item \textbf{Noms propres :} ajout de -s sans apostrophe : \all{Annas Erbe}, \all{Thomas' Haus} (apostrophe seulement si le nom finit par -s, -ß, -z, -x).
\item \textbf{Prépositions régissant le génitif :} \all{während, wegen, trotz, statt, anlässlich} : \all{wegen des Erbes} (à cause de l'héritage).
\item \textbf{Alternative orale :} \all{von + Dativ} (\all{das Haus von meinem Vater}) est courante à l'oral, mais à l'écrit (Bac), le génitif est la norme.
\end{enumerate}
\end{propriete}

\begin{methode}[Traduire un « de » français par le génitif allemand]
\begin{enumerate}
\item Repérer le nom « possesseur » (celui qui suit « de »).
\item Déterminer son genre en allemand (masc., fém., neutre).
\item Mettre l'article (défini/indéfini/possessif) au génitif : \all{des, eines, meines} (masc./neutre) ; \all{der, einer, meiner} (fém.).
\item Ajouter \textbf{-(e)s} au nom si masculin ou neutre singulier (\all{Vaters, Kindes, Hauses, Testaments}).
\item Placer le groupe génitif \textbf{avant} le nom possédé : \all{des Vaters Haus} (mais ordre usuel : \all{das Haus des Vaters}).
\end{enumerate}
\end{methode}

\begin{exoresolu}[Thème — le génitif en traduction]
\textbf{Énoncé :} Traduisez en allemand.

1. La maison de mon grand-père est très vieille.
2. Le testament du défunt était clair.
3. À cause de l'héritage, les enfants de la voisine se sont disputés.
4. La décision du tribunal a surpris la famille de Thomas.
\tcblower
\textbf{Solutions détaillées :}

1. \all{Das Haus meines Großvaters ist sehr alt.} — \all{der Großvater} masculin : génitif \all{meines Großvaters} (article \all{meines} + -s au nom). Piège : ne pas écrire \all{von meinem Großvater} à l'écrit.

2. \all{Das Testament des Verstorbenen war klar.} — \all{der Verstorbene} est un adjectif nominalisé (déclinaison faible) : génitif \all{des Verstorbenen} (et non \all{des Verstorbenes}).

3. \all{Wegen des Erbes haben sich die Kinder der Nachbarin gestritten.} — \all{wegen + Genitiv} : \all{des Erbes} (\all{das Erbe}, neutre, + -s). \all{die Kinder der Nachbarin} : génitif féminin sans -s. Verbe réfléchi \all{sich streiten}, Perfekt avec \all{haben}.

4. \all{Die Entscheidung des Gerichts hat die Familie von Thomas überrascht.} — \all{das Gericht} : génitif \all{des Gerichts}. Pour le nom propre, deux options : \all{Thomas' Familie} (génitif) ou \all{die Familie von Thomas} (von + Dativ) ; les deux sont acceptés à l'écrit pour les noms propres.

\textbf{Conclusion :} trois réflexes : repérer le « de » français, déterminer le genre du nom allemand, appliquer la marque -(e)s uniquement au masculin et au neutre singulier.
\end{exoresolu}

\courssub{Grammaire 2 : Parfait et prétérit pour raconter}

\begin{propriete}[Emploi et formation du Perfekt et du Präteritum]
\begin{enumerate}
\item \textbf{Partage des rôles :} le \cle{Perfekt} s'emploie surtout à l'oral et dans les récits personnels ; le \cle{Präteritum} domine dans les textes écrits (articles, littérature, commentaire de Bac). Au Bac, dans une production écrite, le Präteritum est privilégié.
\item \textbf{Exceptions majeures (quasi toujours au Präteritum) :} \all{sein, haben, werden} et les verbes modaux (\all{können, müssen, wollen, sollen, dürfen, mögen}) : \all{er war krank}, \all{sie hatte ein Testament}, \all{er konnte nicht klar denken}.
\item \textbf{Perfekt :} auxiliaire \all{haben} ou \all{sein} + Participe II (à la fin de la phrase).
\begin{itemize}
\item \all{sein} : verbes de déplacement (\all{gehen, kommen, laufen}) ou changement d'état (\all{sterben, aufwachen, einschlafen, werden}) : \all{Der Vater \textbf{ist} gestorben.}
\item \all{haben} : tous les autres verbes (transitifs, réflexifs, statifs) : \all{Thomas \textbf{hat} das Testament \textbf{gefunden}.}
\end{itemize}
\item \textbf{Präteritum des verbes faibles (réguliers) :} radical + \textbf{-te} + terminaisons : \all{machen $\rightarrow$ ich machte}, \all{sagen $\rightarrow$ ich sagte}, \all{arbeiten $\rightarrow$ ich arbeitete}.
\item \textbf{Verbes forts (irréguliers) — formes à connaître (niveau Tle) :}
\begin{center}
\begin{tabular}{|l|l|l|l|}
\hline
\rowcolor{popblueL} Infinitif & Präteritum (ich/er/sie/es) & Participe II & Auxiliaire \\
\hline
\all{finden} & \all{fand} & \all{gefunden} & \all{haben} \\
\all{geben} & \all{gab} & \all{gegeben} & \all{haben} \\
\all{sterben} & \all{starb} & \all{gestorben} & \all{sein} \\
\all{hinterlassen} & \all{hinterließ} & \all{hinterlassen} & \all{haben} \\
\all{streiten} & \all{stritt} & \all{gestritten} & \all{haben} \\
\all{entscheiden} & \all{entschied} & \all{entschieden} & \all{haben} \\
\all{behalten} & \all{behielt} & \all{behalten} & \all{haben} \\
\all{bekommen} & \all{bekam} & \all{bekommen} & \all{haben} \\
\hline
\end{tabular}
\end{center}
\item \textbf{Cohérence du récit :} une fois le temps choisi (Präteritum pour l'écrit), ne pas alterner sans raison. Le \cle{Plusquamperfekt} (\all{hatte gefunden / war gestorben}) marque l'antériorité par rapport au Präteritum.
\end{enumerate}
\end{propriete}

\begin{methode}[Transformer un récit oral (Perfekt) en récit écrit (Präteritum)]
\begin{enumerate}
\item Identifier chaque verbe conjugué au Perfekt.
\item Déterminer s'il est faible (rajoute -te) ou fort (changement de voyelle, tableau ci-dessus).
\item Remplacer l'auxiliaire + Participe II par la forme simple du Präteritum.
\item Vérifier les verbes modaux/\all{sein}/\all{haben} : ils restent au Präteritum (\all{war, hatte, konnte}).
\item Conserver l'ordre des mots (verbe en 2\textsuperscript{e} position, Participe II disparaît).
\end{enumerate}
\end{methode}

\begin{exoresolu}[Version et transformation — les temps du récit]
\textbf{Énoncé A :} Mettez le texte suivant au Präteritum (récit écrit) :

\all{Mein Onkel ist letztes Jahr gestorben. Er hat uns ein großes Haus hinterlassen. Meine Brüder und ich haben uns zuerst gestritten, aber dann haben wir eine Lösung gefunden.}

\textbf{Énoncé B :} Traduisez : « Mon grand-père est mort en 2020. Il avait fait un testament, mais personne ne l'a trouvé. »
\tcblower
\textbf{Solution A :}

\all{Mein Onkel \textbf{starb} letztes Jahr. Er \textbf{hinterließ} uns ein großes Haus. Meine Brüder und ich \textbf{stritten} uns zuerst, aber dann \textbf{fanden} wir eine Lösung.}

Justifications : \all{sterben – starb} (fort, aux. \all{sein}) ; \all{hinterlassen – hinterließ} (fort, aux. \all{haben}) ; \all{sich streiten – stritten} (fort, pluriel) ; \all{finden – fanden} (fort). Dans un récit écrit, on remplace tout le Perfekt par le Präteritum.

\textbf{Solution B :}

\all{Mein Großvater \textbf{ist} 2020 \textbf{gestorben}. Er \textbf{hatte} ein Testament \textbf{gemacht}, aber niemand \textbf{hat} es \textbf{gefunden}.}

Justifications : \all{sterben} = changement d'état $\rightarrow$ Perfekt avec \all{sein} : \all{ist gestorben}. « Avait fait » = antériorité $\rightarrow$ Plusquamperfekt : \all{hatte gemacht} (aux. \all{haben} au Präteritum + Part. II). « Personne ne l'a trouvé » : \all{niemand} + verbe 3\textsuperscript{e} pers. sg. ; « le » = \all{das Testament} $\rightarrow$ \all{es} ; Perfekt \all{hat gefunden} (récit oral/standard) ou Präteritum \all{fand} (récit écrit). Les deux sont acceptables selon le registre visé ; ici Perfekt car phrase isolée.

\textbf{Conclusion :} vérifier pour chaque verbe : auxiliaire (\all{haben} ou \all{sein}), forme du participe ou du Präteritum, et place du verbe conjugué en deuxième position.
\end{exoresolu}

\courssub{Expression guidée : Résumé et production écrite}

\begin{methode}[Résumer et rédiger une production guidée (Zusammenfassung \& Stellungnahme)]
\begin{enumerate}
\item \textbf{Résumé (Zusammenfassung) :} garder uniquement les informations essentielles (personnages, problème, dénouement) ; 3 à 5 phrases ; utiliser ses propres mots ; débuter par \all{In dem Text geht es um...} (Le texte traite de...). Temps : Présent (standard) ou Perfekt.
\item \textbf{Opinion personnelle (Stellungnahme) :} amorces : \all{Meiner Meinung nach...}, \all{Ich finde es wichtig / ungerecht, dass...}, \all{Man sollte...}, \all{Es ist besser, ... zu...}.
\item \textbf{Structure :} Introduction (1 phrase) $\rightarrow$ Arguments (2-3 phrases avec connecteurs \all{zuerst, außerdem, deshalb, aber}) $\rightarrow$ Conclusion (\all{Zum Schluss...}, \all{Insgesamt...}).
\item \textbf{Auto-correction avant de rendre :} verbe conjugué en 2\textsuperscript{e} position ; verbe à la fin dans les subordonnées (\all{weil, dass, obwohl}) ; majuscule à \textbf{tous} les noms ; génitifs corrects (article + -(e)s masc./neutre) ; accords adjectifs.
\end{enumerate}
\end{methode}

\begin{exoresolu}[Production écrite guidée — type Bac]
\textbf{Énoncé :} Fassen Sie den Text „Der Streit um das Erbe“ in 3-4 Sätzen zusammen und geben Sie Ihre Meinung: Sollen Kinder immer gleich viel erben? (environ 60 mots)
\tcblower
\textbf{Démarche :} 1) Relever les faits : père mort, testament partage maison/argent, sœur conteste, tribunal valide. 2) Rédiger le résumé au Présent. 3) Argumenter avec \all{meiner Meinung nach} + subordonnée \all{weil/dass}.

\textbf{Production modèle :}

\all{In dem Text geht es um einen Streit zwischen zwei Geschwistern. Nach dem Tod des Vaters findet Thomas ein Testament, aber seine Schwester Anna akzeptiert es nicht. Am Ende entscheidet ein Gericht, dass das Testament gültig ist. Meiner Meinung nach sollen Kinder immer gleich viel erben, weil Ungleichheit oft Konflikte in der Familie verursacht. Man sollte auch früh mit der Familie über das Erbe sprechen.}

(Le texte traite d'un conflit entre deux frère et sœur. Après la mort du père, Thomas trouve un testament, mais sa sœur Anna ne l'accepte pas. Finalement, un tribunal décide que le testament est valable. À mon avis, les enfants devraient toujours hériter de parts égales, parce que l'inégalité provoque souvent des conflits dans la famille. On devrait aussi parler tôt de l'héritage avec sa famille.)

\textbf{Points de grammaire vérifiés :} génitif \all{nach dem Tod des Vaters} ; verbe à la fin après \all{dass} et \all{weil} ; \all{gleich viel erben} ; majuscules aux noms (\all{Streit, Geschwistern, Tod, Vaters, Testament, Schwester, Gericht, Meinung, Kinder, Ungleichheit, Konflikte, Familie, Erbe}).
\end{exoresolu}

\begin{savaistu}[Droit successoral : Allemagne vs Cameroun]
En Allemagne, le Code civil (\all{Bürgerliches Gesetzbuch - BGB}) régit la succession. Principe clé : la \cle{réserve héréditaire} (\all{Pflichtteil}). Même déshérités par testament, les enfants (et le conjoint) ont droit à la moitié de leur part légale (\all{gesetzlicher Erbteil}). Cela limite la liberté de tester (\all{Testierfreiheit}). Au Cameroun, la coutume prime souvent : l'héritage suit la lignée (matrilinéaire ou patrilinéaire), et le testament écrit est rare. Le Code civil camerounais (inspiré du droit français) s'applique aux successions « modernes », mais les conflits entre droit coutumier et droit écrit sont fréquents. En Allemagne, l'État (\all{Nachlassgericht}) surveille la procédure ; au Cameroun, le tribunal de grande instance ou le chef de village peuvent être saisis.
\end{savaistu}

\begin{attention}[Les 5 erreurs qui coûtent des points au Bac]
\begin{enumerate}
\item \textbf{Oublier la majuscule aux noms :} \all{das erbe}, \all{das testament} sont fautifs. Tout nom allemand prend une majuscule : \all{das Erbe, das Testament, das Gesetz}. C'est la faute la plus sanctionnée.
\item \textbf{Confondre \all{der Erbe} et \all{das Erbe} :} \all{der Erbe} = l'héritier (la personne), \all{das Erbe} = l'héritage (la chose). Faux ami supplémentaire : \all{das Gesetz} = la loi (pas « le gazet ») ; \all{das Gericht} = le tribunal \emph{et} le plat (contexte !).
\item \textbf{Génitif sans marque :} écrire \all{das Haus mein Vater} ou \all{das Haus meinem Vater} (datif !) au lieu de \all{das Haus meines Vaters}. Masculin et neutre singulier : article en \all{-es} + nom en \all{-(e)s}.
\item \textbf{Mauvais auxiliaire au Perfekt :} \all{Der Vater hat gestorben} est faux : \all{sterben} exprime un changement d'état $\rightarrow$ \all{Der Vater \textbf{ist} gestorben}. De même : \all{ist gegangen, ist gekommen, ist aufgewacht}.
\item \textbf{Place du verbe calquée du français :} \all{Weil der Vater war krank} est un calque. En subordonnée, le verbe conjugué va à la fin : \all{weil der Vater krank \textbf{war}}. Et après un complément en tête de phrase principale : \all{Nach dem Tod des Vaters \textbf{fand} Thomas ein Testament} (verbe en 2\textsuperscript{e} position, sujet après).
\end{enumerate}
\end{attention}

\newpage

\coursec{Exercices}

\courssub{Je vérifie mes connaissances}

\begin{exercice}[Compréhension globale : QCM]
Lisez le texte \og Der Streit um das Erbe \fg{} (page précédente) et choisissez la bonne réponse.
\begin{enumerate}
    \item Quel est le sujet principal du conflit familial ?
    \begin{itemize}
        \item[a)] Le partage de l'argent liquide.
        \item[b)] La maison familiale et le testament.
        \item[c)] Les bijoux de la grand-mère.
        \item[d)] Une dette de jeu du père.
    \end{itemize}
    \item Quel rôle joue le notaire (der Notar) dans l'histoire ?
    \begin{itemize}
        \item[a)] Il prend parti pour le fils aîné.
        \item[b)] Il vend la maison immédiatement.
        \item[c)] Il médie et explique la loi.
        \item[d)] Il détruit le testament.
    \end{itemize}
    \item Comment se termine le conflit ?
    \begin{itemize}
        \item[a)] Par un procès long et coûteux.
        \item[b)] Par une vente de la maison et un partage équitable.
        \item[c)] Le fils cadet hérite de tout.
        \item[d)] La famille se sépare sans parler.
    \end{itemize}
    \item Quel temps verbal domine dans la partie narrative du texte ?
    \begin{itemize}
        \item[a)] Le présent.
        \item[b)] Le parfait (Perfekt).
        \item[c)] Le prétérit (Präteritum).
        \item[d)] Le futur.
    \end{itemize}
\end{enumerate}
\end{exercice}

\begin{exercice}[Compréhension détaillée : Vrai ou Faux ? Justifiez par une phrase du texte]
Déterminez si les affirmations sont vraies (Richtig) ou fausses (Falsch). Citez le passage du texte qui justifie votre choix.
\begin{enumerate}
    \item Le grand-père a laissé un testament clair et sans ambiguïté.
    \item La fille aînée (die älteste Tochter) habite encore dans la maison familiale.
    \item Les héritiers ont dû vendre la maison parce qu'ils ne pouvaient pas s'entendre.
    \item Le notaire a décidé que le fils cadet recevrait la plus grande part.
    \item La loi allemande (das deutsche Erbrecht) prévoit une part obligatoire (Pflichtteil) pour les enfants.
\end{enumerate}
\end{exercice}

\begin{exercice}[Lexique : Familles de mots et articles]
Complétez le tableau suivant. Donnez l'article défini, le pluriel (avec Umlaut si nécessaire) et la traduction française. Pour les verbes, donnez l'infinitif, le Partizip II et le Präteritum (ich-Form).
\begin{center}
\begin{tabularx}{\linewidth}{|X|X|X|X|}
\hline
\rowcolor{popblueL} \textbf{Mot de base} & \textbf{Formes grammaticales} & \textbf{Dérivés / Composés} & \textbf{Traduction} \\
\hline
\all{der Erbe} & Pluriel : \dots & die Erbin, das Erbe, vererben & \dots \\
\hline
\all{das Testament} & Pluriel : \dots & testieren, testamentarisch & \dots \\
\hline
\all{der Streit} & Pluriel : \dots & streiten, streitsüchtig, der Erbstreit & \dots \\
\hline
\all{der Notar} & Pluriel : \dots & notariell, die Notarin, beurkunden & \dots \\
\hline
\all{vererben} & Partizip II: \dots ; Präteritum: \dots & die Vererbung, der Vererbungsweg & \dots \\
\hline
\all{regeln} & Partizip II: \dots ; Präteritum: \dots & die Regelung, die Erbregelung & \dots \\
\hline
\all{teilen} & Partizip II: \dots ; Präteritum: \dots & der Teil, die Teilung, aufteilen & \dots \\
\hline
\all{der Pflichtteil} & Pluriel : \dots & pflichtteilsberechtigt, enterben & \dots \\
\hline
\end{tabularx}
\end{center}
\end{exercice}

\begin{exercice}[Conjugaison : Perfekt et Präteritum — Verbes du thème]
Conjuguez les verbes entre parenthèses soit au \textbf{Perfekt} (avec l'auxiliaire correct \all{haben} ou \all{sein}), soit au \textbf{Präteritum}, selon le temps indiqué.
\begin{enumerate}
    \item (Perfekt) Der Großvater \dots\ (sterben) vor zwei Jahren.
    \item (Perfekt) Er \dots\ (vererben) das Haus seiner Enkelin.
    \item (Perfekt) Die Geschwister \dots\ (kommen) nicht zu einer Einigung.
    \item (Präteritum) Der Notar \dots\ (machen) einen Vorschlag.
    \item (Präteritum) Die Erben \dots\ (warten) lange auf die Entscheidung.
    \item (Perfekt) Schließlich \dots\ sie das Haus \dots\ (verkaufen).
    \item (Präteritum) Der Streit \dots\ (werden) beigelegt.
    \item (Perfekt) Wir \dots\ (arbeiten) mit einem Anwalt zusammen.
\end{enumerate}
\end{exercice}

\courssub{Je m'entraîne}

\begin{exercice}[Grammaire : Le génitif — Compléments du nom]
Mettez les groupes nominaux entre parenthèses au \textbf{génitif} (article + nom + terminaison si nécessaire). Attention aux noms en \emph{-(e)n} (n-Deklination) et aux noms forts (terminaison \emph{-(e)s}).
\begin{enumerate}
    \item Das Testament \dots\ (der Großvater) war unklar.
    \item Die Ansprüche \dots\ (die Erbin) wurden geprüft.
    \item Der Notar erklärte die Regelung \dots\ (die Erbschaft).
    \item Wegen \dots\ (der Streit) ging die Familie vor Gericht.
    \item Das Auto \dots\ (mein Vater) wurde verkauft.
    \item Die Kinder \dots\ (der Verstorbene) haben ein Anrecht auf den Pflichtteil.
    \item Trotz \dots\ (der Widerspruch) wurde das Haus versteigert.
    \item Die Unterschrift \dots\ (der Zeuge) fehlte auf dem Dokument.
\end{enumerate}
\end{exercice}

\begin{exercice}[Traduction thématique : Français $\rightarrow$ Allemand]
Traduisez les phrases suivantes en allemand. Utilisez le vocabulaire de la leçon (Erbe, Testament, Pflichtteil, Notar, vererben, regeln, streiten).
\begin{enumerate}
    \item Le grand-père a légué la maison à sa petite-fille.
    \item Le testament n'était pas valide (gültig) car il manquait la signature.
    \item Les héritiers se disputent l'héritage depuis des mois.
    \item Selon la loi, les enfants ont droit à la réserve héréditaire (Pflichtteil).
    \item Le notaire a réglé la succession (die Nachlassregelung) équitablement.
    \item Nous devons partager l'argent entre les trois frères.
    \item Après la mort du père, la famille s'est réunie chez le notaire.
\end{enumerate}
\end{exercice}

\begin{exercice}[Transformation : Perfekt $\leftrightarrow$ Präteritum (Récit)]
\textbf{A.} Reécrivez le dialogue suivant (oral / Perfekt) sous forme d'un \textbf{récit au Präteritum} (écrit littéraire). Adaptez les pronoms et l'ordre des mots si nécessaire.
\begin{textebox}[Dialogue chez le notaire]
\all{-- Was haben Sie entschieden, Herr Müller?}
\all{-- Ich habe das Testament geprüft. Der Großvater hat das Haus seiner Enkelin vererbt.}
\all{-- Haben die anderen Erben zugestimmt?}
\all{-- Nein, sie haben gestritten. Aber der Notar hat vermittelt. Schließlich haben sie das Haus verkauft.}
\all{-- Und haben sie das Geld geteilt?}
\all{-- Ja, jeder hat seinen Anteil bekommen.}
\end{textebox}

\textbf{B.} Transformez ce récit au Prétérit en \textbf{compte rendu oral au Perfekt} (style parlé, \all{haben/sein} + Partizip II).
\begin{textebox}[Récit]
\all{Der Erbstreit begann nach der Beerdigung. Der älteste Sohn öffnete das Testament. Er las die Klauseln vor. Die Schwester wurde wütend. Sie forderte ihren Pflichtteil. Der Notar kam hinzu. Er erklärte das Gesetz. Die Geschwister einigten sich. Sie verkauften das Grundstück.}
\end{textebox}
\end{exercice}

\begin{exercice}[Texte à trous : Grammaire et vocabulaire mixtes]
Complétez le texte suivant avec les mots manquants (articles, prépositions, terminaisons de cas, formes verbales au Perfekt ou Präteritum).
\begin{textebox}[Die Nachlassregelung]
Nach dem Tod \dots\ (der Großvater, Genitiv) fand die Familie ein altes Testament. Darin \dots\ (stehen, Präteritum), dass das Haus \dots\ (die Enkelin, Dativ) gehöre. Die beiden Söhne \dots\ (sein, Präteritum) empört. Sie \dots\ (können, Präteritum) nicht glauben, dass ihr Vater sie \dots\ (enterben, Partizip II) habe. Ein Anwalt \dots\ (müssen, Präteritum) die Lage prüfen. Er \dots\ (erklären, Perfekt) den Söhnen, dass das Testament formell richtig \dots\ (sein, Perfekt). Aber: \dots\ (wegen, Präposition + Genitiv) dem Gesetz \dots\ (haben, Perfekt) die Kinder Anspruch auf den Pflichtteil. Schließlich \dots\ (einigen, Perfekt, refl.) sich alle Parteien. Das Haus \dots\ (werden, Perfekt, Passiversatz mit \all{werden} oder Aktiv: \all{man verkauft}) verkauft und der Erlös \dots\ (teilen, Perfekt, Passiv).
\end{textebox}
\end{exercice}

\courssub{Je me prépare au Bac}

\begin{exercice}[Compréhension et expression : Texte inédit type Bac]
Lisez attentivement le texte suivant puis répondez \textbf{en allemand} et \textbf{en phrases complètes} aux questions.
\begin{textebox}[Der letzte Wille von Herrn Keller]
Herr Keller war ein reicher Kaufmann aus München. Er hatte drei Kinder: zwei Söhne und eine Tochter. Vor seinem Tod schrieb er ein handschriftliches Testament. Darin bestimmte er: „Mein ältester Sohn Thomas bekommt das Geschäftshaus in der Innenstadt. Meine Tochter Anna erhält das Ferienhaus am See. Mein jüngster Sohn Peter bekommt das Bargeld und die Aktien.“

Nach der Beerdigung trafen sich die Geschwister beim Notar. Thomas war zufrieden, aber Anna und Peter protestierten. „Das ist ungerecht!“, sagte Anna. „Das Geschäftshaus ist viel mehr wert als das Ferienhaus.“ Peter fügte hinzu: „Und ich muss Steuern für die Aktien zahlen. Der Wert schwankt ständig.“

Der Notar, Herr Dr. Weber, beruhigte sie. Er erklärte das Pflichtteilsrecht. „Nach § 2303 BGB steht jedem Kind die Hälfte des gesetzlichen Erbteils zu. Wenn der Wert der Sachen sehr unterschiedlich ist, muss ein Ausgleich gezahlt werden.“ Nach langen Verhandlungen einigten sie sich: Thomas zahlte einen Ausgleich an seine Geschwister. Der Streit war beigelegt, aber das Vertrauen unter den Geschwistern blieb beschädigt.
\end{textebox}

\textbf{Fragen de compréhension (en allemand, phrases complètes) :}
\begin{enumerate}
    \item Was hatte Herr Keller vor seinem Tod verfasst und was war der Inhalt für den ältesten Sohn Thomas?
    \item Warum waren Anna und Peter mit dem Testament nicht einverstanden? Nennen Sie zwei Gründe.
    \item Welche rechtliche Regelung erklärte der Notar Dr. Weber den Geschwistern? (Zitieren Sie den Paragraphen).
    \item Wie wurde der Konflikt schließlich gelöst?
    \item Welche Folge hatte der Erbstreit für das Verhältnis der Geschwister zueinander?
\end{enumerate}

\textbf{Questions de langue (sur le texte) :}
\begin{enumerate}
    \setcounter{enumi}{5}
    \item Relevez dans le texte \textbf{trois verbes au Präteritum} (forme \emph{ich/er/sie/es}) et donnez leur infinitif.
    \item Relevez dans le texte \textbf{deux verbes au Perfekt} (auxiliaire + Partizip II) et donnez leur infinitif.
    \item Mettez au \textbf{génitif} : \og der Erbteil \dots\ (Anna) \fg{} ; \og das Testament \dots\ (Herr Keller) \fg{} ; \og die Entscheidung \dots\ (der Notar) \fg{}.
    \item Transformez la phrase directe en discours indirect (Konjunktiv I ou Präteritum selon le contexte) : \all{Anna sagte: „Das ist ungerecht!“} $\rightarrow$ Anna sagte, \dots
\end{enumerate}
\end{exercice}

\begin{exercice}[Thème et Version : Traduction Bac]
\textbf{A. Version (Allemand $\rightarrow$ Français) :}
Traduisez en français le passage suivant (lignes 1 à 9 du texte ci-dessus, jusqu'à \og \dots schwankt ständig.\fg{}).
\begin{textebox}[Extrait à traduire]
\all{Herr Keller war ein reicher Kaufmann aus München. Er hatte drei Kinder: zwei Söhne und eine Tochter. Vor seinem Tod schrieb er ein handschriftliches Testament. Darin bestimmte er: „Mein ältester Sohn Thomas bekommt das Geschäftshaus in der Innenstadt. Meine Tochter Anna erhält das Ferienhaus am See. Mein jüngster Sohn Peter bekommt das Bargeld und die Aktien.“ Nach der Beerdigung trafen sich die Geschwister beim Notar. Thomas war zufrieden, aber Anna und Peter protestierten. „Das ist ungerecht!“, sagte Anna. „Das Geschäftshaus ist viel mehr wert als das Ferienhaus.“ Peter fügte hinzu: „Und ich muss Steuern für die Aktien zahlen. Der Wert schwankt ständig.“}
\end{textebox}

\textbf{B. Thème (Français $\rightarrow$ Allemand) :}
Traduisez en allemand le paragraphe suivant. Soignez l'ordre des mots, les cas (génitif !) et les temps (Präteritet pour le récit).
\begin{quote}
« Le testament du grand-père n'était pas clair. Les héritiers se sont disputé la maison de la famille. Finalement, un notaire a réglé la succession. Il a expliqué la loi aux enfants. La maison a été vendue et l'argent a été partagé équitablement. »
\end{quote}
\end{exercice}

\begin{exercice}[Rédaction : Sujet type Bac — Expression écrite]
\textbf{Sujet :} \all{Ein Erbstreit in Ihrer Familie oder in Ihrer Nachbarschaft : erzählen Sie und geben Sie Ihre Meinung.}

\textbf{Consignes :}
\begin{itemize}
    \item Rédigez un texte cohérent de \textbf{80 à 100 mots} environ.
    \item Utilisez le \textbf{Präteritum} pour le récit des faits passés et le \textbf{Présent/Perfekt} pour votre opinion actuelle.
    \item Intégrez au moins \textbf{5 mots de vocabulaire} de la leçon (ex: \all{der Erbe, das Testament, der Pflichtteil, der Notar, vererben, streiten, regeln, enterben}).
    \item Structurez votre texte : Introduction (contexte) -- Récit (faits) -- Conséquences/Résolution -- Opinion personnelle.
\end{itemize}

\textbf{Aide au démarrage (optionnelle) :}
\all{In meiner Nachbarschaft gab es vor zwei Jahren einen großen Erbstreit. Als der alte Herr Meier starb, \dots}
\end{exercice}

\newpage

\coursec{Corrigés des exercices}

\coursec{Corrigés des exercices — Leçon AL02 : Les problèmes d'héritage}

\begin{corrige}[Exercice 1 — Compréhension globale : QCM]
\begin{enumerate}
    \item \textbf{Réponse b)} — Le conflit porte sur la maison familiale et le testament du grand-père, dont les clauses étaient ambiguës.
    \item \textbf{Réponse c)} — Le notaire joue le rôle de médiateur : il explique la loi (notamment le \all{Pflichtteil}) et propose une solution équitable, sans prendre parti.
    \item \textbf{Réponse b)} — Le conflit se termine par la vente de la maison et le partage équitable du prix de vente entre les héritiers.
    \item \textbf{Réponse c)} — La partie narrative est rédigée au \all{Präteritum} (\all{starb, fand, erklärte, stritten, verkauften}), temps du récit écrit ; le \all{Perfekt} apparaît surtout dans les dialogues.
\end{enumerate}
\textbf{Méthode :} au QCM de compréhension, éliminez d'abord les réponses qui contredisent explicitement le texte (a et d à la question 1), puis vérifiez la réponse retenue par une citation précise.
\end{corrige}

\begin{corrige}[Exercice 2 — Compréhension détaillée : Vrai ou Faux]
\begin{enumerate}
    \item \textbf{Falsch.} Le testament était ambigu. Justification : \all{„Das Testament des Großvaters war unklar und führte zu einem langen Streit.“} (Le testament du grand-père était peu clair et a provoqué une longue dispute.)
    \item \textbf{Richtig.} Justification : \all{„Die älteste Tochter wohnte noch in dem Haus und wollte es behalten.“} (La fille aînée habitait encore dans la maison et voulait la garder.)
    \item \textbf{Richtig.} Justification : \all{„Da sich die Erben nicht einigen konnten, wurde das Haus schließlich verkauft und der Erlös geteilt.“} (Comme les héritiers ne pouvaient pas se mettre d'accord, la maison a finalement été vendue et le produit de la vente partagé.)
    \item \textbf{Falsch.} Le notaire n'a favorisé personne ; il a appliqué la loi. Justification : \all{„Der Notar erklärte, dass alle Kinder nach dem Gesetz gleichberechtigt sind.“} (Le notaire a expliqué que tous les enfants ont les mêmes droits selon la loi.)
    \item \textbf{Richtig.} Justification : \all{„Nach dem deutschen Erbrecht steht den Kindern ein Pflichtteil zu.“} (Selon le droit successoral allemand, les enfants ont droit à une réserve héréditaire.)
\end{enumerate}
\textbf{Point de vigilance :} la justification doit être une \emph{citation exacte} du texte, entre guillemets allemands „…“, et non une paraphrase de mémoire.
\end{corrige}

\begin{corrige}[Exercice 3 — Lexique : Familles de mots et articles]
\begin{center}
\begin{tabularx}{\linewidth}{|X|X|X|X|}
\hline
\rowcolor{popblueL} \textbf{Mot de base} & \textbf{Formes grammaticales} & \textbf{Dérivés / Composés} & \textbf{Traduction} \\
\hline
\all{der Erbe} & Pluriel : \all{die Erben} & \all{die Erbin, das Erbe, vererben} & l'héritier \\
\hline
\all{das Testament} & Pluriel : \all{die Testamente} & \all{testieren, testamentarisch} & le testament \\
\hline
\all{der Streit} & Pluriel : \all{die Streitigkeiten} & \all{streiten, streitsüchtig, der Erbstreit} & la dispute, le conflit \\
\hline
\all{der Notar} & Pluriel : \all{die Notare} & \all{notariell, die Notarin, beurkunden} & le notaire \\
\hline
\all{vererben} & Partizip II : \all{vererbt} ; Präteritum : \all{vererbte} & \all{die Vererbung} & léguer, transmettre \\
\hline
\all{regeln} & Partizip II : \all{geregelt} ; Präteritum : \all{regelte} & \all{die Regelung, die Erbregelung} & régler, organiser \\
\hline
\all{teilen} & Partizip II : \all{geteilt} ; Präteritum : \all{teilte} & \all{der Teil, die Teilung, aufteilen} & partager, diviser \\
\hline
\all{der Pflichtteil} & Pluriel : \all{die Pflichtteile} & \all{pflichtteilsberechtigt, enterben} & la réserve héréditaire \\
\hline
\end{tabularx}
\end{center}
\textbf{Remarques :}
\begin{itemize}
    \item \all{der Erbe} suit la \emph{n-Deklination} : \all{des Erben, dem Erben, den Erben} ; pluriel \all{die Erben}. Ne pas confondre avec \all{das Erbe} (l'héritage, la chose transmise).
    \item \all{vererben} est un verbe faible à particule \emph{inséparable} : \all{vererbt} sans \all{ge-}. Même logique pour \all{enterben} $\rightarrow$ \all{enterbt}.
    \item \all{streiten} est fort : \all{stritt, hat gestritten}.
\end{itemize}
\end{corrige}

\begin{corrige}[Exercice 4 — Conjugaison : Perfekt et Präteritum]
\begin{enumerate}
    \item \all{Der Großvater \textbf{ist vor zwei Jahren gestorben}.} — \all{sterben} exprime un changement d'état $\rightarrow$ auxiliaire \all{sein}. (Ordre : temps avant verbe en fin de phrase).
    \item \all{Er \textbf{hat} das Haus seiner Enkelin \textbf{vererbt}.} — Verbe faible à particule inséparable : pas de \all{ge-}.
    \item \all{Die Geschwister \textbf{sind} nicht zu einer Einigung \textbf{gekommen}.} — \all{kommen} (mouvement / aboutissement) $\rightarrow$ auxiliaire \all{sein}.
    \item \all{Der Notar \textbf{machte} einen Vorschlag.} — Verbe faible : radical + \all{-te}.
    \item \all{Die Erben \textbf{warteten} lange auf die Entscheidung.} — Radical en \all{-t} $\rightarrow$ \all{-eten} (euphonie).
    \item \all{Schließlich \textbf{haben} sie das Haus \textbf{verkauft}.} — Verbe faible, auxiliaire \all{haben} (verbe transitif).
    \item \all{Der Streit \textbf{wurde} beigelegt.} — \all{werden} est fort et irrégulier : \all{wurde} (passif Präteritum).
    \item \all{Wir \textbf{haben} mit einem Anwalt \textbf{zusammengearbeitet}.} — Particule séparable : \all{zusammen} + \all{ge} + \all{arbeitet}.
\end{enumerate}
\textbf{Point de vigilance :} au \all{Perfekt}, le participe II se place en \emph{fin de phrase} ; l'auxiliaire conjugué occupe la deuxième position.
\end{corrige}

\begin{corrige}[Exercice 5 — Grammaire : Le génitif]
\begin{enumerate}
    \item \all{Das Testament \textbf{des Großvaters} war unklar.} — Nom fort masculin : \all{des} + \all{-s}.
    \item \all{Die Ansprüche \textbf{der Erbin} wurden geprüft.} — Nom féminin : \all{der}, sans terminaison sur le nom.
    \item \all{Der Notar erklärte die Regelung \textbf{der Erbschaft}.} — Féminin : \all{der}.
    \item \all{Wegen \textbf{des Streits} ging die Familie vor Gericht.} — \all{wegen} + génitif ; masculin fort : \all{des Streits} (ou \all{des Streites}).
    \item \all{Das Auto \textbf{meines Vaters} wurde verkauft.} — Possessif au génitif : \all{meines} + \all{-s} sur \all{Vater}.
    \item \all{Die Kinder \textbf{des Verstorbenen} haben ein Anrecht auf den Pflichtteil.} — Adjectif substantivé, déclinaison faible : \all{des Verstorbenen}.
    \item \all{Trotz \textbf{des Widerspruchs} wurde das Haus versteigert.} — \all{trotz} + génitif ; masculin fort : \all{des Widerspruchs}.
    \item \all{Die Unterschrift \textbf{des Zeugen} fehlte auf dem Dokument.} — n-Deklination : \all{der Zeuge} $\rightarrow$ \all{des Zeugen}.
\end{enumerate}
\textbf{Rappel :} au génitif, les noms masculins et neutres forts prennent \all{-(e)s} ; les noms de la n-Deklination (\all{der Zeuge, der Erbe, der Student}) prennent \all{-(e)n} ; les noms féminins restent inchangés.
\end{corrige}

\begin{corrige}[Exercice 6 — Traduction thématique : Français $\rightarrow$ Allemand]
\begin{enumerate}
    \item \all{Der Großvater hat das Haus seiner Enkelin vererbt.} — Datif de la personne : \all{seiner Enkelin}.
    \item \all{Das Testament war nicht gültig, weil die Unterschrift fehlte.} — \all{weil} renvoie le verbe en fin de proposition (Präteritum ici pour le récit).
    \item \all{Die Erben streiten (sich) seit Monaten um das Erbe.} — \all{sich streiten um} + accusatif ; \all{seit} + datif (\all{Monaten}).
    \item \all{Nach dem Gesetz haben die Kinder Anspruch auf den Pflichtteil.} — \all{Anspruch auf} + accusatif.
    \item \all{Der Notar hat die Nachlassregelung gerecht geregelt.}
    \item \all{Wir müssen das Geld unter den drei Brüdern teilen.} — \all{unter} + datif (répartition) ; \all{Brüdern} (datif pluriel -n).
    \item \all{Nach dem Tod des Vaters hat sich die Familie beim Notar versammelt.} — Génitif : \all{des Vaters} ; \all{beim} = \all{bei dem} ; \all{sich versammeln} $\rightarrow$ auxiliaire \all{haben}.
\end{enumerate}
\textbf{Points de vigilance :} majuscule à tous les noms ; auxiliaire \all{haben} avec les verbes transitifs (\all{vererben, regeln, teilen}) ; verbe en deuxième position dans la phrase principale.
\end{corrige}

\begin{corrige}[Exercice 7 — Transformation : Perfekt $\leftrightarrow$ Präteritum]
\textbf{A. Récit au Präteritum (avec discours indirect au Konjunktiv I) :}
\begin{quote}
\all{Der Notar fragte Herrn Müller, was er entschieden habe. Herr Müller antwortete, dass er das Testament geprüft habe: Der Großvater \textbf{habe} das Haus seiner Enkelin vererbt. Auf die Frage, ob die anderen Erben zugestimmt \textbf{hätten}, sagte er, dass sie gestritten \textbf{hätten}. Aber der Notar vermittelte. Schließlich verkauften sie das Haus. Jeder bekam seinen Anteil.}
\end{quote}
Formes clés (Indicatif Präteritum) : \all{entschied, prüfte, vererbte, stimmten zu, stritten, vermittelte, verkauften, bekam}. 
\textbf{Attention :} en discours indirect, l'indicatif Prétérit/Parfait/Plus-que-parfait devient \textbf{Konjunktiv I Perfekt} (\all{habe vererbt, habe gestritten}) ou, pour le plus-que-parfait, souvent \textbf{Konjunktiv II Plusquamperfekt} (\all{hätten zugestimmt, hätten gestritten}) à l'oral/écrit courant.

\textbf{B. Compte rendu oral au Perfekt :}
\begin{quote}
\all{Der Erbstreit \textbf{ist} nach der Beerdigung \textbf{begonnen}. Der älteste Sohn \textbf{hat} das Testament \textbf{geöffnet} und die Klauseln \textbf{vorgelesen}. Die Schwester \textbf{ist} wütend \textbf{geworden} und \textbf{hat} ihren Pflichtteil \textbf{gefordert}. Der Notar \textbf{ist} \textbf{hinzugekommen} und \textbf{hat} das Gesetz \textbf{erklärt}. Die Geschwister \textbf{haben} sich \textbf{geeinigt} und das Grundstück \textbf{verkauft}.}
\end{quote}
\textbf{Points de vigilance :} \all{werden} $\rightarrow$ \all{ist geworden} (auxiliaire \all{sein}) ; \all{hinzukommen} $\rightarrow$ \all{ist hinzugekommen} (mouvement) ; \all{vorlesen} $\rightarrow$ \all{vorgelesen} (particule séparable) ; \all{beginnen} (changement d'état) $\rightarrow$ \all{ist begonnen}.
\end{corrige}

\begin{corrige}[Exercice 8 — Texte à trous]
\begin{quote}
\all{Nach dem Tod \textbf{des Großvaters} fand die Familie ein altes Testament. Darin \textbf{stand}, dass das Haus \textbf{der Enkelin} gehöre. Die beiden Söhne \textbf{waren} empört. Sie \textbf{konnten} nicht glauben, dass ihr Vater sie \textbf{enterbt} habe. Ein Anwalt \textbf{musste} die Lage prüfen. Er \textbf{hat} den Söhnen \textbf{erklärt}, dass das Testament formell richtig \textbf{gewesen sei}. Aber: \textbf{wegen des Gesetzes} \textbf{haben} die Kinder Anspruch auf den Pflichtteil. Schließlich \textbf{haben} sich alle Parteien \textbf{geeinigt}. Das Haus \textbf{ist verkauft worden} und der Erlös \textbf{ist geteilt worden}.}
\end{quote}
\textbf{Justifications :}
\begin{itemize}
    \item \all{des Großvaters} : génitif après \all{der Tod} (complément du nom).
    \item \all{stand} : prétérit fort de \all{stehen}.
    \item \all{der Enkelin} : datif exigé par \all{gehören}.
    \item \all{waren / konnten / musste} : prétérits irréguliers de \all{sein, können, müssen}.
    \item \all{enterbt} : participe II sans \all{ge-} (particule inséparable \all{ent-}).
    \item \all{gewesen sei} : Konjunktiv I Perfekt de \all{sein} (discours indirect).
    \item \all{wegen des Gesetzes} : \all{wegen} + \textbf{génitif} (standard) ; \emph{attention :} \all{wegen dem Gesetz} (datif) est une faute courante à l'oral.
    \item \all{ist verkauft worden} : passif au Perfekt — \all{worden} (et non \all{geworden}) au passif.
\end{itemize}
\end{corrige}

\begin{corrige}[Exercice 9 — Compréhension et expression : Texte type Bac]
\textbf{Fragen de compréhension — réponses modèles :}
\begin{enumerate}
    \item \all{Vor seinem Tod hatte Herr Keller ein handschriftliches Testament verfasst. Darin bestimmte er, dass sein ältester Sohn Thomas das Geschäftshaus in der Innenstadt bekommt.}
    \item \all{Anna und Peter waren nicht einverstanden, weil das Geschäftshaus viel mehr wert war als das Ferienhaus. Außerdem musste Peter Steuern für die Aktien zahlen, deren Wert ständig schwankte.}
    \item \all{Der Notar erklärte das Pflichtteilsrecht: Nach § 2303 BGB steht jedem Kind die Hälfte des gesetzlichen Erbteils zu.}
    \item \all{Der Konflikt wurde gelöst, indem Thomas nach langen Verhandlungen einen Ausgleich an seine Geschwister zahlte.}
    \item \all{Der Streit war zwar beigelegt, aber das Vertrauen unter den Geschwistern blieb beschädigt.}
\end{enumerate}
\textbf{Questions de langue :}
\begin{enumerate}
    \setcounter{enumi}{5}
    \item \textbf{Präteritum (récit) :} \all{schrieb} (\all{schreiben}), \all{bestimmte} (\all{bestimmen}), \all{trafen} (\all{treffen}), \all{sagte}, \all{fügte hinzu}, \all{beruhigte}, \all{erklärte}, \all{zahlte}.
    \item \textbf{Perfekt (relevé dans le texte ou forme correspondante) :} \all{hat verfasst} (\all{verfassen} — Perfekt), \all{ist geworden} (\all{werden}), \all{haben gezahlt} (\all{zahlen}), \all{habe sich geeinigt} (\all{sich einigen} — KI Perfekt si discours indirect). \emph{Note : le texte support est majoritairement au Präteritum.}
    \item \textbf{Génitif :} \all{der Erbteil \textbf{Annas}} ; \all{das Testament \textbf{Herrn Kellers}} ; \all{die Entscheidung \textbf{des Notars}}.
    \item \textbf{Discours indirect (Konjunktiv I) :} \all{Anna sagte, das \textbf{sei} ungerecht.} (\all{sei} = KI de \all{sein} pour \all{ist}).
\end{enumerate}
\textbf{Barème indicatif (20 pts) :} compréhension 10 pts (2 pts par question : 1 pt contenu, 1 pt langue) ; questions de langue 10 pts (6 : 3 pts ; 7 : 2 pts ; 8 : 3 pts ; 9 : 2 pts).
\textbf{Points de vigilance :} répondre par des \emph{phrases complètes} avec verbe conjugué en position 2 ; reprendre le temps de la question ; au discours indirect, \all{ist} devient \all{sei}.
\end{corrige}

\begin{corrige}[Exercice 10 — Thème et Version : Traduction Bac]
\textbf{A. Version — proposition de traduction :}
\begin{quote}
Monsieur Keller était un riche commerçant de Munich. Il avait trois enfants : deux fils et une fille. Avant sa mort, il écrivit un testament manuscrit. Il y stipulait : « Mon fils aîné Thomas recevra l'immeuble commercial du centre-ville. Ma fille Anna recevra la maison de vacances au bord du lac. Mon fils cadet Peter recevra l'argent liquide et les actions. » Après l'enterrement, les frères et sœurs se retrouvèrent chez le notaire. Thomas était satisfait, mais Anna et Peter protestèrent. « C'est injuste ! », dit Anna. « L'immeuble commercial vaut beaucoup plus que la maison de vacances. » Peter ajouta : « Et moi, je dois payer des impôts sur les actions. Leur valeur varie constamment. »
\end{quote}
\textbf{Points de vigilance (version) :} \all{der Kaufmann} = commerçant (et non « homme d'achat ») ; \all{das Geschäftshaus} = immeuble commercial ; \all{hinzufügen} = ajouter ; rendre le prétérit allemand par le passé simple ou l'imparfait selon l'aspect (punctuel / duratif).

\textbf{B. Thème — proposition de traduction :}
\begin{quote}
\all{Das Testament des Großvaters war nicht klar. Die Erben haben sich um das Haus der Familie gestritten. Schließlich hat ein Notar die Nachlassregelung geregelt. Er hat den Kindern das Gesetz erklärt. Das Haus wurde verkauft und das Geld wurde gerecht geteilt.}
\end{quote}
\textbf{Justifications :} \all{des Großvaters} (génitif) ; \all{sich streiten um} + accusatif ; récit au \all{Präteritum} possible aussi : \all{stritten sich, regelte, erklärte} ; passif au prétérit : \all{wurde verkauft / wurde geteilt}.
\textbf{Barème indicatif (20 pts) :} version 10 pts (fidélité 6, qualité du français 4) ; thème 10 pts (vocabulaire 4, grammaire et cas 4, ordre des mots 2).
\end{corrige}

\begin{corrige}[Exercice 11 — Rédaction : Sujet type Bac]
\textbf{Proposition de rédaction modèle (environ 95 mots) :}
\begin{quote}
\all{In meiner Nachbarschaft gab es vor zwei Jahren einen großen Erbstreit. Als der alte Herr Meier starb, fanden seine drei Kinder kein Testament. Der älteste Sohn wollte das Haus behalten, aber seine Schwestern forderten ihren Pflichtteil. Sie stritten monatelang und sprachen nicht mehr miteinander. Schließlich schaltete die Familie einen Notar ein. Er erklärte das Gesetz und regelte die Nachlassfrage: Das Haus wurde verkauft und der Erlös gerecht geteilt. Meiner Meinung nach ist ein klares Testament sehr wichtig. Es schützt die Familie vor Streit und traurigen Konflikten. Geld ist weniger wichtig als Frieden in der Familie.}
\end{quote}
\textbf{Traduction :} Dans mon quartier, il y a eu il y a deux ans un grand conflit d'héritage. Quand le vieux monsieur Meier est mort, ses trois enfants n'ont pas trouvé de testament. Le fils aîné voulait garder la maison, mais ses sœurs ont réclamé leur réserve héréditaire. Ils se sont disputés pendant des mois et ne se parlaient plus. Finalement, la famille a fait appel à un notaire. Il a expliqué la loi et réglé la succession : la maison a été vendue et le produit partagé équitablement. À mon avis, un testament clair est très important. Il protège la famille des disputes et des tristes conflits. L'argent est moins important que la paix dans la famille.

\textbf{Barème indicatif (20 pts) :} respect des consignes et structure 4 pts ; vocabulaire thématique (au moins 5 mots) 4 pts ; correction grammaticale (cas, génitif, temps) 6 pts ; ordre des mots et cohérence 4 pts ; opinion personnelle 2 pts.

\textbf{Points de vigilance :}
\begin{itemize}
    \item Récit au \all{Präteritum} (\all{starb, fanden, wollte, stritten, erklärte, schaltete, regelte}) ; opinion au présent (\all{ist, schützt}).
    \item Vocabulaire exigé présent : \all{der Erbstreit, das Testament, der Pflichtteil, der Notar, regeln, streiten, einschalten}.
    \item Éviter les calques du français : on dit \all{einen Notar einschalten} (faire appel à un notaire), \all{meiner Meinung nach} (à mon avis), \all{sich um etwas streiten} (se disputer pour quelque chose).
\end{itemize}
\end{corrige}

\newpage

\coursec{Fiche bilan}

\begin{popfigure}
\begin{tikzpicture}[mindmap, grow cyclic, every node/.style={concept, align=center, text=white, font=\small}, concept color=popblue, root concept/.append style={font=\small\bfseries, minimum size=2.6cm}, level 1 concept/.append style={sibling angle=60, minimum size=2.2cm, level distance=3.4cm}, level 2 concept/.append style={sibling angle=45, minimum size=1.8cm, level distance=2.2cm, font=\scriptsize}]
\node[root concept, align=center]{Der Streit um das Erbe\\(Le conflit autour de l'héritage)}
 child[concept color=poporange]{ node[align=center]{Lexique\\das Erbe, das Testament, der Erbe, die Erbin, das Gesetz} }
 child[concept color=popgreen]{ node[align=center]{Le génitif\\des Vaters, der Erbin\\possession} }
 child[concept color=poppurple]{ node[align=center]{Perfekt\\haben / sein + Partizip II\\oral, récit} }
 child[concept color=poppink]{ node[align=center]{Präteritum\\verbes faibles : -te\\récit écrit} }
 child[concept color=popgold]{ node[align=center]{Compréhension\\global puis détaillé\\questions} }
 child[concept color=popblue!70]{ node[align=center]{Production\\résumé, avis personnel,\\court paragraphe} };
\end{tikzpicture}
\legende{Carte mentale de la leçon : les problèmes d'héritage}
\end{popfigure}

\begin{bilan}
\textbf{1. Lexique clé (à connaître avec article et pluriel)}
\begin{itemize}
\item \all{das Erbe, -n} : l'héritage ; \all{erben} : hériter ; \all{vererben} : léguer
\item \all{der Erbe, -n} : l'héritier ; \all{die Erbin, -nen} : l'héritière
\item \all{das Testament, -e} : le testament ; \all{das Gesetz, -e} : la loi
\item \all{der Streit, -e} : le conflit ; \all{sich streiten (stritt, hat gestritten)} : se disputer
\item \all{die Familie, -n} : la famille ; \all{der Notar, -e} : le notaire ; \all{der Anwalt, die Anwälte} : l'avocat
\item \all{gerecht} : juste ; \all{die Lösung, -en} : la solution ; \all{teilen} : partager
\end{itemize}

\textbf{2. Grammaire à connaître par cœur}
\begin{center}
\begin{tabularx}{\linewidth}{|l|X|X|}
\hline
\rowcolor{popblueL} Point & Règle & Exemple \\
\hline
Génitif & possession, complément du nom ; \all{des + nom (m/n) + -s/-es} ; \all{der} (f/pl) & \all{das Haus des Vaters} « la maison du père » ; \all{der Wille der Erbin} \\
\hline
Perfekt & auxiliaire \all{haben} (ou \all{sein} : mouvement/changement) + Partizip II en fin & \all{Der Vater hat ein Testament gemacht.} ; \all{Sie ist nach Douala gefahren.} \\
\hline
Präteritum (faibles) & radical + \all{-te, -test, -te, -ten, -tet, -ten} ; récit écrit & \all{machte, arbeitete, streitete} \\
\hline
Präteritum (forts) & radical modifié, à apprendre : \all{stritt, kam, gab, starb} & \all{Die Brüder stritten um das Erbe.} \\
\hline
\end{tabularx}
\end{center}

\textbf{3. Méthodes express}
\begin{itemize}
\item Compréhension : lire le titre et la 1\iere{} phrase $\rightarrow$ idée globale ; souligner les mots du champ lexical (\all{Erbe, Testament, Streit}) ; répondre en phrases complètes avec \all{weil} (verbe en fin).
\item Raconter au passé : à l'oral et dans une lettre $\rightarrow$ Perfekt ; dans un récit écrit $\rightarrow$ Präteritum (\all{sein, haben, werden} toujours au Präteritum : \all{war, hatte, wurde}).
\item Production : 4 à 6 phrases, connecteurs \all{zuerst, dann, danach, schließlich}, un génitif et un verbe au Perfekt minimum.
\end{itemize}
\end{bilan}

\begin{aretenir}[Checklist avant le Bac]
\begin{itemize}
\item[$\square$] Je connais l'article et le pluriel de \all{das Erbe, der Erbe, die Erbin, das Testament, der Streit, das Gesetz}.
\item[$\square$] Je sais former le génitif : \all{des Vaters, der Mutter, der Kinder}.
\item[$\square$] Je sais que le génitif exprime la possession : \all{das Erbe des Großvaters}.
\item[$\square$] Je choisis correctement l'auxiliaire au Perfekt : \all{haben} ou \all{sein}.
\item[$\square$] Je place le Partizip II en fin de phrase au Perfekt.
\item[$\square$] Je conjugue un verbe faible au Präteritum : \all{machte, arbeitete, öffnete}.
\item[$\square$] Je connais le Präteritum de \all{sein (war), haben (hatte), kommen (kam), geben (gab), streiten (stritt)}.
\item[$\square$] Je sais que le verbe conjugué est en 2\ieme{} position et le sujet juste avant ou après.
\item[$\square$] Je mets une majuscule à tous les noms : \all{das Testament, die Familie}.
\item[$\square$] J'évite le faux ami : \all{das Gesetz} = la loi (et non « le geste »).
\item[$\square$] Je sais résumer le texte en 3 phrases avec \all{zuerst, dann, schließlich}.
\item[$\square$] Je sais donner mon avis : \all{Ich finde / Meiner Meinung nach ...}.
\end{itemize}
\end{aretenir}

\findecours

\end{document}
